\section{Formally smooth algebras and Cohen rings}
\label{section:0.19}
\label{0.19}

\setcounter{subsection}{-1}
\subsection{Introduction}
\label{subsection:0.19.0}
\label{0.19.0}

\begin{env}[19.0.1]
\label{0.19.0.1}
\oldpage[0\textsubscript{IV-1}]{69}
In Chapter~IV, we will introduce and study, among other things, an important class of morphisms of preschemes, the \emph{smooth} morphisms\footnote{Also called \emph{simple} morphisms in certain recent works (inspired by the classical terminology of ``simple point''); this terminology lends itself however to confusion, notably in the theory of algebraic groups.}.
One of their fundamental properties (and which, together with a finiteness condition, characterizes them) is a property of \emph{lifting of morphisms}:
if $f:X\to Y$ is a smooth morphism, $g:Y'\to Y$ a morphism, then for every morphism $h:Y''\to Y'$ making $Y''$ a prescheme ``slightly different'' from $Y'$, every $Y$-morphism $Y''\to X$ factorizes as $Y''\xrightarrow{h}Y'\to X$.
More precisely, let us restrict to the case where $Y=\Spec(A)$, $X=\Spec(B)$ are affine;
$B$ is then called a \emph{formally smooth} $A$-algebra if, for every $A$-algebra $C$ and every nilpotent ideal $\mathfrak{J}$ of $C$, every $A$-homomorphism $B\to C/\mathfrak{J}$ \emph{lifts to} $B\to C\to C/\mathfrak{J}$.
In other words, the map
\[
  \Hom_A(B,C) \to \Hom_A(B,C/\mathfrak{J})
\]
is \emph{surjective}.
In many applications, $X$ will appear as a contravariant representable functor $Y'\mapsto F(Y')$ from the category of $Y$-preschemes to that of sets, so that ($\mathbf{0_{III}}$, 8.1.8) one will have $F(Y')=\Hom_Y(Y',X)$.
In the affine case, setting $F^0(C)=F(\Spec(C))$, the verification of the fact that, under the conditions above, $F^0(C)\to F^0(C/\mathfrak{J})$ is \emph{surjective} (which can be done even without knowing that $F$ is representable) will establish that $f$ is smooth, provided the additional finiteness condition is satisfied.
\end{env}

\begin{env}[19.0.2]
\label{0.19.0.2}
For \emph{topological algebras} over a topological ring $A$, there is an analogous notion of \emph{formally smooth algebra} that we will not specify here (cf.~\sref{0.19.3.1}).
The study of these notions is first done from an elementary point of view in \S~19, then, with the aid of the properties of differentials that will be developed in \S\S~20 and 21, more delicate theorems will be proved in \S~22.
Let us summarize here the main results on formally smooth algebras:

\medskip\noindent\textbf{I.} ---
Let $A$, $B$ be two Noetherian local rings, $\varphi:A\to B$ a local homomorphism, $k$ the residue field of $A$, and let $B_0=B\otimes_A k$;
we equip $A$, $B$, and $B_0$ with their preadic topologies and $k$ with the discrete topology.
Then, for $B$ to be a formally smooth $A$-algebra, it is necessary and sufficient that $B$ be a \emph{flat} $A$-module and that $B_0$ be a formally smooth $k$-algebra \sref{0.19.7.1}.
This theorem thus reduces formal smoothness, for Noetherian local rings, to the same question for Noetherian local rings that are algebras over a field.

\medskip\noindent\textbf{II.} ---
Let $k$ be a field, $A$ a Noetherian local ring that is a $k$-algebra.
For $A$ to be formally smooth, it is necessary and sufficient that $A$ be \emph{geometrically regular} over $k$, that is to say that for every finite extension $k'$ of $k$, the semi-local ring $A\otimes_k k'$ be

\oldpage[0\textsubscript{IV-1}]{70}

regular (22.5.8); the completion $\widehat{A}$ of $A$ is then isomorphic to a power series ring $K[[T_1,\ldots,T_n]]$ \sref{0.19.6.5}.
Moreover, the $k$-algebra structure of $A$, when $A$ is assumed complete and formally smooth over $k$, is entirely determined by the residue field $K$ of $A$ and the dimension of $A$;
the latter can furthermore be arbitrary subject to the inequality $\dim(A)\geq\rg_k(\Gamma_{K/k})$, where $\Gamma_{K/k}$ is the ``module of imperfection'' of $K$ (22.2.6).

In particular, for an extension $K$ of $k$ to be formally smooth, it is necessary and sufficient that $K$ be a \emph{separable} extension of $k$ \sref{0.19.6.1}.

\medskip\noindent\textbf{III.} ---
Let $A$ be a Noetherian local ring, $\mathfrak{J}$ an ideal of $A$ distinct from $A$, $A_0=A/\mathfrak{J}$, $B_0$ a complete Noetherian local ring, $A_0\to B_0$ a local homomorphism making $B_0$ a formally smooth $A_0$-algebra.
Then there exists a complete Noetherian local ring $B$, a local homomorphism $A\to B$ making $B$ a flat $A$-module, and an $A_0$-isomorphism $B\otimes_A A_0\xrightarrow{\sim}B_0$ (so that, by~I), $B$ is a formally smooth $A$-algebra);
moreover $B$ is determined by these properties up to isomorphism \sref{0.19.7.2}.
This theorem contains in particular the \emph{theorems of Cohen} on the structure of complete Noetherian local rings \sref{0.19.8}, which will play an important role in \S\S~6 and 7 of Chapter~IV.

\medskip\noindent\textbf{IV.} ---
The interest in the study of formally smooth Noetherian local rings over another comes from the following ``pointwise'' characterization of smooth morphisms:
if $X$ and $Y$ are locally Noetherian preschemes, $f:X\to Y$ a morphism locally of finite type, then, for $f$ to be smooth, it is necessary and sufficient that for every $x\in X$, the ring $\sh{O}_x$ be formally smooth over $\sh{O}_{f(x)}$.
In particular, if $Y=\Spec(k)$, where $k$ is a perfect field, to say that $f$ is smooth is equivalent to saying that $X$ is a \emph{regular} prescheme.

\medskip\noindent\textbf{V.} ---
Finally, we will see in \S\S~20, 21, and 22 that the notion of formally smooth algebra is introduced naturally in the theory of \emph{K\"ahler differentials}, the two theories shedding light on each other.
\end{env}

\begin{env}[19.0.3]
\label{0.19.0.3}
Throughout this section and the following ones, topological rings and modules will be assumed to be \emph{linearly topologized} ($\mathbf{0_I}$, 7.1.1);
the topological rings considered will be assumed to be \emph{commutative}, unless explicit mention to the contrary.
Recall that if $A$ and $B$ are two topological rings, $\rho:A\to B$ a ring homomorphism defining on $B$ a structure of $A$-algebra, we say that $B$ is a \emph{topological $A$-algebra} if $\rho$ is continuous for the topologies considered.

For brevity, in a topological ring $A$ (resp. a topological $A$-module $M$), we will say ``fundamental system of open ideals (resp. submodules)'' instead of ``fundamental system of neighbourhoods of $0$ consisting of ideals (resp. submodules)''.

Given a topological ring $A$ and a topological $A$-module $M$, the sets $\mathfrak{J}M$, where $\mathfrak{J}$ ranges over a fundamental system of open ideals in $A$, form a fundamental system of open submodules of a topology on $M$ making $M$ a topological $A$-module, called the topology \emph{derived} from that of $A$.

Let $M$ be a topological $A$-module whose topology is \emph{less fine} than the topology derived from that of $A$;
then, if $N$ is an open submodule of $M$, the discrete $A$-module $M/N$ is annihilated by an open ideal $\mathfrak{R}$ of $A$, since by hypothesis there exists such an ideal $\mathfrak{R}$ with $\mathfrak{R}\cdot M\subset N$.

\oldpage[0\textsubscript{IV-1}]{71}

If $M$ and $N$ are two topological $A$-modules whose topologies are both derived from that of $A$, then every $A$-homomorphism $u:M\to N$ is \emph{continuous}, since for every neighbourhood $V$ of $0$ in $N$, there exists by hypothesis an open ideal $\mathfrak{J}$ of $A$ such that $\mathfrak{J}N\subset V$, and one thus has $u(\mathfrak{J}M)\subset\mathfrak{J}N\subset V$.

When $B$ is a topological $A$-algebra, the topology on $B$ \emph{derived} from that of $A$ is \emph{finer} than the topology given on $B$, since for every open ideal $\mathfrak{R}$ of $B$, there is by hypothesis an open ideal $\mathfrak{J}$ of $A$ such that $\mathfrak{J}B\subset\mathfrak{R}$.
\end{env}

\subsection{Formal epimorphisms and monomorphisms}
\label{subsection:0.19.1}

\begin{proposition}[19.1.1]
\label{0.19.1.1}
Let $A$ be a topological ring, $M$, $N$ two topological $A$-modules, $(W_\lambda)$ a fundamental system of open submodules in $N$, $u:M\to N$ a continuous $A$-homomorphism, $\widehat{M}$, $\widehat{N}$ the separated completions of $M$ and $N$, $\hat{u}:\widehat{M}\to\widehat{N}$ the continuous extension of $u$ to the separated completions.
\begin{enumerate}
  \item[(i)] The following conditions are equivalent:
  \begin{enumerate}
    \item[a)] $u(M)$ is dense in $N$.
    \item[a$'$)] $\hat{u}(\widehat{M})$ is dense in $\widehat{N}$.
    \item[a$''$)] For all $\lambda$, the composite homomorphism $M\xrightarrow{u}N\to N/W_\lambda$ is surjective.
  \end{enumerate}
  \item[(ii)] The following conditions are equivalent:
  \begin{enumerate}
    \item[b)] The inverse image by $u$ of the topology of $N$ is equal to the topology of $M$.
    \item[b$'$)] $\hat{u}$ is an isomorphism of the topological $\widehat{A}$-module $\widehat{M}$ onto the sub-$\widehat{A}$-module $\hat{u}(\widehat{M})$ of $\widehat{N}$ (which is necessarily closed).
    \item[b$''$)] The $u^{-1}(W_\lambda)$ form a fundamental system of neighbourhoods of $0$ in $M$.
  \end{enumerate}
\end{enumerate}
This follows immediately from the definition of separated completions, and from the fact that $N/W_\lambda$ is discrete.
\end{proposition}

\begin{env}[19.1.2]
\label{0.19.1.2}
When the equivalent conditions of (i) (resp. (ii)) in \sref{0.19.1.1} are satisfied, we say that $u$ is a \emph{formal epimorphism} (resp. a \emph{formal monomorphism}).
We say that $u$ is a \emph{formal bimorphism} if it is both a formal epimorphism and a formal monomorphism;
this amounts to the same, by virtue of \sref{0.19.1.1}, as saying that $\hat{u}$ is an \emph{isomorphism} of the topological $\widehat{A}$-module $\widehat{M}$ onto the topological $\widehat{A}$-module $\widehat{N}$.
\end{env}

\begin{proposition}[19.1.3]
\label{0.19.1.3}
Let $A$ be a topological ring, $M$, $N$ two topological $A$-modules, $u:M\to N$ a continuous $A$-homomorphism, $(V_\lambda)$, $(W_\lambda)$ two fundamental systems of open submodules in $M$ and $N$ respectively, having the same index set, and such that $u(V_\lambda)\subset W_\lambda$ for all $\lambda$;
let $u_\lambda:M/V_\lambda\to N/W_\lambda$ be the homomorphism induced by $u$ by passing to quotients.
Then:
\begin{enumerate}
  \item[(i)] For $u$ to be a formal epimorphism, it is necessary and sufficient that, for all $\lambda$, $u_\lambda$ be surjective.
  \item[(ii)] For $u$ to be a formal monomorphism, it is necessary and sufficient that, for all $\lambda$, there exist $\mu$ such that $V_\mu\subset V_\lambda$ and that $\Ker(u_\mu)$ be contained in the kernel $V_\lambda/V_\mu$ of the canonical map $M/V_\mu\to M/V_\lambda$.
  \item[(iii)] For $u$ to be a formal bimorphism, it is necessary and sufficient that, for all $\lambda$, $u_\lambda$ be surjective and that there exist $\mu$ such that $V_\mu\subset V_\lambda$ and that the canonical homomorphism $M/V_\mu\to M/V_\lambda$

  \oldpage[0\textsubscript{IV-1}]{72}

  factorize as $M/V_\mu\xrightarrow{u_\mu}N/W_\mu\xrightarrow{h}M/V_\lambda$ (where the homomorphism $N/W_\mu\to M/V_\lambda$ is necessarily unique).
\end{enumerate}
\end{proposition}

\begin{proof}
Assertion (i) is immediate;
on the other hand, $\Ker(u_\mu)=u^{-1}(W_\mu)/V_\mu$, and the kernel of the canonical map $M/V_\mu\to M/V_\lambda$ is $V_\lambda/V_\mu$;
the condition of (ii) thus amounts to saying that $u^{-1}(W_\mu)\subset V_\lambda$, and assertion (ii) follows immediately;
finally, when $u_\mu$ is surjective, to say that $\Ker(u_\mu)$ is contained in $V_\lambda/V_\mu$, or to say that $M/V_\mu\to M/V_\lambda$ factorizes as $v\circ u_\mu$, where $v$ is a homomorphism $N/W_\mu\to M/V_\lambda$, is the same thing, since $N/W_\mu$ then identifies with $(M/V_\mu)/\Ker(u_\mu)$.
\end{proof}

\begin{corollary}[19.1.4]
\label{0.19.1.4}
Let $A$ be a topological ring, $M$, $N$ two topological $A$-modules whose topologies are derived from that of $A$, $u:M\to N$ a formal epimorphism.
Let $(\mathfrak{J}_\lambda)$ be a fundamental system of open ideals in $A$.
For $u$ to be a formal bimorphism, it is necessary and sufficient that for all $\lambda$, the homomorphism $u_\lambda:M/\mathfrak{J}_\lambda M\to N/\mathfrak{J}_\lambda N$ induced by $u$ by passing to quotients is bijective.
\end{corollary}

\begin{proof}
Indeed, $u(\mathfrak{J}_\lambda M)\subset\mathfrak{J}_\lambda N$ and one can apply the criterion \sref{0.19.1.3}, (iii);
but if one has a factorization $M/\mathfrak{J}_\mu M\xrightarrow{u_\mu}N/\mathfrak{J}_\mu N\xrightarrow{v}M/\mathfrak{J}_\lambda M$, the image of $\mathfrak{J}_\lambda M/\mathfrak{J}_\mu M$ by $u_\mu$ is $\mathfrak{J}_\lambda N/\mathfrak{J}_\mu N$, hence the image of $\mathfrak{J}_\lambda N/\mathfrak{J}_\mu N$ by $v$ is $0$ and $v$ factorizes as
\[
  N/\mathfrak{J}_\mu N\to N/\mathfrak{J}_\lambda N\xrightarrow{v'}M/\mathfrak{J}_\lambda M;
\]
but then the automorphism identical to that of $M/\mathfrak{J}_\lambda M$ factorizes as
\[
  M/\mathfrak{J}_\lambda M\xrightarrow{u_\lambda}N/\mathfrak{J}_\lambda N\xrightarrow{v'}M/\mathfrak{J}_\lambda M,
\]
which shows that $u_\lambda$ is injective, and as we already know it is surjective, this proves the corollary.
\end{proof}

\begin{proposition}[19.1.5]
\label{0.19.1.5}
Let $A$ be a topological ring, $M$, $N$ two topological $A$-modules, and $u:M\to N$ a continuous $A$-homomorphism.
The following conditions are equivalent:
\begin{enumerate}
  \item[a)] For every discrete topological $A$-module $E$ and every continuous $A$-homomorphism $v:M\to E$, there exists a continuous $A$-homomorphism $w:N\to E$ such that $v=w\circ u$.
  \item[b)] For every open submodule $V'$ of $M$, there exist an open submodule $W''$ of $N$, an open submodule $V''\subset V'\cap u^{-1}(W'')$, and a homomorphism $h:N/W''\to M/V'$ such that the canonical homomorphism $M/V''\to M/V'$ factorizes as
  \[
    M/V''\xrightarrow{u''}N/W''\xrightarrow{h}M/V',
  \]
  where $u''$ is induced by $u$ by passing to quotients.
\end{enumerate}
\end{proposition}

\begin{proof}
To see that a) implies b), apply a) to $E=M/V'$ and to the canonical homomorphism $v:M\to M/V'$.
Then $W''=w^{-1}(0)$ is an open submodule of $N$ and $V''=u^{-1}(W'')$ is an open submodule of $M$ contained in $V'$;
these submodules and the homomorphism $h:N/W''\to M/V'$ induced by $w$ answer the question.
Conversely, to show that b) implies a), one may restrict to the case where $v$ is surjective, replacing $E$ by $v(M)$;
then $V'=\operatorname{Ker}(v)$ is an open submodule of $M$, so $E$ is isomorphic to the discrete $A$-module $M/V'$, and applying b) gives

\oldpage[0\textsubscript{IV-1}]{73}

the desired continuous $A$-homomorphism $w$ by taking the composite $N\to N/W''\xrightarrow{h}M/V'$, the diagram
\[
  \xymatrix{
    M\ar[r]^u\ar[d] & N\ar[d]\\
    M/V''\ar[r]_{u''} & N/W''
  }
\]
being commutative.
\end{proof}

When the equivalent conditions of \sref{0.19.1.5} are satisfied, we say that $u$ is \emph{formally left invertible};
since condition~b) of \sref{0.19.1.5} implies $u^{-1}(W'')\subset V'$, $u$ is then a formal monomorphism.
The terminology is further justified by the following corollaries.

\begin{corollary}[19.1.6]
\label{0.19.1.6}
If there exists a continuous $\widehat A$-homomorphism $s:\widehat N\to\widehat M$ such that $s\circ\widehat u=1_{\widehat M}$, then $u$ is formally left invertible.
\end{corollary}

Note that $\widehat u$ is then a topological isomorphism of $\widehat M$ onto a \emph{topological direct summand} of $\widehat N$.
To prove \sref{0.19.1.6}, it suffices to note that, with the notation of \sref{0.19.1.5}, a), $v:M\to E$ extends to a continuous $\widehat A$-homomorphism $\widehat v:\widehat M\to E$ since $E$ is discrete;
let $j:M\to\widehat M$ and $j':N\to\widehat N$ be the canonical homomorphisms.
Then, setting $w=\widehat v\circ s\circ j'$, one has $w\circ u=\widehat v\circ s\circ\widehat u\circ j=\widehat v\circ j=v$.

\begin{corollary}[19.1.7]
\label{0.19.1.7}
Suppose that the topologies of $M$ and $N$ are derived from that of $A$, and let $(\mathfrak J_\lambda)$ be a fundamental system of open ideals in $A$.
For $u$ to be formally left invertible, it is necessary and sufficient that, for every $\lambda$, the homomorphism
\[
  u_\lambda:M/\mathfrak J_\lambda M\longrightarrow N/\mathfrak J_\lambda N
\]
induced by $u$ by passing to quotients be left invertible (in other words, be an isomorphism of $M/\mathfrak J_\lambda M$ onto a direct summand of $N/\mathfrak J_\lambda N$).
\end{corollary}

Indeed, the condition is sufficient, since for $V'=\mathfrak J_\lambda M$ in condition~b) of \sref{0.19.1.5}, one answers the question by taking $W''=\mathfrak J_\lambda N$, $V''=V'$, and $h$ such that $h\circ u_\lambda$ is the identity on $M/\mathfrak J_\lambda M$.
Conversely, if $u$ is formally left invertible, then for every $\lambda$, by \sref{0.19.1.5}, b), there are $\mu$ with $\mathfrak J_\mu\subset\mathfrak J_\lambda$ and a homomorphism $h:N/\mathfrak J_\mu N\to M/\mathfrak J_\lambda M$ such that the canonical homomorphism $M/\mathfrak J_\mu M\to M/\mathfrak J_\lambda M$ factorizes as
\[
  M/\mathfrak J_\mu M\xrightarrow{u_\mu}N/\mathfrak J_\mu N
  \xrightarrow{h}M/\mathfrak J_\lambda M.
\]
But since
$h(\mathfrak J_\lambda(N/\mathfrak J_\mu N))
 \subset\mathfrak J_\lambda(M/\mathfrak J_\lambda M)=0$,
$h$ factorizes canonically as
\[
  N/\mathfrak J_\mu N\longrightarrow N/\mathfrak J_\lambda N
  \xrightarrow{s}M/\mathfrak J_\lambda M,
\]
and it is immediate that $s$ is a left inverse of $u_\lambda$.

\begin{proposition}[19.1.8]
\label{0.19.1.8}
Let $A$ be a topological ring admitting a countable decreasing fundamental system $(\mathfrak J_n)$ of open ideals.
Let $M$, $N$ be two topological $A$-modules whose topologies are derived from that of $A$, and let $u:M\to N$ be an $A$-homomorphism.
Set $M_n=M/\mathfrak J_nM$, $N_n=N/\mathfrak J_nN$, and let $u_n:M_n\to N_n$ be the $(A/\mathfrak J_n)$-homomorphism induced by $u$ by passing to quotients.
Suppose that, for every $n$, $P_n=\operatorname{Coker}(u_n)$ is a projective $(A/\mathfrak J_n)$-module and that $M$ is separated and complete.
Then $u$ is formally left invertible if and only if $u$ is left invertible (and $u$ is then a topological isomorphism of $M$ onto a topological direct summand of $N$).
\end{proposition}

\oldpage[0\textsubscript{IV-1}]{74}

\begin{proof}
By \sref{0.19.1.7}, one has commutative diagrams
\[
  \xymatrix{
    0\ar[r] & M_n\ar[r]^{u_n} & N_n\ar[r]^{p_n} & P_n\ar[r] & 0\\
    0\ar[r] & M_{n+1}\ar[u]^f\ar[r]_{u_{n+1}} &
      N_{n+1}\ar[u]^g\ar[r]_{p_{n+1}} & P_{n+1}\ar[u]^h\ar[r] & 0
  }
\]
whose rows are split exact sequences;
in other words, for every $n$ there exists a homomorphism $s_n:P_n\to N_n$ such that $p_n\circ s_n=1_{P_n}$.
We show by induction on $n$ that homomorphisms $s'_n:P_n\to N_n$ can be defined so that $p_n\circ s'_n=1_{P_n}$ and the diagrams
\[
  \xymatrix{
    P_n\ar[r]^{s'_n} & N_n\\
    P_{n+1}\ar[u]^h\ar[r]_{s'_{n+1}} & N_{n+1}\ar[u]^g
  }
\]
are commutative.
Indeed, $g\circ s_{n+1}-s'_n\circ h$ is a homomorphism from $P_{n+1}$ into
\[
  u_n(M_n)=u_{n+1}(M_{n+1})/
  \mathfrak J_nu_{n+1}(M_{n+1});
\]
since $P_{n+1}$ is a projective $(A/\mathfrak J_{n+1})$-module, this homomorphism factorizes as
\[
  P_{n+1}\xrightarrow{t_{n+1}}u_{n+1}(M_{n+1})
  \longrightarrow
  u_{n+1}(M_{n+1})/\mathfrak J_nu_{n+1}(M_{n+1}),
\]
whence $s'_{n+1}=s_{n+1}-t_{n+1}$ answers the question.
The decomposition of $N_n$ as the direct sum of $u_n(M_n)$ and $s'_n(P_n)$ then immediately gives a homomorphism $w_n:N_n\to M_n$ which is a left inverse of $u_n$ and makes the diagrams
\[
  \xymatrix{
    N_n\ar[r]^{w_n} & M_n\\
    N_{n+1}\ar[u]^g\ar[r]_{w_{n+1}} & M_{n+1}\ar[u]^f
  }
\]
commutative.
The inverse system $(w_n)$ therefore has an inverse limit $\widehat w:\widehat N\to\widehat M=M$.
Composing with the canonical homomorphism $N\to\widehat N$ gives a homomorphism $w:N\to M$ such that, for every $n$, the endomorphism $(w\circ u)_n$ of $M_n=M/\mathfrak J_nM$ induced by $w\circ u$ is the identity.
Thus $w\circ u=1_M$ because $M$ is separated and complete.
\end{proof}

\begin{proposition}[19.1.9]
\label{0.19.1.9}
Let $A$ be a preadmissible topological ring ($\mathbf{0_I}$, 7.1.2), $\mathfrak Q$ an ideal of definition of $A$, and $(\mathfrak J_\lambda)$ a fundamental system of open ideals in $A$.
Let $M$, $N$ be two topological $A$-modules whose topologies are derived from that of $A$, such that for every $\lambda$, $N/\mathfrak J_\lambda N$ is a projective $(A/\mathfrak J_\lambda)$-module (cf.~\sref{0.19.2.3}).
Let $u:M\to N$ be an $A$-homomorphism.
Then the following conditions are equivalent:
\begin{enumerate}
  \item[a)] $u$ is formally left invertible.

  \oldpage[0\textsubscript{IV-1}]{75}

  \item[b)] The homomorphism $u_0:M/\mathfrak QM\to N/\mathfrak QN$ induced by $u$ by passing to quotients is left invertible.
\end{enumerate}
\end{proposition}

\begin{proof}
By \sref{0.19.1.7}, condition~a) is equivalent to $u_\lambda$ being left invertible for every $\lambda$.
Since $\mathfrak Q$ is open, and therefore contains some $\mathfrak J_\lambda$, it follows at once that $u_0$ is left invertible.
Conversely, to show that b) implies a), note that for every $\lambda$, $\mathfrak Q/\mathfrak J_\lambda$ is by hypothesis a nilpotent ideal of $A/\mathfrak J_\lambda$.
The assertion follows from the next proposition.
\end{proof}

\begin{proposition}[19.1.10]
\label{0.19.1.10}
Let $A$ be a ring, $M$, $N$ two $A$-modules, with $N$ projective, and let $u:M\to N$ be an $A$-homomorphism.
Let $\mathfrak J$ be an ideal of $A$;
suppose that one of the following conditions holds:
\begin{enumerate}
  \item[(i)] $\mathfrak J$ is nilpotent.
  \item[(ii)] $\mathfrak J$ is contained in the radical of $A$ and $M$ is of finite type.
\end{enumerate}
Then $u$ is left invertible if and only if the homomorphism
$u_0:M/\mathfrak JM\to N/\mathfrak JN$ of $(A/\mathfrak J)$-modules induced by $u$ by passing to quotients is left invertible.
\end{proposition}

\begin{proof}
The condition is plainly necessary; let us prove that it is sufficient.
Let $v_0$ be a left inverse of $u_0$.
The composite $N\to N/\mathfrak JN\xrightarrow{v_0}M/\mathfrak JM$ factorizes as $N\xrightarrow{v}M\to M/\mathfrak JM$ because $N$ is projective.
Then $w=v\circ u$ is an endomorphism of $M$ whose induced endomorphism $w_0$ of $M/\mathfrak JM$ is the identity, and it is enough to prove that $w$ is bijective (for then $w^{-1}\circ v$ is a left inverse of $u$).
We distinguish the two cases.

(i) For every $n$ one has the commutative diagram
\[
  \xymatrix{
    (\mathfrak J^n/\mathfrak J^{n+1})\otimes_{A/\mathfrak J}
      (M/\mathfrak JM)\ar[r]\ar[d]_{1\otimes\operatorname{gr}^0(w)} &
      \mathfrak J^nM/\mathfrak J^{n+1}M\ar[d]^{\operatorname{gr}^n(w)}\\
    (\mathfrak J^n/\mathfrak J^{n+1})\otimes_{A/\mathfrak J}
      (M/\mathfrak JM)\ar[r] & \mathfrak J^nM/\mathfrak J^{n+1}M
  }
\]
in which the horizontal arrows are surjective.
Since $\operatorname{gr}^0(w)=w_0$ is the identity, so is $\operatorname{gr}^n(w)$, which is therefore, a fortiori, bijective.
The $\mathfrak J$-preadic filtration on $M$ is finite because $\mathfrak J$ is nilpotent, so $w$ is bijective (Bourbaki, \emph{Alg. comm.}, Chap.~III, \S~2, no.~8, Cor.~3 of Th.~1).

(ii) It suffices to show that, for every maximal ideal $\mathfrak m$ of $A$, the endomorphism $w_{\mathfrak m}$ of $M_{\mathfrak m}$ is bijective (Bourbaki, \emph{Alg. comm.}, Chap.~II, \S~3, no.~3, Th.~1).
Since $\mathfrak JA_{\mathfrak m}\subset\mathfrak mA_{\mathfrak m}$ and
$A_{\mathfrak m}/\mathfrak JA_{\mathfrak m}=(A/\mathfrak J)_{\mathfrak m}$, we are reduced to proving the proposition when $A$ is local.
Moreover, one may suppose that $\mathfrak J$ is the maximal ideal $\mathfrak J_0$ of $A$, because if $u_0$ is left invertible, so is

\oldpage[0\textsubscript{IV-1}]{76}

$u_{00}:M/\mathfrak J_0M\to N/\mathfrak J_0N$, obtained by tensoring $u_0$ with $1_{A/\mathfrak J_0}$, since
\[
  M/\mathfrak J_0M=(M/\mathfrak JM)\otimes_{A/\mathfrak J}(A/\mathfrak J_0).
\]
\footnote{The French source prints $N/\mathfrak J_0M$ as the codomain of $u_{00}$; the module type requires $N/\mathfrak J_0N$.}
Suppose then that $\mathfrak J$ is maximal, so that $A/\mathfrak J$ is a field.
It plainly suffices to prove that $M$ is a free $A$-module under the hypotheses of the statement:
indeed, $\det(w_0)$ is then the canonical image of $\det(w)$ in $A/\mathfrak J$, hence $\det(w)\notin\mathfrak J$ and is therefore invertible.
Now $M/\mathfrak JM$ is a finite-dimensional free $(A/\mathfrak J)$-vector space, so there are a finite free $A$-module $L$ and an $A$-homomorphism $f:L\to M$ such that the induced homomorphism $f_0:L/\mathfrak JL\to M/\mathfrak JM$ is bijective.
Since $M$ is of finite type, $f$ is surjective by Nakayama's lemma (Bourbaki, \emph{Alg. comm.}, Chap.~II, \S~3, no.~2, Cor.~1 of Prop.~4).
Moreover, if $g=u\circ f$, then the induced homomorphism $g_0:L/\mathfrak JL\to N/\mathfrak JN$ is left invertible, and since $L$ is free, the preceding observation shows that $g$ itself is left invertible.
This plainly implies that $f$ is injective, completing the proof.
\end{proof}

Let us mention in passing the following corollaries of \sref{0.19.1.10}.

\begin{corollary}[19.1.11]
\label{0.19.1.11}
Let $A$ be a local ring, $k$ its residue field, $M$ a finite-type $A$-module, $N$ a projective $A$-module, and $u:M\to N$ a homomorphism.
Then $u$ is left invertible if and only if there exists a system of generators $(x_i)_{1\leq i\leq m}$ of $M$ such that the images of the $x_i\otimes1$ under $u\otimes1:M\otimes_Ak\to N\otimes_Ak$ are linearly independent in $N\otimes_Ak$;
the $x_i$ then form a basis of $M$.
\end{corollary}

The condition is plainly necessary, since if $u$ is left invertible, $M$ is a finite-type projective $A$-module and hence free.
Conversely, if the condition holds, the $x_i\otimes1$ plainly form a basis of $M\otimes_Ak$ and $u\otimes1$ is left invertible;
apply \sref{0.19.1.10} to the maximal ideal $\mathfrak J$ of $A$.

\begin{corollary}[19.1.12]
\label{0.19.1.12}
Let $A$ be a ring, $M$ a finite-type $A$-module, $N$ a projective $A$-module, and $u:M\to N$ a homomorphism.
For every prime ideal $\mathfrak p\in\Spec(A)$, the following conditions are equivalent:
\begin{enumerate}
  \item[a)] $u_{\mathfrak p}:M_{\mathfrak p}\to N_{\mathfrak p}$ is left invertible.
  \item[b)] $u\otimes1:M\otimes_Ak(\mathfrak p)\to N\otimes_Ak(\mathfrak p)$ is injective (where $k(\mathfrak p)$ denotes the residue field of $A$ at $\mathfrak p$).
  \item[c)] There is a finite system of elements $x_i\in M$ $(1\leq i\leq m)$ whose images generate $M_{\mathfrak p}$, and a system of $m$ linear forms $y_i^*$ on $N$ $(1\leq i\leq m)$ such that $\det(\langle y_i^*,u(x_j)\rangle)\notin\mathfrak p$.
  \item[d)] There exists $f\in A-\mathfrak p$ such that $u_f:M_f\to N_f$ is left invertible.
\end{enumerate}
Moreover, the set of $\mathfrak p\in\Spec(A)$ satisfying these conditions is open in $\Spec(A)$.
\end{corollary}

The last assertion is an immediate consequence of d).
Since $N$ is projective, it is a direct summand of a free $A$-module $A^{(I)}$;
and since $M$ is of finite type, $u(M)$ is contained in a submodule of $A^{(I)}$ of the form $A^n$ for finite $n$.
Since each of a), b), c), and d) is equivalent to the corresponding assertion with $N$ replaced by $N\oplus P$ (with $u$ still mapping $M$ into $N$), we are reduced to the case where $N$ is free of finite type.
Clearly d) implies a), and a) and b) are equivalent by \sref{0.19.1.11}.
Moreover, a) implies that $M_{\mathfrak p}$ is free by \sref{0.19.1.11}, so a) implies c) by choosing the $x_i$ so that their images form a basis of this

\oldpage[0\textsubscript{IV-1}]{77}

$A_{\mathfrak p}$-module and noting that, because $N$ is free of finite type, every linear form on $N_{\mathfrak p}$ can be written $v/f$, where $f\in A-\mathfrak p$ and $v$ is a linear form on $N$.
Clearly c) implies b), so it remains to see that a) implies d).
Since $N$ is finitely presented, $(\operatorname{Hom}_A(N,M))_{\mathfrak p}$ identifies canonically with $\operatorname{Hom}_{A_{\mathfrak p}}(N_{\mathfrak p},M_{\mathfrak p})$ (Bourbaki, \emph{Alg. comm.}, Chap.~II, \S~2, no.~7, Prop.~19).
If $w_{\mathfrak p}$ is a left inverse of $u_{\mathfrak p}$, there are a homomorphism $w:N\to M$ and an element $f\in A-\mathfrak p$ such that $w_{\mathfrak p}=w\otimes(1/f)1_{A_{\mathfrak p}}$.
The relation $w_{\mathfrak p}\circ u_{\mathfrak p}=1_{M_{\mathfrak p}}$ can therefore also be written $(w\circ u)\otimes1_{A_{\mathfrak p}}=f.1_{M_{\mathfrak p}}$.
Since $M$ is of finite type, there is $g\in A-\mathfrak p$ such that the endomorphism $g((w\circ u)-f.1_M)$ vanishes on all generators of $M$, and hence is zero.
Put $h=fg$, $u_h=u\otimes1_{A_h}$, and $w_h=w\otimes1_{A_h}$.
Then $w_h(u_h(z))=(f/1)z$ for every $z\in M_h$;
since $h/1$ is invertible in $A_h$, so is $f/1$, and $(f/1)^{-1}w_h$ is a left inverse of $u_h$.
This completes the proof.

\begin{proposition}[19.1.13]
\label{0.19.1.13}
Suppose the hypotheses of \sref{0.19.1.9} hold, and suppose in addition that $(\mathfrak J_\lambda)$ is a decreasing sequence $(\mathfrak J_n)$ and that $M$ is separated and complete.
Then conditions~a) and~b) of \sref{0.19.1.9} are also equivalent to:
\begin{enumerate}
  \item[c)] $u$ is left invertible.
\end{enumerate}
\end{proposition}

We already know that c) implies a) by \sref{0.19.1.6}.
Conversely, if a) holds, then with the notation of \sref{0.19.1.8}, $M_n$ is a direct summand of the projective $(A/\mathfrak J_n)$-module $N_n$;
hence $P_n$ is also isomorphic to a direct summand of $N_n$ and is therefore projective.
It remains only to apply \sref{0.19.1.8}.

Finally, note the following proposition.

\begin{proposition}[19.1.14]
\label{0.19.1.14}
Let $A$ be a ring, $M$ a finite-type $A$-module, $N$ a projective $A$-module, and $u:M\to N$ an $A$-homomorphism.
\begin{enumerate}
  \item[(i)] The homomorphism $u$ is left invertible if and only if, for every maximal ideal $\mathfrak m$ of $A$, $u_{\mathfrak m}:M_{\mathfrak m}\to N_{\mathfrak m}$ is left invertible.
  \item[(ii)] Let $A'$ be an $A$-algebra which is a faithfully flat $A$-module.
  Then $u$ is left invertible if and only if $u\otimes1:M\otimes_AA'\to N\otimes_AA'$ is left invertible.
\end{enumerate}
\end{proposition}

As in \sref{0.19.1.12}, it suffices to consider the case where $N$ is free of finite type.
Then saying that $u$ is left invertible means that $u$ is injective and that the quotient module $P=N/u(M)$ is projective, since $u(M)$ is then a direct summand of $N$.
Also, since $M$ is of finite type, $P$ is finitely presented.

(i) Necessity is clear.
Conversely, if the condition holds, $u$ is injective (Bourbaki, \emph{Alg. comm.}, Chap.~II, \S~3, no.~3, Th.~1), and since
$P_{\mathfrak m}=N_{\mathfrak m}/u_{\mathfrak m}(M_{\mathfrak m})$ is projective for every $\mathfrak m$, it follows that $P$ is projective (\emph{loc. cit.}, \S~5, no.~2, Th.~1).

(ii) Again, necessity is immediate.
Conversely, if the condition holds, $u$ is injective ($\mathbf{0_I}$, 6.4.1), and since
$P\otimes_AA'=\operatorname{Coker}(u\otimes1)$ is projective, hence flat, it follows that $P$ is a flat $A$-module ($\mathbf{0_I}$, 6.6.3), and therefore projective because it is finitely presented (Bourbaki, \emph{Alg. comm.}, Chap.~II, \S~5, no.~2, Cor.~2 of Th.~1).

\oldpage[0\textsubscript{IV-1}]{78}

\begin{remark}[19.1.15]
\label{0.19.1.15}
The notion dual to that of a formally left invertible homomorphism is that of a \emph{formally right invertible} homomorphism.
Such a continuous $A$-homomorphism $u:M\to N$ satisfies, by definition, the following condition:
for every open submodule $W'$ of $N$, there exist an open submodule $V'\subset u^{-1}(W')$ of $M$, an open submodule $W''\subset W'$ of $N$, and a homomorphism $h:N/W''\to M/V'$ such that the canonical homomorphism $N/W''\to N/W'$ factorizes as
\[
  N/W''\xrightarrow{h}M/V'\xrightarrow{u'}N/W',
\]
where $u'$ is induced by $u$ by passing to quotients.
This implies that $u$ is a formal epimorphism.
If there exists a continuous $\widehat A$-homomorphism $r:\widehat N\to\widehat M$ such that $\widehat u\circ r=1_{\widehat N}$, then $u$ is immediately seen to be formally right invertible.
Under the hypotheses of \sref{0.19.1.7}, $u$ is formally right invertible if and only if every $u_\lambda$ is right invertible (that is, if $\operatorname{Ker}(u_\lambda)$ is a direct summand of $M/\mathfrak J_\lambda M$ and $u_\lambda$ is an isomorphism from a complement of $\operatorname{Ker}(u_\lambda)$ onto $N/\mathfrak J_\lambda N$).
We leave to the reader the statements and proofs of the propositions corresponding to \sref{0.19.1.8} and \sref{0.19.1.9}:
in the analogue of \sref{0.19.1.8}, one must suppose $M$ separated and complete and $M_n$ a projective $(A/\mathfrak J_n)$-module;
in the analogue of \sref{0.19.1.9}, one must omit the hypothesis on the $N/\mathfrak J_\lambda N$, but instead suppose that, for every $\lambda$, $M/\mathfrak J_\lambda M$ is a projective $(A/\mathfrak J_\lambda)$-module.
\end{remark}

\subsection{Formally projective modules}
\label{subsection:0.19.2}

\begin{definition}[19.2.1]
\label{0.19.2.1}
Let $A$ be a topological ring.
A topological $A$-module $M$ is said to be \emph{formally projective} if it satisfies the following condition:
for every open ideal $\mathfrak J$ of $A$, every pair of discrete $(A/\mathfrak J)$-modules $P$, $Q$, every surjective $A$-homomorphism $u:P\to Q$, and every continuous $A$-homomorphism $v:M\to Q$, there exists a continuous $A$-homomorphism $w:M\to P$ such that $v=u\circ w$.
\end{definition}

\begin{env}[19.2.2]
\label{0.19.2.2}
To verify the condition of \sref{0.19.2.1}, one may plainly restrict, replacing $Q$ by $v(M)$ and $P$ by $u^{-1}(v(M))$, to the case where $v$ itself is surjective.
Then $Q$ is isomorphic to $M/V$, where $V$ is an open submodule of $M$ such that $\mathfrak JM\subset V$.
The condition of \sref{0.19.2.1} is therefore equivalent to requiring that, for every discrete $(A/\mathfrak J)$-module $P$ and every surjective $A$-homomorphism $u:P\to M/V$, there exist an open submodule $V'\subset V$ of $M$ and an $A$-homomorphism $w:M/V'\to P$ such that the canonical homomorphism $M/V'\to M/V$ factorizes as
\[
  M/V'\xrightarrow{w}P\xrightarrow{u}M/V.
\]
It suffices to verify this condition when $\mathfrak J$ ranges over a fundamental system of neighbourhoods of $0$ in $A$ formed of ideals.
\end{env}

\begin{env}[19.2.3]
\label{0.19.2.3}
Suppose that there exist a fundamental system $(\mathfrak J_\lambda)$ of open ideals in $A$ and a fundamental system $(V_\lambda)$ of open submodules of $M$, with the same index set, such that for every $\lambda$, $M/V_\lambda$ is a projective $(A/\mathfrak J_\lambda)$-module.
Then $M$ is formally projective.
Indeed, with the notation of \sref{0.19.2.2}, choose $\lambda$ such that $\mathfrak J_\lambda\subset\mathfrak J$ and $V_\lambda\subset V$.
Since $P$ is also an $(A/\mathfrak J_\lambda)$-module, the factorization of the canonical homomorphism $M/V_\lambda\to M/V$ as $M/V_\lambda\to P\to M/V$ follows from its being an $(A/\mathfrak J_\lambda)$-homomorphism and from the projectivity of $M/V_\lambda$.

When the stronger condition in this paragraph holds, $M$ is said to be \emph{strictly formally projective}.
\end{env}

\begin{proposition}[19.2.4]
\label{0.19.2.4}
Let $A$ be a topological ring, and let $M$ be a topological $A$-module whose topology is derived from that of $A$.
Then $M$ is formally projective if and only if it is strictly formally projective.
\end{proposition}

\oldpage[0\textsubscript{IV-1}]{79}

\begin{proof}
We have just seen that the condition is sufficient.
Conversely, suppose that $M$ is formally projective, and let $(\mathfrak J_\lambda)$ be a fundamental system of open ideals in $A$.
For each $\lambda$, let $P$ be a free $(A/\mathfrak J_\lambda)$-module and let $u:P\to M/\mathfrak J_\lambda M$ be a surjective $(A/\mathfrak J_\lambda)$-homomorphism.
There is then a $\mathfrak J_\mu\subset\mathfrak J_\lambda$ such that the canonical homomorphism $M/\mathfrak J_\mu M\to M/\mathfrak J_\lambda M$ factorizes as
\[
  M/\mathfrak J_\mu M\xrightarrow{w}P\xrightarrow{u}M/\mathfrak J_\lambda M.
\]
But since $w(\mathfrak J_\lambda(M/\mathfrak J_\mu M))\subset\mathfrak J_\lambda P=0$, $w$ factorizes through a map $v:M/\mathfrak J_\lambda M\to P$, and it is clear that $u\circ v$ is the identity on $M/\mathfrak J_\lambda M$.
Thus this $(A/\mathfrak J_\lambda)$-module is projective.
\end{proof}

\begin{proposition}[19.2.5]
\label{0.19.2.5}
Let $A$ be a topological ring, $M$ a topological $A$-module.
\begin{enumerate}
  \item[(i)] The module $M$ is formally projective (resp. strictly formally projective) if and only if the topological $\widehat A$-module $\widehat M$ is so.
  \item[(ii)] Let $A'$ be a topological $A$-algebra.
  If $M$ is formally projective (resp. strictly formally projective), then $M\otimes_A A'$, equipped with the tensor-product topology, is a formally projective (resp. strictly formally projective) topological $A'$-module.
\end{enumerate}
\end{proposition}

\begin{proof}
\begin{enumerate}
  \item[(i)] When $\mathfrak J$ (resp. $V$) ranges over the open ideals of $A$ (resp. the open submodules of $M$), the separated completion $\widehat{\mathfrak J}$ (resp. $\widehat V$) ranges over the open ideals of $\widehat A$ (resp. the open submodules of $\widehat M$), and $\widehat A/\widehat{\mathfrak J}=A/\mathfrak J$ (resp. $\widehat M/\widehat V=M/V$), up to canonical isomorphism.
  Since formal projectivity (resp. strict formal projectivity) involves only the $A/\mathfrak J$ and the $M/V$, assertion~(i) follows immediately.

  \item[(ii)] Suppose first that $M$ is formally projective and put $M'=M\otimes_A A'$.
  Let $\mathfrak J'$ be an open ideal of $A'$, let $P'$ and $Q'$ be two discrete $(A'/\mathfrak J')$-modules, let $u':P'\to Q'$ be a surjective $A'$-homomorphism, and let $v':M'\to Q'$ be a continuous $A'$-homomorphism.
  There is an open ideal $\mathfrak J$ of $A$ such that $\mathfrak JA'\subset\mathfrak J'$, so $P'$ and $Q'$ are also discrete $(A/\mathfrak J)$-modules.
  Applying formal projectivity to the continuous composite $v:M\xrightarrow{j}M'\xrightarrow{v'}Q'$, one obtains a continuous $A$-homomorphism $w:M\to P'$ such that $v=u'\circ w$.
  Since $P'$ is a topological $A'$-module, $w$ factorizes as $M\xrightarrow{j}M'\xrightarrow{w'}P'$, where $w'$ is a continuous $A'$-homomorphism; from $v'\circ j=(u'\circ w')\circ j$ one concludes that $v'=u'\circ w'$.

  Now suppose that $M$ is strictly formally projective.
  Let $(\mathfrak J_\lambda)$ (resp. $(V_\lambda)$) be a fundamental system of open ideals (resp. open submodules) in $A$ (resp. $M$), with $M/V_\lambda$ projective over $A/\mathfrak J_\lambda$, and let $(\mathfrak J'_\mu)$ be a fundamental system of open ideals in $A'$.
  A fundamental system of neighbourhoods of $0$ in $M'$ is formed by
  \[
    W_{\lambda\mu}=\operatorname{Im}(M\otimes_A\mathfrak J'_\mu)+\operatorname{Im}(V_\lambda\otimes_A A'),
  \]
  where $\lambda$ and $\mu$ satisfy $\mathfrak J_\lambda A'\subset\mathfrak J'_\mu$.
  Since
  \[
    (M\otimes_A A')/W_{\lambda\mu}=(M/V_\lambda)\otimes_{A/\mathfrak J_\lambda}(A'/\mathfrak J'_\mu)
  \]
  and $M/V_\lambda$ is projective over $A/\mathfrak J_\lambda$, the module $M'/W_{\lambda\mu}$ is projective over $A'/\mathfrak J'_\mu$.
\end{enumerate}
\end{proof}

\subsection{Formally smooth algebras}
\label{subsection:0.19.3}

\begin{definition}[19.3.1]
\label{0.19.3.1}
Let $A$ be a topological ring, $B$ a topological $A$-algebra.
We say that $B$ is a \emph{formally smooth} $A$-algebra if, for every discrete topological $A$-algebra $C$, and every nilpotent ideal $\mathfrak{J}$ of $C$, every continuous $A$-homomorphism $u:B\to C/\mathfrak{J}$ \emph{factorizes} as $B\xrightarrow{v}C\xrightarrow{\varphi}C/\mathfrak{J}$, where $v$ is a continuous $A$-homomorphism and $\varphi$ the canonical homomorphism.
\end{definition}

\oldpage[0\textsubscript{IV-1}]{80}
The definition \sref{0.19.3.1} amounts to saying that the following property holds:
\begin{enumerate}
  \item[(P)] For every open ideal $\mathfrak{R}$ of the $A$-algebra $B$ and every $A$-homomorphism $u':B/\mathfrak{R}\to C/\mathfrak{J}$, there exists an open ideal $\mathfrak{R}'\subset\mathfrak{R}$ of $B$ such that the homomorphism
  $B\to B/\mathfrak{R}\xrightarrow{u'}C/\mathfrak{J}$ factorizes as
  \[
    B\to B/\mathfrak{R}'\xrightarrow{v'}C\to C/\mathfrak{J},
  \]
  where $v'$ is an $A$-homomorphism.
\end{enumerate}

Indeed, if $u:B\to C/\mathfrak{J}$ is a continuous $A$-homomorphism, its kernel $\mathfrak{R}$ is an open ideal of $B$, and if (P) is satisfied, it suffices to take $u'$ and to take for $v$ the composition $B\to B/\mathfrak{R}\xrightarrow{v'}C$ to satisfy the conditions of \sref{0.19.3.1}.
Conversely, suppose $B$ is a formally smooth $A$-algebra;
let there be given an open ideal $\mathfrak{R}$ of $B$ and a continuous $A$-homomorphism $u':B/\mathfrak{R}\to C/\mathfrak{J}$;
if $v:B\to C$ is a continuous $A$-homomorphism such that $u$ factorizes as $B\xrightarrow{v}C\to C/\mathfrak{J}$, the ideal $\mathfrak{R}'=\Ker(v)\cap\mathfrak{R}$ of $B$ is open and contained in $\mathfrak{R}$;
by taking $v'$ to factorize as $B\to B/\mathfrak{R}'\xrightarrow{v'}C$, and $v'$ clearly satisfies the condition (P).

\begin{proposition}[19.3.2]
\label{0.19.3.2}
Let $A$ be a discrete ring, $V$ a projective $A$-module;
the symmetric algebra $B=\mathrm{S}_A^\bullet(V)$, equipped with the discrete topology, is a formally smooth $A$-algebra.
\end{proposition}

\begin{proof}
Indeed, the notations $C$ and $\mathfrak{J}$ having the same meaning as in \sref{0.19.3.1}, let $u:B\to C/\mathfrak{J}$;
as $V$ is projective, $u_1$, the restriction of $u$ to $V$, lifts to a homomorphism of $A$-modules $v_1:V\to C$ such that $u_1$ factorizes as $V\xrightarrow{v_1}C\xrightarrow{\varphi}C/\mathfrak{J}$;
and $v_1$ extends to a homomorphism of $A$-algebras $v:\mathrm{S}_A^\bullet(V)\to C$, and $v$ coincides with $u$ in $V$, hence $\varphi\circ v=u$, which proves the proposition.
\end{proof}

\begin{corollary}[19.3.3]
\label{0.19.3.3}
If $A$ is a discrete ring, every polynomial algebra $B=A[T_\alpha]_{\alpha\in I}$, equipped with the discrete topology, is a formally smooth $A$-algebra.
\end{corollary}

\begin{proposition}[19.3.4]
\label{0.19.3.4}
Let $A$ be a topological ring, and let $B=A[[T_\alpha]]_{\alpha\in I}$ be a ring of formal power series $\sum c_\lambda T^\lambda$ (algebra large over $A$ of the monoid $\bb{N}^{(I)}$, identified as a topological $A$-module with the product $A^{\bb{N}^{(I)}}$).
If one equips $B$ with the product topology, $B$ is a formally smooth topological $A$-algebra.
\end{proposition}

Let $(\mathfrak{L}_\mu)_{\mu\in M}$ be a fundamental system of open ideals of $A$.
For every finite subset $H$ of $I$, every $\mu\in M$, and every integer $n$, let $\mathfrak{K}_{H,\mu,n}$ be the neighbourhood of $0$ in $B$ consisting of the families $(c_\lambda)$ such that $c_\lambda\in\mathfrak{L}_\mu$ whenever $\lambda=(\lambda_\alpha)_{\alpha\in I}$ satisfies $\lambda_\alpha=0$ for $\alpha\notin H$ and $\sum_{\alpha\in H}\lambda_\alpha\leq n$.
The $\mathfrak{K}_{H,\mu,n}$ plainly form a fundamental system of neighbourhoods of $0$ in $B$ and are ideals of $B$; hence the product topology is compatible with the $A$-algebra structure of $B$.

We first record the following lemma.

\begin{lemma}[19.3.4.1]
\label{0.19.3.4.1}
Let $E$ be a discrete $A$-algebra.
\begin{enumerate}
  \item[\rm(i)] If $f:B\to E$ is a continuous $A$-homomorphism, there exists a finite subset $H$ of $I$ such that $f(T_\alpha)=0$ for $\alpha\notin H$, and $f(T_\alpha)$ is nilpotent in $E$ for every $\alpha\in H$.
  \item[\rm(ii)] Conversely, let $H$ be a finite subset of $I$ and $(z_\alpha)_{\alpha\in H}$ a family of nilpotent elements of $E$.
  There exists a unique continuous $A$-homomorphism $g:B\to E$ such that $g(T_\alpha)=z_\alpha$ for $\alpha\in H$ and $g(T_\alpha)=0$ for $\alpha\notin H$.
\end{enumerate}
\end{lemma}

\oldpage[0\textsubscript{IV-1}]{81}

\begin{proof}
Assertion (i) follows immediately from the fact that $f^{-1}(0)$ is a neighbourhood of $0$ in the product $A$-module $A^{\bb{N}^{(I)}}$: thus $f(T^\lambda)=0$ except for finitely many $\lambda\in\bb{N}^{(I)}$.

For (ii), the polynomial ring $B'=A[T_\alpha]_{\alpha\in I}$ is dense in $B$.
Existence and uniqueness on $B'$ are immediate.
Continuity follows from nilpotence: if $z_\alpha^n=0$ for all $\alpha\in H$, then $g(T^\lambda)=0$ unless $\lambda_\alpha=0$ for $\alpha\notin H$ and $\lambda_\alpha\leq n$ for $\alpha\in H$, and there are only finitely many such $\lambda$.
\end{proof}

\begin{proof}[Proof of Proposition~{\rm\sref{0.19.3.4}}]
Let $C$ and $\mathfrak{J}$ have the meaning of \sref{0.19.3.1}, and let $u:B\to C/\mathfrak{J}$ be a continuous $A$-homomorphism.
By \sref{0.19.3.4.1}, $u(T_\alpha)=0$ except for the indices in a finite subset $H$ of $I$, and the elements $z_\alpha=u(T_\alpha)$, $\alpha\in H$, are nilpotent in $C/\mathfrak{J}$.
Since $\mathfrak{J}$ is nilpotent, choose nilpotent elements $x_\alpha\in C$ lifting the $z_\alpha$.
The lemma gives a continuous $A$-homomorphism $v:B\to C$ such that $v(T_\alpha)=x_\alpha$ for $\alpha\in H$ and $v(T_\alpha)=0$ otherwise.
Then $u$ is the composite $B\xrightarrow{v}C\to C/\mathfrak{J}$.
\end{proof}

\begin{proposition}[19.3.5]
\label{0.19.3.5}
Let $A$ be a topological ring.
\begin{enumerate}
  \item[(i)] $A$ is a formally smooth $A$-algebra.
  \item[(ii)] If $B$ is a formally smooth $A$-algebra and $C$ a formally smooth $B$-algebra, then $C$ is a formally smooth $A$-algebra.
  \item[(iii)] Let $B$ be a formally smooth $A$-algebra, $A'$ a topological $A$-algebra;
  then the topological $A'$-algebra $B\otimes_A A'$ ($\mathbf{0_I}$, 7.7.5 and 7.7.6) is formally smooth.
  \item[(iv)] Let $B$ be a topological $A$-algebra, $S$ (resp. $T$) a multiplicative subset of $A$ (resp. $B$) such that the canonical image of $S$ in $B$ is contained in $T$.
  If $B$ is a formally smooth $A$-algebra, then $T^{-1}B$ is an $S^{-1}A$-algebra formally smooth.
  \item[(v)] Let $B_i$ ($1\leq i\leq n$) be topological $A$-algebras.
  For $\prod_{i=1}^n B_i$ to be a formally smooth $A$-algebra, it is necessary and sufficient that each of the $B_i$ be so.
\end{enumerate}
\end{proposition}

\oldpage[0\textsubscript{IV-1}]{82}

\begin{proof}
\begin{enumerate}
  \item[(i)] If $C$ is a discrete $A$-algebra and
  $\varphi:C\to C/\mathfrak{J}$ is the canonical homomorphism onto any
  quotient $A$-algebra of $C$, the \emph{only} $A$-homomorphism from $A$ to
  $C/\mathfrak{J}$ is the composite
  $A\xrightarrow{\psi}C\xrightarrow{\varphi}C/\mathfrak{J}$, where $\psi$
  defines the $A$-algebra structure of $C$; since $\psi$ is continuous, the
  condition of \sref{0.19.3.1} is trivially satisfied.

  \item[(ii)] Let $\alpha:A\to B$ and $\beta:B\to C$ be the continuous
  homomorphisms defining respectively the $A$-algebra structure on $B$ and the
  $B$-algebra structure on $C$, so that $\beta\circ\alpha$ defines the
  $A$-algebra structure on $C$. Let $E$ be a discrete $A$-algebra,
  $\mathfrak{L}$ a nilpotent ideal of $E$, and
  $u:C\to E/\mathfrak{L}$ a continuous $A$-homomorphism. Since $B$ is a
  formally smooth $A$-algebra, $u\circ\beta:B\to E/\mathfrak{L}$ factorizes as
  $B\xrightarrow{v}E\to E/\mathfrak{L}$, where $v$ is a continuous
  $A$-homomorphism. The homomorphisms $v$ and $u\circ\beta$ then define on $E$
  and $E/\mathfrak{L}$ respectively topological $B$-algebra structures for
  which $E/\mathfrak{L}$ is still the discrete quotient algebra of $E$.
  Moreover, $u$ is a continuous $B$-homomorphism, and hence factorizes as
  $C\xrightarrow{w}E\to E/\mathfrak{L}$, where $w$ is a continuous
  $B$-homomorphism; since $w\circ\beta\circ\alpha$ is the homomorphism defining
  the $A$-algebra structure on $E$, $w$ is indeed a continuous
  $A$-homomorphism.

  \item[(iii)] Let $C$ be a discrete topological $A'$-algebra,
  $\mathfrak{J}$ a nilpotent ideal of $C$, and
  $u:B\otimes_A A'\to C/\mathfrak{J}$ a continuous $A'$-homomorphism. Composing
  $u$ with the canonical homomorphism
  $\rho:B\to B\otimes_A A'$ gives a continuous $A$-homomorphism, which by
  hypothesis factorizes as $B\xrightarrow{v}C\to C/\mathfrak{J}$. If
  $\sigma:A'\to B\otimes_A A'$ and $w:A'\to C$ are the canonical
  homomorphisms, equality of the composites $A\to B\xrightarrow{v}C$ and
  $A\to A'\xrightarrow{w}C$ gives a continuous ring homomorphism
  $f:B\otimes_A A'\to C$ such that $v=f\circ\rho$ and $w=f\circ\sigma$
  ($\mathbf{0_I}$, 7.7.6). Comparing the composites from $B$ and $A'$ shows
  that $u$ factorizes as
  $B\otimes_A A'\xrightarrow{f}C\to C/\mathfrak{J}$, proving~(iii).

  \item[(iv)] The topology on $S^{-1}A$ (resp. $T^{-1}B$) has as a fundamental
  system of neighbourhoods of $0$ the ideals $S^{-1}\mathfrak{J}$ (resp.
  $T^{-1}\mathfrak{R}$), where $\mathfrak{J}$ (resp. $\mathfrak{R}$) ranges
  over a fundamental system of open ideals of $A$ (resp. $B$)
  ($\mathbf{0_I}$, 7.6.1). Let $C$ be a discrete topological
  $S^{-1}A$-algebra, $\mathfrak{J}$ a nilpotent ideal of $C$, and
  $u:T^{-1}B\to C/\mathfrak{J}$ a continuous $S^{-1}A$-homomorphism. The
  composite $B\to T^{-1}B\xrightarrow{u}C/\mathfrak{J}$ factorizes as
  $B\xrightarrow{v}C\to C/\mathfrak{J}$. For every $t\in T$, $u(t/1)$ is
  invertible in $C/\mathfrak{J}$; since $\mathfrak{J}$ is nilpotent, $v(t)$ is
  invertible in $C$ ($\mathbf{0_I}$, 7.1.12). Thus $v$ factorizes through a
  continuous homomorphism $w:T^{-1}B\to C$ ($\mathbf{0_I}$, 7.6.6). The
  canonical diagrams show that $w$ is an $S^{-1}A$-homomorphism and that
  $u$ is the composite $T^{-1}B\xrightarrow{w}C\to C/\mathfrak{J}$.

  \item[(v)] This is immediate: giving a continuous $A$-homomorphism from
  $B=\prod_{i=1}^n B_i$ to $C$ (resp. $C/\mathfrak{J}$) is equivalent to
  giving $n$ continuous $A$-homomorphisms $B_i\to C$ (resp.
  $B_i\to C/\mathfrak{J}$), and every such homomorphism from a factor gives by
  composition a continuous homomorphism from $B$.
\end{enumerate}
\end{proof}

\oldpage[0\textsubscript{IV-1}]{83}

\begin{proposition}[19.3.6]
\label{0.19.3.6}
Let $A$ be a topological ring, $B$ a topological $A$-algebra, $\widehat{A}$ and $\widehat{B}$ the separated completions of $A$ and $B$, so that $\widehat{B}$ is a topological $\widehat{A}$-algebra.
The following conditions are equivalent:
\begin{enumerate}
  \item[a)] $B$ is a formally smooth $A$-algebra.
  \item[b)] $\widehat{B}$ is a formally smooth $A$-algebra.
  \item[c)] $\widehat{B}$ is a formally smooth $\widehat{A}$-algebra.
\end{enumerate}
\end{proposition}

\begin{proof}
Of course, the structure of $\widehat{A}$-algebra on $\widehat{B}$ is defined by the homomorphism $\hat{\varphi}$, where $\varphi:A\to B$ is the homomorphism defining the structure of $A$-algebra on $B$.
As every discrete $A$-algebra $C$ is separated and complete, it is also a $\widehat{A}$-algebra (by canonical prolongation of the homomorphism from $A$ to $C$), and a continuous $A$-homomorphism from $B$ to $C$ gives by extension a canonical $\widehat{A}$-homomorphism (and \emph{a fortiori} an $A$-homomorphism) from $\widehat{B}$ to $C$ (in other words, every continuous $A$-homomorphism $B\to C$ factorizes as $B\to\widehat{B}\to C$ uniquely).
These remarks and the definition \sref{0.19.3.1} immediately imply the proposition.
\end{proof}

\begin{corollary}[19.3.7]
\label{0.19.3.7}
Under the hypotheses of \sref{0.19.3.5}, (iv), the topological $A\{S^{-1}\}$-algebra $B\{T^{-1}\}$ is formally smooth.
\end{corollary}

This follows from the definitions ($\mathbf{0_I}$, 7.6.1 and 7.6.7),
\sref{0.19.3.5},~(iv), and \sref{0.19.3.6}.

\begin{proposition}[19.3.8]
\label{0.19.3.8}
Let $A$ be a topological ring, $B$ a topological $A$-algebra, and suppose that for every open ideal $\mathfrak{R}$ of $B$, $\mathfrak{R}^2$ is also open.
Let $A'$, $B'$ be the rings $A$ and $B$ equipped with topologies finer than the initial topologies, and suppose that the canonical homomorphism $\varphi:A\to B$ is still a continuous homomorphism from $A'$ to $B'$.
Then, if $B'$ is a formally smooth $A'$-algebra, $B$ is a formally smooth $A$-algebra.
\end{proposition}

\begin{proof}
It suffices to apply the following lemma:
\end{proof}

\begin{lemma}[19.3.8.1]
\label{0.19.3.8.1}
Let $C$ be a discrete $A$-algebra, $\mathfrak{J}$ a nilpotent ideal of $C$.
Suppose that the square of every open ideal of $B$ is also open.
Then, if a composite homomorphism $v:B\xrightarrow{u}C\to C/\mathfrak{J}$ is continuous, the homomorphism $u$ is continuous.
\end{lemma}

\begin{proof}
Indeed, $\mathfrak{R}=u^{-1}(\mathfrak{J})$ is the kernel of $v$, and it is open by hypothesis;
as $\mathfrak{J}^n=0$, $\mathfrak{R}^n$ is contained in $\Ker(u)$;
but by induction on $h$, it follows from the hypothesis that $\mathfrak{R}^{2^h}$ is also open, hence $\mathfrak{R}^n$ for all $n$, and consequently $\Ker(u)$ is also open, which proves our assertion.
\end{proof}

Note that the hypothesis on $B$ means that the topology of $B$ is the
supremum of the $\mathfrak{R}$-preadic topologies, where $\mathfrak{R}$ ranges
over the open ideals of $B$. If $B$ is a preadic ring ($\mathbf{0_I}$, 7.1.9),
this condition is always satisfied.

\begin{env}[19.3.9]
\label{0.19.3.9}
We will now see that the property of being formally smooth implies properties of ``lifting'' of homomorphisms under more general conditions than those of the definition \sref{0.19.3.1}.
\end{env}

\oldpage[0\textsubscript{IV-1}]{84}

\begin{proposition}[19.3.10]
\label{0.19.3.10}
Let $A$ be a topological ring, $B$ a formally smooth topological $A$-algebra.
Let $C$ be a topological $A$-algebra, $\mathfrak{J}$ an ideal of $C$, satisfying the following conditions:
\begin{enumerate}
  \item[$1^\circ$] $C$ is metrizable and complete.
  \item[$2^\circ$] $\mathfrak{J}$ is closed and the sequence $(\mathfrak{J}^n)$ tends to $0$.
\end{enumerate}
Then, every continuous $A$-homomorphism $u:B\to C/\mathfrak{J}$ factorizes as $B\xrightarrow{v}C\to C/\mathfrak{J}$, where $v$ is a continuous $A$-homomorphism.
\end{proposition}

\begin{proof}
Let $(\mathfrak{Q}_n)$ be a decreasing sequence of ideals of $C$ forming a fundamental system of neighbourhoods of $0$ in this algebra;
as there exists $k$ such that $\mathfrak{J}^k\subset\mathfrak{Q}_n$, $(\mathfrak{J}+\mathfrak{Q}_n)/\mathfrak{Q}_n$ is nilpotent.
For all $n$, let $u_n$ be the continuous $A$-homomorphism $B\xrightarrow{u}C/\mathfrak{J}\to C/(\mathfrak{J}+\mathfrak{Q}_n)=(C/\mathfrak{Q}_n)/((\mathfrak{J}+\mathfrak{Q}_n)/\mathfrak{Q}_n)$;
by hypothesis $u_n$ factorizes as $B\xrightarrow{v_n}C/\mathfrak{Q}_n\xrightarrow{\varphi_n}C/(\mathfrak{J}+\mathfrak{Q}_n)$, where $v_n$ is a continuous $A$-homomorphism.
Let us show that one can choose step by step the $v_n$ so that $v_n$ factorizes as
\[
  B\xrightarrow{v_{n+1}}C/\mathfrak{Q}_{n+1}\to C/\mathfrak{Q}_n
\]
for all $n$ (in other words, the $v_n$ form a \emph{projective system} of homomorphisms).
This will follow from the following lemma:
\end{proof}

\begin{lemma}[19.3.10.1]
\label{0.19.3.10.1}
Let $B$ be a formally smooth $A$-algebra, $E$, $E'$ two discrete topological $A$-algebras, $\mathfrak{R}$ (resp. $\mathfrak{R}'$) a nilpotent ideal of $E$ (resp. $E'$), $f:E\to E'$ a surjective $A$-homomorphism such that $f(\mathfrak{R})\subset\mathfrak{R}'$, $f':E/\mathfrak{R}\to E'/\mathfrak{R}'$ the homomorphism induced by $f$ by passing to quotients.
Let $u:B\to E/\mathfrak{R}$ be a continuous $A$-homomorphism, $u'$ the composite homomorphism $B\xrightarrow{u}E/\mathfrak{R}\xrightarrow{f'}E'/\mathfrak{R}'$, and let $v':B\to E'$ be a continuous $A$-homomorphism such that $u'$ factorizes as $B\xrightarrow{v'}E'\to E'/\mathfrak{R}'$.
Then there exists a continuous $A$-homomorphism $v:B\to E$ such that $v'$ factorizes as $B\xrightarrow{v}E\xrightarrow{f}E'$.
\end{lemma}

\begin{proof}
In the commutative diagram
\[
\xymatrix{
  E \ar[r]^f \ar[d] & E' \ar[d] \\
  E/\mathfrak{R} \ar[r]_{f'} & E'/\mathfrak{R}'
}
\]
all the homomorphisms are surjective.
If $\mathfrak{L}=\Ker(f)$, then $E'$ is identified with $E/\mathfrak{L}$ and $\mathfrak{R}'$ with $(\mathfrak{R}+\mathfrak{L})/\mathfrak{L}$.
We shall use the following elementary lemma.
\end{proof}

\begin{lemma}[19.3.10.2]
\label{0.19.3.10.2}
Let $F$ be a ring, not necessarily commutative, and let $\mathfrak{a}$ and $\mathfrak{b}$ be two-sided ideals of $F$.
Then the fibre product
\[
  (F/\mathfrak{a})\times_{F/(\mathfrak{a}+\mathfrak{b})}(F/\mathfrak{b})
\]
is canonically identified with $F/(\mathfrak{a}\cap\mathfrak{b})$.
\end{lemma}

\begin{proof}
This is the special case of \sref{0.18.1.7} obtained by replacing, in diagram \sref{0.18.1.7.1}, $G$ by $F/(\mathfrak{a}\cap\mathfrak{b})$, $\mathfrak{J}'$ by $\mathfrak{b}/(\mathfrak{a}\cap\mathfrak{b})$, $F$ by $F/\mathfrak{b}$, $\mathfrak{K}$ by $(\mathfrak{a}+\mathfrak{b})/\mathfrak{b}$, $B$ by $F/(\mathfrak{a}+\mathfrak{b})$, and $\theta$ by the canonical isomorphism
\[
  (\mathfrak{a}+\mathfrak{b})/\mathfrak{b}\xrightarrow{\sim}\mathfrak{a}/(\mathfrak{a}\cap\mathfrak{b}).
\]
\end{proof}

\oldpage[0\textsubscript{IV-1}]{85}

\begin{proof}[Continuation of the proof of Lemma~{\rm\sref{0.19.3.10.1}}]
By \sref{0.19.3.10.2}, the commutativity of
\[
\xymatrix{
  B \ar[r]^{v'} \ar[d]_u & E'=E/\mathfrak{L} \ar[d] \\
  E/\mathfrak{R} \ar[r]_{f'} & E'/\mathfrak{R}'
}
\]
gives a homomorphism $w:B\to E/(\mathfrak{R}\cap\mathfrak{L})$ through which both $v'$ and $u$ factor.
Its kernel contains $\Ker(u)\cap\Ker(v')$, and is therefore open; thus $w$ is continuous.
The ideal $\mathfrak{R}\cap\mathfrak{L}$ is nilpotent, so formal smoothness of $B$ gives a continuous $A$-homomorphism $v:B\to E$ lifting $w$.
It follows at once that $f\circ v=v'$, as required.
\end{proof}

\begin{proof}[Continuation of the proof of Proposition~{\rm\sref{0.19.3.10}}]
Apply \sref{0.19.3.10.1} with
\[
  E=C/\mathfrak{Q}_{n+1},\quad E'=C/\mathfrak{Q}_n,
\]
\[
  \mathfrak{R}=(\mathfrak{J}+\mathfrak{Q}_{n+1})/\mathfrak{Q}_{n+1},\quad
  \mathfrak{R}'=(\mathfrak{J}+\mathfrak{Q}_n)/\mathfrak{Q}_n,
\]
and with $u,u',v'$ replaced respectively by $u_{n+1},u_n,v_n$.
It produces $v_{n+1}$ compatible with $v_n$.

Since $C$ is separated and complete, $C=\varprojlim C/\mathfrak{Q}_n$, and $v=\varprojlim v_n$ is a continuous $A$-homomorphism $B\to C$.
Because $\mathfrak{J}$ is closed, $\mathfrak{J}=\bigcap_n(\mathfrak{J}+\mathfrak{Q}_n)$.
Moreover, $C/\mathfrak{J}$ is metrizable and complete and the ideals $(\mathfrak{J}+\mathfrak{Q}_n)/\mathfrak{J}$ form a fundamental system of neighbourhoods of $0$; hence
\[
  C/\mathfrak{J}=\varprojlim C/(\mathfrak{J}+\mathfrak{Q}_n).
\]
Thus $\varprojlim\varphi_n$ is the canonical map $\varphi:C\to C/\mathfrak{J}$, while $\varprojlim u_n=u$.
Consequently $u=\varphi\circ v$.
\end{proof}

\begin{corollary}[19.3.11]
\label{0.19.3.11}
Let $A$ be a topological ring, $B$ a formally smooth topological $A$-algebra, $C$ a Noetherian complete local ring, $\mathfrak{J}$ an ideal distinct from $C$, $\varphi:A\to C$ a continuous homomorphism making $C$ a topological $A$-algebra.
Then every continuous $A$-homomorphism $u_0:B\to C_0=C/\mathfrak{J}$ factorizes as $B\xrightarrow{u}C\to C/\mathfrak{J}$, where $u$ is a continuous $A$-homomorphism.
\end{corollary}

\begin{proof}
All the conditions of \sref{0.19.3.10} are indeed satisfied ($\mathbf{0_I}$, 7.3.5).
\end{proof}

\begin{proposition}[19.3.12]
\label{0.19.3.12}
Let $A$ be a topological ring, $B$ a formally smooth $A$-algebra, $C$ a topological $A$-algebra, and $\mathfrak{J}$ an ideal of $C$, satisfying the following conditions:
\begin{enumerate}
  \item[$1^\circ$] There exists a fundamental system of open ideals $\mathfrak{Q}_\lambda$ of $C$, such that the $C_\lambda=C/\mathfrak{Q}_\lambda$ are Artinian rings and that $C\xrightarrow{\sim}\varprojlim C_\lambda$.
  \item[$2^\circ$] $\mathfrak{J}$ is closed in $C$ and topologically nilpotent.
  \item[$3^\circ$] The square of every open ideal of $B$ is open.
\end{enumerate}
Under these conditions, every continuous $A$-homomorphism $u:B\to C/\mathfrak{J}$ factorizes as $B\xrightarrow{v}C\to C/\mathfrak{J}$, where $v$ is a continuous $A$-homomorphism.
\end{proposition}

\begin{proof}
Let $\mathfrak{J}_\lambda=(\mathfrak{J}+\mathfrak{Q}_\lambda)/\mathfrak{Q}_\lambda$
be the canonical image of $\mathfrak{J}$ in $C_\lambda$, which is a nilpotent
ideal, and let $u_\lambda$ be the composite
$B\xrightarrow{u}C/\mathfrak{J}\to C_\lambda/\mathfrak{J}_\lambda
=C/(\mathfrak{J}+\mathfrak{Q}_\lambda)$. We first show that every
$A$-homomorphism $w:B\to C$ through which $u$ factorizes is necessarily
continuous. For every $\lambda$, there is an open ideal $\mathfrak{R}$ of $B$
such that
$u(\mathfrak{R})\subset(\mathfrak{J}+\mathfrak{Q}_\lambda)/\mathfrak{J}$,
and hence $w(\mathfrak{R})\subset\mathfrak{J}+\mathfrak{Q}_\lambda$. There is
therefore a power $\mathfrak{R}^{2^m}$ such that
$w(\mathfrak{R}^{2^m})\subset\mathfrak{Q}_\lambda$; this proves the assertion,
since $\mathfrak{R}^{2^m}$ is open in $B$. It is thus enough to find an
$A$-homomorphism $w$ with the preceding property without separately requiring
continuity.

Now
\[
  \operatorname{Hom}_A(B,C)
  =\varprojlim_\lambda\operatorname{Hom}_A(B,C_\lambda).
\]
For each $\lambda$, $\operatorname{Hom}_A(B,C_\lambda)$ is closed in the
$C_\lambda$-module $C_\lambda^B$ with its product topology, for which it is
linearly compact because $C_\lambda$ is Artinian. For each $\lambda$, let
$W_\lambda$
\oldpage[0\textsubscript{IV-1}]{86}
be the set of $w_\lambda\in\operatorname{Hom}_A(B,C_\lambda)$ such that
$u_\lambda$ factorizes as
$B\xrightarrow{w_\lambda}C_\lambda\to C_\lambda/\mathfrak{J}_\lambda$.
The set $W_\lambda$ is a closed affine subspace of the $C_\lambda$-module
$\operatorname{Hom}_A(B,C_\lambda)$ and is nonempty because $B$ is formally
smooth. In
$E_\lambda=\prod_{\mu\leq\lambda}\operatorname{Hom}_A(B,C_\mu)$, let
$U_\lambda$ be the closed affine subspace of the families
$(w_\mu)_{\mu\leq\lambda}$ such that
$w_\mu=\varphi_{\mu\lambda}\circ w_\lambda$ for $\mu\leq\lambda$ and
$w_\lambda\in W_\lambda$, where
$\varphi_{\mu\lambda}:C_\lambda\to C_\mu$ is the canonical homomorphism.
Finally let $V_\lambda$ be the closed affine subspace in
$\prod_\mu\operatorname{Hom}_A(B,C_\mu)$ equal to the product of $U_\lambda$
and the $\operatorname{Hom}_A(B,C_\mu)$ for $\mu\not\leq\lambda$. It remains to
observe that the intersection of the $V_\lambda$ is nonempty, for an element
$w$ of this intersection belongs by definition to
$\varprojlim_\lambda\operatorname{Hom}_A(B,C_\lambda)$.
\footnote{The printed proof has a gap here: sets of algebra homomorphisms are
not modules under pointwise addition, and for a general nilpotent ideal the
lift fibre need not be an affine subspace. Thus the stated linear-compactness
argument does not by itself prove that this intersection is nonempty; a
separate compatible-lifting argument is required.}

Moreover, since $\mathfrak{J}$ is closed and
$C=\varprojlim_\lambda C_\lambda$ is linearly compact,
$C/\mathfrak{J}$ identifies with
$\varprojlim_\lambda C_\lambda/\mathfrak{J}_\lambda$, and the canonical map
$\psi:C\to C/\mathfrak{J}$ with $\varprojlim_\lambda\psi_\lambda$, where
$\psi_\lambda:C_\lambda\to C_\lambda/\mathfrak{J}_\lambda$ is canonical.
As in \sref{0.19.3.10}, we conclude that $\psi\circ w=u$.
\end{proof}

\subsection{First criteria for formal smoothness}
\label{subsection:0.19.4}

\begin{proposition}[19.4.1]
\label{0.19.4.1}
Let $A$ be a topological ring, $B$ a topological $A$-algebra;
suppose that there exist two decreasing filtering families $(\mathfrak{J}_\alpha)_{\alpha\in I}$, $(\mathfrak{R}_\alpha)_{\alpha\in I}$ of ideals of $A$ and $B$ respectively, such that:
\begin{enumerate}
  \item[$1^\circ$] $(\mathfrak{J}_\alpha)$ tends to $0$ in $A$ and $(\mathfrak{R}_\alpha)$ tends to $0$ in $B$;
  \item[$2^\circ$] for all $\alpha\in I$, we have $\mathfrak{J}_\alpha B\subset\mathfrak{R}_\alpha$ (so that $B/\mathfrak{R}_\alpha$ is an $(A/\mathfrak{J}_\alpha)$-topological algebra);
  \item[$3^\circ$] for all $\alpha\in I$, $B/\mathfrak{R}_\alpha$ is a formally smooth $(A/\mathfrak{J}_\alpha)$-algebra.
\end{enumerate}
Then $B$ is a formally smooth $A$-algebra.
\end{proposition}

\begin{proof}
Let $C$ be a discrete $A$-algebra, $\mathfrak{L}$ a nilpotent ideal of $C$, and $\mathfrak{R}$ an open ideal of $B$;
by hypothesis there exists $\alpha\in I$ such that $\mathfrak{R}_\alpha\subset\mathfrak{R}$, so $B/\mathfrak{R}$ is a quotient of $B/\mathfrak{R}_\alpha$.
Every $A$-homomorphism $u':B/\mathfrak{R}\to C/\mathfrak{L}$ is also an
$(A/\mathfrak{J}_\alpha)$-homomorphism, so there exists an open ideal
$\mathfrak{R}'$ of $B$ such that
$\mathfrak{R}_\alpha\subset\mathfrak{R}'\subset\mathfrak{R}$ and the composite
$B/\mathfrak{R}_\alpha\to B/\mathfrak{R}\xrightarrow{u'}C/\mathfrak{L}$
factorizes as
$B/\mathfrak{R}_\alpha\to B/\mathfrak{R}'\to C\to C/\mathfrak{L}$, where the
second homomorphism is an $(A/\mathfrak{J}_\alpha)$-homomorphism. The conclusion
follows from the remark after \sref{0.19.3.1}.
\end{proof}

\begin{corollary}[19.4.2]
\label{0.19.4.2}
Let $A$ be a topological ring, $B$ a topological $A$-algebra, $(\mathfrak{J}_\alpha)_{\alpha\in I}$ a decreasing filtering family of ideals of $A$ tending to~$0$.
For $B$ to be a formally smooth $A$-algebra, it is necessary and sufficient that for all $\alpha\in I$, setting $A_\alpha=A/\mathfrak{J}_\alpha$, $B_\alpha=B/\mathfrak{J}_\alpha B$, $B_\alpha$ be a formally smooth $A_\alpha$-algebra.
\end{corollary}

\begin{proof}
The condition is sufficient by \sref{0.19.4.1}, and it is also necessary by \sref{0.19.3.5},~(iii).
\end{proof}

\begin{proposition}[19.4.3]
\label{0.19.4.3}
Let $A$ be a topological ring, $B$ a topological $A$-algebra.
Suppose that for every discrete $A$-algebra $C$ and every ideal $\mathfrak{J}$ of $C$ such that $\mathfrak{J}^2=0$, every continuous $A$-homomorphism $u:B\to C/\mathfrak{J}$ factorizes as $B\xrightarrow{v}C\to C/\mathfrak{J}$, where $v$ is a continuous $A$-homomorphism.
Then $B$ is a formally smooth $A$-algebra.
\end{proposition}

\begin{proof}
With the same notation, let $\mathfrak{R}$ be any nilpotent ideal of $C$,
and consider a continuous $A$-homomorphism $u':B\to C/\mathfrak{R}$.
Suppose $\mathfrak{R}^m=0$, and set $C_i=C/\mathfrak{R}^i$ for $1\leq i\leq m$, so that $C_m=C$, that $\mathfrak{R}^{i-1}/\mathfrak{R}^i$ is an ideal of square zero in $C_i$, and $C_i=C_{i+1}/\mathfrak{J}_{i+1}$;
the hypothesis then gives by induction on $i$ the existence of continuous $A$-homomorphisms $u_i:B\to C_i$ such that $u_1=u'$ and that $u_i$ factorizes as $B\xrightarrow{u_{i+1}}C_{i+1}\to C_{i+1}/\mathfrak{J}_{i+1}=C_i$;
whence the conclusion.
\end{proof}

\begin{proposition}[19.4.4]
\label{0.19.4.4}
Let $A$ be a topological ring, $B$ a topological $A$-algebra (commutative).
For $B$ to be a formally smooth $A$-algebra, it is necessary and sufficient that for
\oldpage[0\textsubscript{IV-1}]{87}
every discrete topological $B$-module $L$, annihilated by an open ideal of $B$, we have (cf.~\sref{0.18.5.1})
\[
  \operatorname{Exalcotop}_A(B,L)=0.
\]
\end{proposition}

\begin{proof}
Let $(\mathfrak{R}_\lambda)$ be a fundamental decreasing system of open ideals of $B$ and set $B_\lambda=B/\mathfrak{R}_\lambda$.
Consider a discrete topological $A$-algebra $C$ and an ideal $\mathfrak{J}$ of $C$ such that $\mathfrak{J}^2=0$, so that $C$ is an $A$-extension of $C/\mathfrak{J}$ by $\mathfrak{J}$;
suppose given an $A$-homomorphism $u:B_\lambda\to C/\mathfrak{J}$ and form the $A$-extension $E_\lambda$ of $B_\lambda$ by $\mathfrak{J}$, inverse image of $C$ by the homomorphism $u$ \sref{0.18.2.5};
it is a topological $A$-algebra for the discrete topology.
If $\operatorname{Exalcotop}_A(B,L)=0$, there exists by definition \sref{0.18.5.1} a $\mu$ such that $\mathfrak{R}_\mu\subset\mathfrak{R}_\lambda$ and the inverse image extension $E_\mu$ is $A$-trivial;
but this means \sref{0.18.1.6} that there exists a continuous $A$-homomorphism $v:B_\mu\to E_\lambda$ such that the canonical homomorphism $B_\mu\to B_\lambda$ factorizes as $B_\mu\xrightarrow{v}E_\lambda\to B_\lambda$;
we conclude immediately that $B_\mu\to B_\lambda\xrightarrow{u}C/\mathfrak{J}$ factorizes as $B_\mu\xrightarrow{v}E_\lambda\to C\to C/\mathfrak{J}$,
and this proves, by virtue of \sref{0.19.3.1} and \sref{0.19.4.3}, that $B$ is a formally smooth $A$-algebra.
The converse is immediate, applying \sref{0.19.3.1} to the case where $C$ is a topological $A$-algebra which is an $A$-extension of $B_\lambda$ by $L$, and to the identical homomorphism $B_\lambda\to B_\lambda=C/L$.
\end{proof}

When $A$ and $B$ are discrete rings, the criterion of formal smoothness \sref{0.19.4.4} reduces to
\begin{equation}
\tag{19.4.4.1}
\label{0.19.4.4.1}
\operatorname{Exalcom}_A(B,L)=0 \qquad\text{for all $B$-modules $L$};
\end{equation}
in other words, \emph{every commutative $A$-extension of $B$ by a $B$-module must be $A$-trivial}.

\begin{corollary}[19.4.5]
\label{0.19.4.5}
Let $A$ be a topological ring (commutative), $B$ a topological $A$-algebra (commutative).
\begin{enumerate}
  \item[\rm(i)] Suppose that $B$ is a formally smooth $A$-algebra. Then, for every open ideal $\mathfrak{R}$ of $B$, every $(B/\mathfrak{R})$-module $L$, and every Hochschild $A$-bilinear symmetric $2$-cocycle $f$ of $(B/\mathfrak{R})\times(B/\mathfrak{R})$ in $L$ \sref{0.18.4.3}, there exists an open ideal $\mathfrak{R}'\subset\mathfrak{R}$ such that, if $\varphi:B/\mathfrak{R}'\to B/\mathfrak{R}$ is the canonical homomorphism, $f\circ(\varphi\times\varphi)$ is a Hochschild $A$-bilinear $2$-coboundary of $(B/\mathfrak{R}')\times(B/\mathfrak{R}')$ in $L$.
  \item[\rm(ii)] If $B$ is a formally projective $A$-module \sref{0.19.2.1} and if the condition of~{\rm(i)} is satisfied, $B$ is a formally smooth $A$-algebra.
\end{enumerate}
\end{corollary}

\begin{proof}
\begin{enumerate}
  \item[\rm(i)] The $2$-cocycle $f$ defines an $A$-extension of Hochschild $C$ of $B/\mathfrak{R}$ by $L$ \sref{0.18.4.3}.
  Applying \sref{0.19.3.1} to $C$, to the ideal of square zero $L$ of $C$, and to the identical homomorphism $B/\mathfrak{R}\to B/\mathfrak{R}=C/L$, we deduce the condition~(i) by virtue of \sref{0.18.4.3}.

  \item[\rm(ii)] Let us apply the criterion \sref{0.19.4.3}, considering an open ideal $\mathfrak{J}$ of $A$, an open ideal $\mathfrak{R}$ of $B$, and an $(A/\mathfrak{J})$-extension $E$ of $B/\mathfrak{R}$ by $L$.
  As $B$ is a formally projective $A$-module, the canonical continuous $A$-linear map $\rho:B\to B/\mathfrak{R}$ factorizes as $B\xrightarrow{w}E\to B/\mathfrak{R}$, where $w$ is a continuous $A$-linear map \sref{0.19.2.1},
  which itself factorizes as $B\to B/\mathfrak{R}'\xrightarrow{w'}E$ where $\mathfrak{R}'$ is an open ideal of $B$ contained in $\mathfrak{R}$;
  replacing $\mathfrak{J}$ by a smaller ideal, one can suppose that $\mathfrak{J}B\subset\mathfrak{R}'$.
  Then the inverse image of the extension $E$ of $B/\mathfrak{R}$ by $L$, via the canonical homomorphism $B/\mathfrak{R}'\to B/\mathfrak{R}$, is an $(A/\mathfrak{J})$-extension \emph{of Hochschild} $E'$ of $B/\mathfrak{R}'$ by $L$;
  \oldpage[0\textsubscript{IV-1}]{88}
  applying~(i) to a cocycle defining this extension \sref{0.18.4.3}, we conclude that there exists an open ideal $\mathfrak{R}''\subset\mathfrak{R}'$ of $B$ such that the inverse image $E''$ of $E$ by $B/\mathfrak{R}''\to B/\mathfrak{R}$ is $A$-trivial, Q.E.D.
\end{enumerate}
\end{proof}

\begin{corollary}[19.4.6]
\label{0.19.4.6}
Let $A$ be a topological ring, $B$ a topological $A$-algebra which is a formally projective $A$-module.
Let $A'$ be an $A$-algebra equipped with the topology derived from that of $A$.
Suppose furthermore that $A'$ is a faithfully flat $A$-module, and that one of the following conditions is satisfied:
\begin{enumerate}
  \item[$1^\circ$] There exists a fundamental system $(\mathfrak{J}_\lambda)$ of open ideals of $A$ and a fundamental system $(\mathfrak{M}_\lambda)$ of open ideals of $B$, with the same index set, such that for all $\lambda$ we have $\mathfrak{J}_\lambda B\subset\mathfrak{M}_\lambda$ and $B/\mathfrak{M}_\lambda$ is a projective $(A/\mathfrak{J}_\lambda)$-module of finite type.
  \item[$2^\circ$] $A'$ is a projective $A$-module of finite type.
\end{enumerate}
Then, for $B'=B\otimes_A A'$ (equipped with the tensor product topology) to be a formally smooth $A'$-algebra, it is necessary and sufficient that $B$ be a formally smooth $A$-algebra.
\end{corollary}

\begin{proof}
The sufficiency of the condition is contained in \sref{0.19.3.5},~(iii), without any additional hypothesis on $A'$.
To prove the converse, we apply the criterion \sref{0.19.4.5};
under hypothesis~$2^\circ$, we will again denote by $(\mathfrak{M}_\lambda)$ a fundamental system of open ideals of $B$, and for each $\lambda$, by $\mathfrak{J}_\lambda$ an open ideal of $A$ such that $\mathfrak{J}_\lambda B\subset\mathfrak{M}_\lambda$;
in both cases, we set $A_\lambda=A/\mathfrak{J}_\lambda$, $B_\lambda=B/\mathfrak{M}_\lambda$, $A'_\lambda=A_\lambda\otimes_A A'$, $B'_\lambda=B_\lambda\otimes_A A'$.
Let $f$ be a Hochschild $A_\lambda$-bilinear symmetric $2$-cocycle of $B_\lambda\times B_\lambda$ in a $B_\lambda$-module $L$;
by extension of scalars, we deduce a Hochschild $A'_\lambda$-bilinear symmetric $2$-cocycle of $B'_\lambda\times B'_\lambda$ in $L'=L\otimes_A A'$.
As $B'$ is by hypothesis a formally smooth $A'$-algebra, there exists by \sref{0.19.4.5},~(i), a $\mu$ such that $\mathfrak{M}_\mu\subset\mathfrak{M}_\lambda$ and such that, if $\varphi':B'_\mu\to B'_\lambda$ is the canonical homomorphism, the image $c'$ of $f'\circ(\varphi'\times\varphi')$ in the Hochschild group $\mathrm{H}^2_{A'_\mu}(B'_\mu,L')^s$ is zero;
it is clear that if $\varphi:B_\mu\to B_\lambda$ is the canonical homomorphism, and $c$ the class of $f\circ(\varphi\times\varphi)$ in the group $\mathrm{H}^2_{A_\mu}(B_\mu,L)^s$, then $c'$ is the canonical image of $c$.
Now, if $\mathrm{P}_\bullet$ is the complex relative to the rings $A_\mu$ and $B_\mu$ defined in \sref{0.18.4.5}, the analogous complex relative to the rings $A'_\mu$ and $B'_\mu$ is evidently $\mathrm{P}_\bullet\otimes_A A'$;
under hypothesis~$1^\circ$, the construction of $\mathrm{P}_\bullet$ shows that it is a projective $A_\mu$-module of finite type.
One concludes from Bourbaki, \emph{Alg.}, chap.~II, $3^{\text{e}}$~\'{e}d., \S\,5, n$^\circ$\,3, prop.~7 that, under both hypotheses, one has, up to canonical isomorphism,
\[
  \operatorname{Hom}_{A'_\mu}(\mathrm{P}_\bullet\otimes_A A',L\otimes_A A')
  =\operatorname{Hom}_{A_\mu}(\mathrm{P}_\bullet,L)\otimes_A A'.
\]
As $A'$ is a flat $A$-module, one has by \sref{0.18.4.5}
\[
  \mathrm{H}^2_{A'_\mu}(B'_\mu,L')^s = (\mathrm{H}^2_{A_\mu}(B_\mu,L)^s)\otimes_A A'
\]
and one can therefore write $c'=c\otimes 1$;
but as $A'$ is a faithfully flat $A$-module, the hypothesis $c'=0$ implies $c=0$, which completes the proof.
\end{proof}

\begin{proposition}[19.4.7]
\label{0.19.4.7}
Let $A$ be a preadmissible topological ring, $\mathfrak{J}$ an ideal of definition of $A$ ($\mathbf{0_I}$, 7.1.2), $B$ a topological $A$-algebra which is a formally projective $A$-module.
Consider the topological quotient rings $A_0=A/\mathfrak{J}$, $B_0=B/\mathfrak{J}B$;
then, for $B$ to be a formally smooth $A$-algebra, it is necessary and sufficient that $B_0$ be a formally smooth $A_0$-algebra.
\end{proposition}

\begin{proof}
\oldpage[0\textsubscript{IV-1}]{89}
The necessity of the condition follows from \sref{0.19.4.2}.
To see that it is sufficient, we first note that, by considering a fundamental system of open neighbourhoods of $0$ in $A$ formed by ideals $\mathfrak{J}_\lambda$ contained in $\mathfrak{J}$, one can by virtue of \sref{0.19.4.2} reduce to the case where $A$ is discrete, since $B/\mathfrak{J}_\lambda B$ is a formally projective $(A/\mathfrak{J}_\lambda)$-module \sref{0.19.2.5};
by definition of a preadmissible ring, $\mathfrak{J}$ is then nilpotent.
It furthermore suffices to prove the proposition when $\mathfrak{J}^2=0$, for if $\mathfrak{J}^m=0$, one applies it successively to the rings $A_k=A/\mathfrak{J}^k$ and $B_k=B/\mathfrak{J}^k B$ and to the ideal $\mathfrak{J}^{k-1}/\mathfrak{J}^k$ of $A_k$ for $k=2,3,\ldots,m$, noting \sref{0.19.2.5} that $B_k=B\otimes_A A_k$ is a formally projective $A_k$-module.
Let us apply the criterion \sref{0.19.4.5},~(ii), considering an open ideal $\mathfrak{R}$ of $B$ and a Hochschild $A$-bilinear symmetric $2$-cocycle $f$ of $(B/\mathfrak{R})\times(B/\mathfrak{R})$ in a $(B/\mathfrak{R})$-module $L$.
Consider first the special case where $\mathfrak{J}L=0$, so that $L$ can also be viewed as a $(B_0/\mathfrak{R}B_0)$-module, and $f$ factorizes as
\[
  (B/\mathfrak{R})\times(B/\mathfrak{R}) \to (B_0/\mathfrak{R}B_0)\times(B_0/\mathfrak{R}B_0) \xrightarrow{f_0} L
\]
where $f_0$ is a Hochschild bilinear symmetric $2$-cocycle.
By hypothesis, there exists an open ideal $\mathfrak{R}'\subset\mathfrak{R}$ in $B$ and an $A_0$-linear map $g_0:B_0/\mathfrak{R}'B_0\to L$ such that $f_0(\varphi_0(x),\varphi_0(y))=xg_0(y)-g_0(xy)+g_0(x)y$ for $x,y$ in $B_0/\mathfrak{R}'B_0$, where $\varphi_0:B_0/\mathfrak{R}'B_0\to B_0/\mathfrak{R}B_0$ is the canonical homomorphism.
One concludes immediately that the composite $A$-linear map $g:B/\mathfrak{R}'\to B_0/\mathfrak{R}'B_0\xrightarrow{g_0}L$ is such that $dg=f\circ(\varphi\times\varphi)$, where $\varphi:B/\mathfrak{R}'\to B/\mathfrak{R}$ is the canonical homomorphism.

Let us now pass to the general case, and consider first the $(B/\mathfrak{R})$-module $L/\mathfrak{J}L=L'$, for which $\mathfrak{J}L'=0$;
if $f'$ is the composite $A$-bilinear map $(B/\mathfrak{R})\times(B/\mathfrak{R})\xrightarrow{f}L\to L'$, $f'$ is again a symmetric Hochschild $2$-cocycle, and by the preceding, there exists an open ideal $\mathfrak{R}'\subset\mathfrak{R}$ in $B$ and an $A$-linear map $g':B/\mathfrak{R}'\to L'$ satisfying $dg'=f'\circ(\varphi'\times\varphi')$ for the canonical map $\varphi':B/\mathfrak{R}'\to B/\mathfrak{R}$.
As $B$ is a formally projective $A$-module, there exists an $A$-linear map $g_1:B/\mathfrak{R}'_1\to L$ lifting $B/\mathfrak{R}'_1\to B/\mathfrak{R}'\xrightarrow{g'}L'$.
Consider then the Hochschild $2$-cocycle
\[
  f_1(x,y) = f(\varphi_1(x),\varphi_1(y)) - xg_1(y) + g_1(xy) - g_1(x)y,
\]
a symmetric $A$-bilinear map of $(B/\mathfrak{R}'_1)\times(B/\mathfrak{R}'_1)$ in $L$.
The choice of $g_1$ implies that $f_1$ takes its values in $\mathfrak{J}L$.
As $\mathfrak{J}(\mathfrak{J}L)=0$, one can again apply the first case, so there exists an open ideal $\mathfrak{R}''\subset\mathfrak{R}'_1$ and an $A$-linear map $g_2:B/\mathfrak{R}''\to\mathfrak{J}L$ such that
\[
  f_1(\varphi_2(x),\varphi_2(y)) = xg_2(y) - g_2(xy) + g_2(x)y
\]
in $(B/\mathfrak{R}'')\times(B/\mathfrak{R}'')$, where $\varphi_2:B/\mathfrak{R}''\to B/\mathfrak{R}'_1$ is the canonical map;
the $A$-linear map $g=g_2+g_1\circ\varphi_2:B/\mathfrak{R}''\to L$ therefore satisfies $dg=f\circ(\varphi\times\varphi)$ for the canonical map $\varphi:B/\mathfrak{R}''\to B/\mathfrak{R}$.
\end{proof}

\subsection{Formal smoothness and associated graded rings}
\label{subsection:0.19.5}

\begin{env}[19.5.1]
\label{0.19.5.1}
\oldpage[0\textsubscript{IV-1}]{90}
Let $C$ be a topological ring (commutative), $V$ a topological $C$-module, and consider the symmetric algebra $\mathbf{S}_C^*(V)=\bigoplus_{n\geq 0}\mathbf{S}_C^n(V)$, which we will canonically equip with a \emph{linear topology}, compatible with its $C$-algebra structure.
For this, consider an open submodule $U$ of $V$, and the \emph{graded ideal} $U.\mathbf{S}_C^*(V)$ it generates in $\mathbf{S}_C^*(V)$;
we take as fundamental system of neighbourhoods of $0$ in $\mathbf{S}_C^*(V)$ the set of sums $\mathfrak{a}.\mathbf{S}_C^*(V)+U.\mathbf{S}_C^*(V)$, where $\mathfrak{a}$ (resp.\ $U$) runs through a fundamental system of open ideals (resp.\ of open submodules) of $C$ (resp.\ $V$).
We note that if the topology of $V$ is less fine than the topology derived from that of $C$, one can restrict to pairs $(\mathfrak{a},U)$ such that $\mathfrak{a}V\subset U$, so $\mathfrak{a}.\mathbf{S}_C^*(V)+U.\mathbf{S}_C^*(V)=\mathfrak{a}.1_C+U.\mathbf{S}_C^*(V)$;
in this case the topology induced on each $\mathbf{S}_C^n(V)$ for $n\geq 1$ has as a fundamental system of neighbourhoods of $0$ the $U.\mathbf{S}_C^{n-1}(V)$, where $U$ runs through the open submodules of $V$;
in particular, on $V=\mathbf{S}_C^1(V)$, it coincides with the given topology (in general it is less fine than the latter).
Furthermore, in all cases, the \emph{topological algebra} $\mathbf{S}_C^*(V)$ thus defined is, for the categories of topological $C$-modules and of topological $C$-algebras, the solution of the same universal problem as for the categories of $C$-modules and $C$-algebras.
Indeed, let $E$ be a topological $C$-algebra, $u:V\to E$ a homomorphism of $C$-modules, $\mathbf{S}(u)=\widetilde{u}$ its canonical extension to $\mathbf{S}_C^*(V)$.
Suppose $u$ continuous; then, if $\mathfrak{b}$ is an open ideal of $E$, its inverse image $U=u^{-1}(\mathfrak{b})$ is an open $C$-submodule of $V$, and the image by $\widetilde{u}$ of $U.\mathbf{S}_C^*(V)$ is therefore contained in $\mathfrak{b}$;
as moreover there exists an open ideal $\mathfrak{a}$ of $C$ such that $\mathfrak{a}.E\subset\mathfrak{b}$, hence $\widetilde{u}(\mathfrak{a}.\mathbf{S}_C^*(V))\subset\mathfrak{b}$, this proves that $\widetilde{u}$ is continuous.
Conversely, if $\widetilde{u}$ is continuous, and if $\mathfrak{b}$ is an open ideal of $E$, there exists an open submodule $U$ of $V$ such that $\widetilde{u}(U.\mathbf{S}_C^*(V))\subset\mathfrak{b}$, and in particular $\widetilde{u}(U.1_C)\subset\mathfrak{b}$, that is $u(U)\subset\mathfrak{b}$, so $u$ is continuous.
We also recall that one has a canonical isomorphism of topological $C$-modules (discrete)

\[
\mathbf{S}_C^n(V)/U.\mathbf{S}_C^{n-1}(V) \xrightarrow{\sim} \mathbf{S}_C^n(V/U). \tag{19.5.1.1}
\]

\label{0.19.5.1.1}
\end{env}

\begin{env}[19.5.2]
\label{0.19.5.2}
Let $A$ be a topological ring, $B$ a topological $A$-algebra (commutative), $\mathfrak{J}$ an ideal of $B$ (not necessarily open or closed);
in all that follows, we equip $\mathfrak{J}$ and the $\mathfrak{J}^n$ ($n\geq 2$) with the topology induced from that of $B$, the quotients $C=B/\mathfrak{J}$, $\mathfrak{J}^n/\mathfrak{J}^{n+1}=\gr_\mathfrak{J}^n(B)$ with the quotient topology, so that the $\mathfrak{J}^n/\mathfrak{J}^{n+1}$ are topological $C$-modules;
the canonical injection $\mathfrak{J}/\mathfrak{J}^2\to\gr_\mathfrak{J}^*(B)$ extends to a homomorphism (not topological at first) of $C$-algebras

\[
\varphi:\mathbf{S}_C^*(\mathfrak{J}/\mathfrak{J}^2) \to \gr_\mathfrak{J}^*(B)
\]

which for each $n$ gives a surjective homomorphism of $C$-modules

\[
\varphi_n:\mathbf{S}_C^n(\mathfrak{J}/\mathfrak{J}^2) \to \mathfrak{J}^n/\mathfrak{J}^{n+1}=\gr_\mathfrak{J}^n(B). \tag{19.5.2.1}
\]

\label{0.19.5.2.1}

\oldpage[0\textsubscript{IV-1}]{91}
When $\mathbf{S}_C^n(\mathfrak{J}/\mathfrak{J}^2)$ is equipped with the topology defined in \sref{0.19.5.1}, the homomorphisms $\varphi_n$ are continuous, by virtue of the universal property \sref{0.19.5.1} of $\mathbf{S}_C^*(\mathfrak{J}/\mathfrak{J}^2)$ applied to the topological algebra $E=\gr_{\mathfrak{J}/\mathfrak{J}^{n+1}}(B/\mathfrak{J}^{n+1})$ equipped with the product topology of those of the $\mathfrak{J}^k/\mathfrak{J}^{k+1}$, and to the canonical injection $\mathfrak{J}/\mathfrak{J}^2\to E$.
Note that here the topology of $\mathfrak{J}/\mathfrak{J}^2$ is less fine than that of $C$ (this latter topology on $\mathfrak{J}/\mathfrak{J}^2$ being also the topology derived from that of $B$).
\end{env}

\begin{theorem}[19.5.3]
\label{0.19.5.3}
Let $A$ be a topological ring, $B$ a topological $A$-algebra, $\mathfrak{J}$ an ideal of $B$, $C=B/\mathfrak{J}$ the topological quotient $A$-algebra.
Suppose that the discrete $A$-algebra $C$ is formally smooth.
Then:
\begin{enumerate}
  \item[\rm(i)] If $B$ is a formally smooth $A$-algebra, $\mathfrak{J}/\mathfrak{J}^2$ is a topologically formally projective $C$-module \sref{0.19.2.1}.
  \item[\rm(ii)] If $B$ is a formally smooth $A$-algebra and a preadmissible ring ($\mathbf{0_I}$, 7.1.2), the homomorphisms $\varphi_n$ \sref{0.19.5.2} are formal bimorphisms \sref{0.19.1.2}.
  \item[\rm(iii)] Conversely, suppose that $B$ is preadmissible, that the sequence $(\mathfrak{J}^n)$ tends to $0$ in $B$, that $\mathfrak{J}/\mathfrak{J}^2$ is a formally projective $C$-module, and that the $\varphi_n$ are formal bimorphisms. Then $B$ is a formally smooth $A$-algebra.
\end{enumerate}
\end{theorem}

The proof of this theorem is long and encumbered with technical details;
we will therefore begin by proving a simpler corollary, where the guiding ideas appear more clearly;
this corollary is moreover the particular case of the theorem \sref{0.19.5.3} that will be most frequently used in the sequel.

\begin{corollary}[19.5.4]
\label{0.19.5.4}
Let $A$ be a topological ring, $B$ a topological $A$-algebra, $\mathfrak{J}$ an ideal of $B$ such that the topology of $B$ is the $\mathfrak{J}$-preadic topology.
Suppose that the discrete $A$-algebra $C=B/\mathfrak{J}$ is formally smooth.
Then the three following conditions are equivalent:
\begin{enumerate}
  \item[a)] $B$ is a formally smooth $A$-algebra.
  \item[b)] $\mathfrak{J}/\mathfrak{J}^2$ is a projective $C$-module and the canonical homomorphism
  
\[
\varphi:\mathbf{S}_C^*(\mathfrak{J}/\mathfrak{J}^2) \to \gr_\mathfrak{J}(B)
\]

  is bijective.
  \item[c)] The separated completion $\widehat{B}$ of $B$ is a topological $\widehat{A}$-algebra isomorphic to an $\widehat{A}$-algebra of the form $\widehat{B'}$, where $B'=\mathbf{S}_C^*(V)$, $V$ being a projective $C$-module and $B'$ being equipped with the ${B'}^+$-preadic topology, where ${B'}^+$ is the ideal of augmentation.
\end{enumerate}
\end{corollary}

\begin{env}[19.5.4.1]
\label{0.19.5.4.1}
Let us first prove that $a)$ implies that $\mathfrak{J}/\mathfrak{J}^2$ is a projective $C$-module.
Let $P$ and $Q$ be two $C$-modules, $u:P\to Q$ a surjective $C$-homomorphism, and $v:\mathfrak{J}/\mathfrak{J}^2\to Q$ a $C$-homomorphism.
The ring $B/\mathfrak{J}^2$ being viewed as a $B$-extension of $C$ by $\mathfrak{J}/\mathfrak{J}^2$, we deduce from $v$ a $B$-extension $E=(B/\mathfrak{J}^2)\oplus_{(\mathfrak{J}/\mathfrak{J}^2)}Q$ of $C$ by $Q$ \sref{0.18.2.8}.
As $C$ is a discrete formally smooth $A$-algebra, the extension $E$ is $A$-trivial \sref{0.19.4.4} and can therefore be identified with $\mathrm{D}_C(Q)$ \sref{0.18.2.3}.
The surjective homomorphism $u:P\to Q$ then canonically defines a surjective $A$-homomorphism $\psi:\mathrm{D}_C(P)\to\mathrm{D}_C(Q)$ \sref{0.18.2.9} whose kernel is an ideal $\mathfrak{L}$ of the extension $E'=\mathrm{D}_C(P)$, contained in $P$, and \emph{a fortiori} of square zero.
Let $f:B\to E=E'/\mathfrak{L}$ be the homomorphism defining the $B$-algebra structure of $E$: as $\mathfrak{L}$ is of square zero and $B$ is a formally smooth $A$-algebra,
\oldpage[0\textsubscript{IV-1}]{92}
$f$ factorizes as $B\xrightarrow{g}E'\to E'/\mathfrak{L}$, where $g$ is a continuous $A$-homomorphism.
The diagram

\[
\xymatrix{
    E' \ar[r]^{\psi} & E \\
    B \ar[u]^{g} \ar[r]_{\theta} & B/\mathfrak{J}^2 \ar[u]_{v'}
  }
\]

where $v'$ is deduced from $v$ \sref{0.18.2.8}, is commutative.
Moreover, the image of $\mathfrak{J}$ by $\psi\circ g$ is contained in $Q$, so the image of $\mathfrak{J}$ by $g$ is contained in $P$, and the image of $\mathfrak{J}^2$ by $g$ is zero.
Restricting $g$ and $\theta$ to $\mathfrak{J}$, one obtains a commutative diagram

\[
\xymatrix{
    P \ar[r]^{u} & Q \\
    \mathfrak{J}/\mathfrak{J}^2 \ar[u]^{w} \ar[ur]_{v}
  }
\]

which proves that the $C$-module $\mathfrak{J}/\mathfrak{J}^2$ is projective.
\end{env}

\begin{env}[19.5.4.2]
\label{0.19.5.4.2}
Let us prove secondly that $a)$ implies that $\varphi$ is bijective, which will complete the proof that $a)$ implies $b)$.
Set $E_n=B/\mathfrak{J}^{n+1}$, and denote by $F_n$ the quotient of $\mathbf{S}_C^*(\mathfrak{J}/\mathfrak{J}^2)$ by the $(n+1)$-th power of its ideal of augmentation.
As $\mathfrak{J}^{n+1}$ is nilpotent in $E_n$ and $C$ is a discrete formally smooth $A$-algebra, the identity automorphism of $C$ factorizes as $C\xrightarrow{f}E_n\to C=E_n/(\mathfrak{J}/\mathfrak{J}^{n+1})$ where $f$ is an $A$-homomorphism;
 moreover, $\mathfrak{J}/\mathfrak{J}^2$ being a projective $C$-module, the identity of $\mathfrak{J}/\mathfrak{J}^2$ factorizes as $\mathfrak{J}/\mathfrak{J}^2\xrightarrow{g}\mathfrak{J}/\mathfrak{J}^{n+1}\to\mathfrak{J}/\mathfrak{J}^2$, and from $f$ and $g$ one obtains canonically a homomorphism of $C$-algebras $\mathbf{S}_C^*(\mathfrak{J}/\mathfrak{J}^2)\to E_n$ which annihilates on the $(n+1)$-th power of the ideal of augmentation, whence, passing to the quotient, a surjective homomorphism of algebras $v:F_n\to E_n$ such that $\gr^0(v)$ and $\gr^1(v)$ are the identity automorphisms of $C$ and of $\mathfrak{J}/\mathfrak{J}^2$.
By definition of the canonical homomorphism $\varphi$, we see that $\gr^j(v)=\varphi_j$ for all $j\leq n$.
Note now that the kernel $\mathfrak{R}$ is a nilpotent ideal of $F_n$, so that $E_n$ is identified with $F_n/\mathfrak{R}$.
As $B$ is a formally smooth $A$-algebra, the canonical $A$-homomorphism $p_n:B\to E_n=B/\mathfrak{J}^{n+1}$, which is continuous, factorizes as $B\xrightarrow{w}F_n\xrightarrow{v}E_n$, where $w$ is a continuous $A$-homomorphism;
as $\gr^0(v)$ is the identity, $w(\mathfrak{J})$ is contained in the ideal of augmentation of $F_n$, hence $w(\mathfrak{J}^{n+1})=0$,
so $w$ factorizes as $B\xrightarrow{p_n}E_n\xrightarrow{w'}F_n$, and the composite $E_n\xrightarrow{w'}F_n\xrightarrow{v}E_n$ is the identity.
Moreover, since $\gr^0(v)$ and $\gr^1(v)$ are the identity automorphisms of $C$ and $\mathfrak{J}/\mathfrak{J}^2$, the same is true of $\gr^0(w')$ and $\gr^1(w')$.
Since $\gr^*(E_n)$ is generated by $\gr^1(E_n)$, the composite
\[
  \gr^j(F_n)\xrightarrow{\varphi_j}\gr^j(E_n)
  \xrightarrow{\gr^j(w')}\gr^j(F_n)
\]
is the identity for every $j\leq n$, because this is true for $j=0$ and $j=1$.
Taking $j=n$ proves that $\varphi_n$ is injective, which completes the proof that $a)$ implies $b)$.
\end{env}

\begin{env}[19.5.4.3]
\label{0.19.5.4.3}
\oldpage[0\textsubscript{IV-1}]{93}
Let us next prove that $b)$ implies $a)$.
The same reasoning as at the beginning of \sref{0.19.5.4.2} proves the existence of a surjective $A$-homomorphism of algebras $v:F_n\to E_n$ such that $\gr^j(v)=\varphi_j$ for all $j$ and the filtrations of $F_n$ and $E_n$ are finite; we conclude that $v$ is \emph{bijective} (Bourbaki, \emph{Alg.\ comm.}, chap.~III, \S\,2, n$^\circ$\,8, cor.~3 of th.~1).
Let now $G$ be a discrete topological $A$-algebra, $\mathfrak{R}$ an ideal of square zero in $G$, $f:B\to G/\mathfrak{R}$ a continuous $A$-homomorphism of algebras.
As $G$ is discrete, there exists an integer $m$ such that $f$ annihilates on $\mathfrak{J}^m$, so $f$ factorizes as $B\to E_n\xrightarrow{f_n}G/\mathfrak{R}$, where we take $n=2m$.
One has therefore by composition a continuous $A$-homomorphism $r:C\to F_n\xrightarrow{v}E_n\xrightarrow{f_n}G/\mathfrak{R}$, and as $C$ is a discrete formally smooth $A$-algebra, $r$ factorizes as $C\xrightarrow{r'}G\to G/\mathfrak{R}$, so $G$ is thus equipped with a structure of $C$-algebra.
On the other hand, the restriction to $\mathfrak{J}/\mathfrak{J}^2$ of the $C$-homomorphism $f'_n:F_n\xrightarrow{v}E_n\xrightarrow{f_n}G/\mathfrak{R}$ is a $C$-linear map $g:\mathfrak{J}/\mathfrak{J}^2\to G/\mathfrak{R}$.
As $\mathfrak{J}/\mathfrak{J}^2$ is a projective $C$-module, $g$ factorizes as $\mathfrak{J}/\mathfrak{J}^2\xrightarrow{h}G\to G/\mathfrak{R}$, where $h$ is a $C$-linear map.
By extension, one deduces a homomorphism of $C$-algebras $w:\mathbf{S}_C^*(\mathfrak{J}/\mathfrak{J}^2)\to G$, and by construction every element of degree $m$ has as its image under $w$ an element of $\mathfrak{R}$, so every element of degree $n=2m$ has zero image since $\mathfrak{R}^2=0$;
in other words, $w$ factorizes as $\mathbf{S}_C^*(\mathfrak{J}/\mathfrak{J}^2)\to F_n\xrightarrow{w'}G$.
By construction, the composite $F_n\xrightarrow{w'}G\xrightarrow{\theta}G/\mathfrak{R}$ coincides with $f'_n$ in $C$ and in $\mathfrak{J}/\mathfrak{J}^2$, hence is equal to $f'_n$.
Finally, the composite $B\to E_n\xrightarrow{w'\circ v^{-1}}G$ being equal to $f$, one sees that $B$ is a formally smooth $A$-algebra.
\end{env}

\begin{env}[19.5.4.4]
\label{0.19.5.4.4}
It remains to prove the equivalence of $a)$ and $c)$.
In the first place, $c)$ implies that $B'$ is a formally smooth $C$-algebra for the discrete topology \sref{0.19.3.2}, hence also for the ${B'}^+$-preadic topology \sref{0.19.3.8}.
As $C$ is a formally smooth $A$-algebra, $B'$ is also a formally smooth $A$-algebra \sref{0.19.3.5},~{\rm(ii)}, and finally $\widehat{B'}$ is a formally smooth $\widehat{A}$-algebra \sref{0.19.3.6}; hence $B$ is a formally smooth $A$-algebra \sref{0.19.3.6}.
It remains to see that $a)$ implies $c)$.
Note first that $C$ is a formally smooth $A$-algebra, the identical homomorphism $C\to B/\mathfrak{J}$ factorizes as $C\to\widehat{B}\to B/\mathfrak{J}$, where $C\to\widehat{B}$ is an $A$-homomorphism \sref{0.19.3.10};
hence $\widehat{B}$, and therefore all the $B/\mathfrak{J}^{n+1}$, are equipped with structures of $C$-algebras.
On the other hand, as $\mathfrak{J}/\mathfrak{J}^2$ is a projective $C$-module by $b)$, the canonical injection $\mathfrak{J}/\mathfrak{J}^2\to B/\mathfrak{J}^2$ allows one to form a projective system of $C$-homomorphisms $u_n:\mathfrak{J}/\mathfrak{J}^2\to B/\mathfrak{J}^{n+1}$ for $n\geq 1$, hence also (by the universal property of the symmetric algebra) a projective system of homomorphisms of $C$-algebras $v_n:\mathbf{S}_C^*(\mathfrak{J}/\mathfrak{J}^2)=B'\to B/\mathfrak{J}^{n+1}$;
moreover it is clear that $v_n$ annihilates on $({B'}^+)^{n+1}$ and as $\gr^0(v_n)=\varphi_0$ and $\gr^1(v_n)=\varphi_1$, one has $\gr^j(v_n)=\varphi_j$ for $0\leq j\leq n$.
As the $\varphi_j$ are isomorphisms by $b)$, and the filtrations of $B'/({B'}^+)^{n+1}$ and of $B/\mathfrak{J}^{n+1}$ are finite, one concludes that $v_n$ is \emph{bijective} for
\oldpage[0\textsubscript{IV-1}]{94}
all $n$ (Bourbaki, \emph{Alg.\ comm.}, chap.~III, \S\,2, n$^\circ$\,8, cor.~3 of th.~1); whence $c)$ by passage to the projective limit.
\end{env}

\begin{remark}[19.5.5]
\label{0.19.5.5}
Under the general hypotheses of \sref{0.19.5.4}, suppose furthermore that $\mathfrak{J}$ is a \emph{maximal} ideal of $B$, so that $C=B/\mathfrak{J}$ is a field, and that $\mathfrak{J}/\mathfrak{J}^2$ is a vector space of \emph{finite} dimension.
Then the conditions $a)$, $b)$ and $c)$ of \sref{0.19.5.4} are also equivalent to the following:
\begin{enumerate}
  \item[d)] \emph{Given} a polynomial algebra $E=C[T_1,\ldots,T_n]$ over the field $C$, and denoting by $\mathfrak{n}$ the maximal ideal of $E$ generated by the $T_i$ ($1\leq i\leq n$), then for every ideal $\mathfrak{b}$ of $E$ containing a power $\mathfrak{n}^k$, every homomorphism $v:B\to E/\mathfrak{b}$ of $C$-augmented $A$-algebras factorizes as
  \[
    B\xrightarrow{w}E/\mathfrak{n}^k\to E/\mathfrak{b},
  \]
  where $w$ is an $A$-homomorphism.
\end{enumerate}
Indeed, $\mathfrak{J}/\mathfrak{J}^2$ being here a free $C$-module, it suffices to verify that condition $d)$ implies that the canonical homomorphism
$\varphi:\mathbf{S}_C^*(\mathfrak{J}/\mathfrak{J}^2)\to\gr_\mathfrak{J}^*(B)$
is bijective.
Here $\mathbf{S}_C^*(\mathfrak{J}/\mathfrak{J}^2)$ identifies with the polynomial algebra $E=C[T_1,\ldots,T_n]$, where $n=\rg_C(\mathfrak{J}/\mathfrak{J}^2)$, and the ideal of augmentation of $E$ is the ideal $\mathfrak{n}$ generated by the $T_i$.
For every integer $k\geq 0$, formal smoothness of the discrete $A$-algebra $C$ gives, as in \sref{0.19.5.4.2}, a surjective $A$-homomorphism
\[
  f:E/\mathfrak{n}^{k+1}\to B/\mathfrak{J}^{k+1}
\]
such that $\gr^j(f)=\varphi_j$ for every $j\leq k$.
If $\mathfrak{b}/\mathfrak{n}^{k+1}=\Ker(f)$, then $B/\mathfrak{J}^{k+1}$ identifies with $E/\mathfrak{b}$, and condition $d)$ factorizes the canonical homomorphism $B\to B/\mathfrak{J}^{k+1}$ as
\[
  B\xrightarrow{w}E/\mathfrak{n}^{k+1}
  \xrightarrow{f}B/\mathfrak{J}^{k+1}.
\]
Since $\gr^0(f)$ is the identity, $w(\mathfrak{J})\subset\mathfrak{n}/\mathfrak{n}^{k+1}$, hence $w(\mathfrak{J}^{k+1})=0$.
As in \sref{0.19.5.4.2}, it follows that $\varphi_k$ is injective for every $k$, which proves the assertion.
\end{remark}

\begin{env}[19.5.6]
\label{0.19.5.6}
\emph{Proof of the theorem} \sref{0.19.5.3}. ---
Let $(\mathfrak{b}_\alpha)$ be a fundamental decreasing system of open ideals of $B$.
Set
\[
  B_\alpha=B/\mathfrak{b}_\alpha,\qquad
  C_\alpha=B/(\mathfrak{b}_\alpha+\mathfrak{J}),\qquad
  \mathfrak{J}_\alpha=(\mathfrak{b}_\alpha+\mathfrak{J})/\mathfrak{b}_\alpha,
\]
so that
\[
  C_\alpha=B_\alpha/\mathfrak{J}_\alpha
  \qquad\text{and}\qquad
  \mathfrak{J}_\alpha^n/\mathfrak{J}_\alpha^{n+1}
  =(\mathfrak{b}_\alpha+\mathfrak{J}^n)/
   (\mathfrak{b}_\alpha+\mathfrak{J}^{n+1}).
\]
The $C$-modules
\[
  ((\mathfrak{b}_\alpha\cap\mathfrak{J}^n)+\mathfrak{J}^{n+1})/
  \mathfrak{J}^{n+1}
\]
of $\mathfrak{J}^n/\mathfrak{J}^{n+1}$ form a fundamental system of open submodules of the topological $C$-module $\mathfrak{J}^n/\mathfrak{J}^{n+1}$.
Since
\[
  (\mathfrak{b}_\alpha\cap\mathfrak{J}^n)+\mathfrak{J}^{n+1}
  =\mathfrak{J}^n\cap(\mathfrak{b}_\alpha+\mathfrak{J}^{n+1}),
\]
we have
\[
\begin{aligned}
  \mathfrak{J}^n/((\mathfrak{b}_\alpha\cap\mathfrak{J}^n)+\mathfrak{J}^{n+1})
  &=
  \mathfrak{J}^n/(\mathfrak{J}^n\cap(\mathfrak{b}_\alpha+\mathfrak{J}^{n+1}))\\
  &=
  (\mathfrak{b}_\alpha+\mathfrak{J}^n)/
  (\mathfrak{b}_\alpha+\mathfrak{J}^{n+1})
  =\mathfrak{J}_\alpha^n/\mathfrak{J}_\alpha^{n+1}.
\end{aligned}
\]
\end{env}

\begin{env}[19.5.6.1]
\label{0.19.5.6.1}
Let $P,Q$ be two discrete $C_\alpha$-modules and let
$u:P\to Q$ be a surjective $C_\alpha$-homomorphism.
The discrete ring $B/(\mathfrak{b}_\alpha+\mathfrak{J}^2)$ is a $B$-extension of $C_\alpha$ by the square-zero ideal
$\mathfrak{J}_\alpha/\mathfrak{J}_\alpha^2$.
Let $v:\mathfrak{J}/\mathfrak{J}^2\to Q$ be a continuous $C$-homomorphism.
After replacing $\mathfrak{b}_\alpha$, if necessary, by a smaller open ideal of $B$, we may suppose that the kernel of $v$ contains the open $C$-submodule
$((\mathfrak{b}_\alpha\cap\mathfrak{J})+\mathfrak{J}^2)/\mathfrak{J}^2$.
Passing to the quotient gives a $C_\alpha$-homomorphism of discrete modules
$v':\mathfrak{J}_\alpha/\mathfrak{J}_\alpha^2\to Q$.
Let $E_\alpha$ be the $B_\alpha$-extension of $C_\alpha$ by $Q$ deduced from
$B/(\mathfrak{b}_\alpha+\mathfrak{J}^2)$ by means of $v'$, and let
$v'_*:B/(\mathfrak{b}_\alpha+\mathfrak{J}^2)\to E_\alpha$
be the corresponding $B_\alpha$-homomorphism \sref{0.18.2.8}.
The algebra $E_\alpha$ is a discrete topological $A$-algebra, and the canonical isomorphism
$C_\alpha\xrightarrow{\sim}E_\alpha/Q$ gives by composition a continuous
$A$-homomorphism $f:C\to C_\alpha\to E_\alpha/Q$.
Since $Q$ is square zero in $E_\alpha$ and $C$ is a formally smooth $A$-algebra,
$f$ factorizes as
$C\xrightarrow{g}E_\alpha\to E_\alpha/Q$, where $g$ is a continuous
$A$-homomorphism.
As $g$ is continuous and $E_\alpha$ discrete, there exists $\beta\geq\alpha$
such that $g$ factorizes as
$C\to C_\beta\xrightarrow{g'}E_\alpha$.
Let $E_\beta$ be the $A$-extension of $C_\beta$ by $Q$ which is the inverse image
of $E_\alpha$ under the canonical homomorphism $C_\beta\to C_\alpha$.
The existence of $g'$ (with composite $C_\beta\to E_\alpha\to C_\alpha$ equal
to the canonical homomorphism) is equivalent to $E_\beta$ being an
$A$-\emph{trivial} extension, so it may be identified with
$\mathrm{D}_{C_\beta}(Q)$.
The surjective homomorphism $u:P\to Q$ then canonically defines a

\oldpage[0\textsubscript{IV-1}]{95}

surjective homomorphism
$\mathrm{D}_{C_\beta}(P)\to\mathrm{D}_{C_\beta}(Q)$ \sref{0.18.2.9}, whose
kernel is an ideal $\mathfrak{L}$ of the extension
$E'_\beta=\mathrm{D}_{C_\beta}(P)$, contained in $P$, and \emph{a fortiori}
square zero.
Let $f':B\to E_\beta=E'_\beta/\mathfrak{L}$ be the continuous $A$-homomorphism
defining the topological $B$-algebra structure of $E_\beta$.
Since $\mathfrak{L}$ is square zero and $B$ is a formally smooth $A$-algebra,
$f'$ factorizes as
$B\xrightarrow{h'}E'_\beta\to E'_\beta/\mathfrak{L}$, where $h'$ is a
continuous $A$-homomorphism.
The construction of $E'_\beta$ gives by composition an $A$-homomorphism
$\psi:E'_\beta\to E_\beta\to E_\alpha$, and the diagram
\[
  \xymatrix{
    E'_\beta \ar[r]^\psi & E_\alpha \\
    B \ar[u]^{h'} \ar[r]_\theta &
    B/(\mathfrak{b}_\alpha+\mathfrak{J}^2) \ar[u]_{v'_*}
  }
\]
is commutative.
Moreover, the image of $\mathfrak{J}$ under $\psi\circ h'$ is contained in $Q$,
so its image under $h'$ is contained in $P$, and the image of
$\mathfrak{J}^2$ under $h'$ is zero.
Restricting $h'$ and $\theta$ to $\mathfrak{J}$ gives a commutative diagram
\[
  \xymatrix{
    P \ar[r]^u & Q \\
    \mathfrak{J}/\mathfrak{J}^2 \ar[u]^w \ar[ur]_v
  }
\]
where $w$ is continuous.
Thus $\mathfrak{J}/\mathfrak{J}^2$ is a formally projective $C$-module.
\end{env}

\begin{env}[19.5.6.2]
\label{0.19.5.6.2}
For all $n\geq 0$, we set $E_n=B/\mathfrak{J}^{n+1}$, so that $E_0=C$;
and we set

\[
E_{\alpha,n}=B/(\mathfrak{b}_\alpha+\mathfrak{J}^{n+1})=B_\alpha/\mathfrak{J}_\alpha^{n+1},
\]

a discrete quotient ring.
Likewise, in $\gr_\mathfrak{J}^n(B)=\mathfrak{J}^n/\mathfrak{J}^{n+1}$, we have seen that the $((\mathfrak{b}_\alpha\cap\mathfrak{J}^n)+\mathfrak{J}^{n+1})/\mathfrak{J}^{n+1}$ form a fundamental system of open submodules.
Consider the symmetric algebra $\mathbf{S}_{C_\alpha}^*(\mathfrak{J}_\alpha/\mathfrak{J}_\alpha^2)$;
we denote by $F_{\alpha,n}$ the quotient of $\mathbf{S}_{C_\alpha}^*(\mathfrak{J}_\alpha/\mathfrak{J}_\alpha^2)$ by the $(n+1)$-th power of its ideal of augmentation.
For a fixed $n\geq 1$, it follows from \sref{0.19.5.1.1} that the $\mathbf{S}_{C_\alpha}^n(\mathfrak{J}_\alpha/\mathfrak{J}_\alpha^2)$ are the quotients of $\mathbf{S}_C^n(\mathfrak{J}/\mathfrak{J}^2)$ by a fundamental system of open submodules within this topological $C$-module.

To abbreviate the language, we will say that for $\alpha\leq\beta$, the canonical homomorphisms $B_\beta\to B_\alpha$, $C_\beta\to C_\alpha$, $\mathfrak{J}_\beta/\mathfrak{J}_\beta^2\to\mathfrak{J}_\alpha/\mathfrak{J}_\alpha^2$, $E_{\beta,n}\to E_{\alpha,n}$, $F_{\beta,n}\to F_{\alpha,n}$, etc., are \emph{homomorphisms of transition}.
\end{env}

\begin{lemma}[19.5.6.3]
\label{0.19.5.6.3}
Suppose that $C$ is a formally smooth $A$-algebra, that the ring $B$ is preadmissible, and that the $C$-module $\mathfrak{J}/\mathfrak{J}^2$ is formally projective.
Then:
\begin{enumerate}
  \item[\rm(i)] For every $\alpha$, there exists $\beta\geq\alpha$ and a surjective $A$-algebra homomorphism
  
\[
v_{\alpha\beta}:F_{\beta,n}\to E_{\alpha,n}
\]

  such that $\gr^0(v_{\alpha\beta}):C_\beta\to C_\alpha$ and $\gr^1(v_{\alpha\beta}):\mathfrak{J}_\beta/\mathfrak{J}_\beta^2\to\mathfrak{J}_\alpha/\mathfrak{J}_\alpha^2$ are the homomorphisms of transition;
  \item[\rm(ii)] If $\beta$ and $v_{\alpha\beta}$ satisfy the conditions of~{\rm(i)}, then for every $\gamma\geq\beta$, there exists $\delta\geq\gamma$ and a surjective $A$-algebra homomorphism $v_{\gamma\delta}:F_{\delta,n}\to E_{\gamma,n}$ satisfying the conditions of~{\rm(i)} (for $\gamma$ and $\delta$) and making commutative the diagram
  
\[
\xymatrix@C=3em{
      F_{\beta,n} \ar[r]^{v_{\alpha\beta}} & E_{\alpha,n} \\
      F_{\delta,n} \ar[u] \ar[r]_{v_{\gamma\delta}} & E_{\gamma,n} \ar[u]
    }
\]

  where the vertical arrows are the homomorphisms of transition.
\end{enumerate}
\end{lemma}

\begin{env}[19.5.6.4]
\label{0.19.5.6.4}
\begin{enumerate}
  \item[\rm(i)] In the discrete topological $A$-algebra $E_{\alpha,n}$, the ideal $\mathfrak{J}_\alpha/\mathfrak{J}_\alpha^{n+1}$ is nilpotent, and the isomorphism $C_\alpha\xrightarrow{\sim}E_{\alpha,n}/(\mathfrak{J}_\alpha/\mathfrak{J}_\alpha^{n+1})$ gives by composition a continuous $A$-homomorphism $C\to C_\alpha\to E_{\alpha,n}/(\mathfrak{J}_\alpha/\mathfrak{J}_\alpha^{n+1})$;

\oldpage[0\textsubscript{IV-1}]{96}

  as $C$ is a formally smooth $A$-algebra, this homomorphism factorizes as $C\xrightarrow{f_\alpha}E_{\alpha,n}\to E_{\alpha,n}/(\mathfrak{J}_\alpha/\mathfrak{J}_\alpha^{n+1})$,
  where $f_\alpha$ is a continuous $A$-homomorphism;
  hence $\mathfrak{J}_\alpha/\mathfrak{J}_\alpha^{n+1}$ becomes, by means of $f_\alpha$, a topological $C$-module annihilated by an open ideal of $C$.
  The hypothesis that $\mathfrak{J}/\mathfrak{J}^2$ is a formally projective $C$-module then implies the existence of a continuous $C$-linear map $g_\alpha:\mathfrak{J}/\mathfrak{J}^2\to\mathfrak{J}_\alpha/\mathfrak{J}_\alpha^{n+1}$ making commutative the diagram
  
\[
\xymatrix{
      \mathfrak{J}_\alpha/\mathfrak{J}_\alpha^{n+1} \ar[r] & \mathfrak{J}_\alpha/\mathfrak{J}_\alpha^2 \\
      & \mathfrak{J}/\mathfrak{J}^2 \ar[ul]^{g_\alpha} \ar[u]
    }
\]


  As $f_\alpha$ and $g_\alpha$ are continuous, there exists $\beta\geq\alpha$ such that these homomorphisms factorize through
  
\[
C\to C_\beta\xrightarrow{f'_{\alpha\beta}}E_{\alpha,n}
\]

  
\[
\mathfrak{J}/\mathfrak{J}^2\to\mathfrak{J}_\beta/\mathfrak{J}_\beta^2\xrightarrow{g'_{\alpha\beta}}\mathfrak{J}_\alpha/\mathfrak{J}_\alpha^{n+1}
\]

  and from $f'_{\alpha\beta}$ and $g'_{\alpha\beta}$, one obtains a homomorphism of $C_\beta$-algebras
  
\[
\mathbf{S}_{C_\beta}^*(\mathfrak{J}_\beta/\mathfrak{J}_\beta^2) \to E_{\alpha,n}
\]

  which, by the definition of $g'_{\alpha\beta}$, vanishes on the $(n+1)$-th power of the augmentation ideal of $\mathbf{S}_{C_\beta}^*(\mathfrak{J}_\beta/\mathfrak{J}_\beta^2)$; passing to the quotient by this power gives the desired homomorphism $v_{\alpha\beta}$, taking into account the definitions of $f_\alpha$ and $g_\alpha$;
  the surjectivity of $v_{\alpha\beta}$ follows from that of the two homomorphisms $\gr^0(v_{\alpha\beta})$ and $\gr^1(v_{\alpha\beta})$,
  since this implies that $\gr(v_{\alpha\beta})$ is surjective (the graded algebra $\gr^*(E_{\alpha,n})$ being generated by $\gr^0(E_{\alpha,n})$ and $\gr^1(E_{\alpha,n})$), and
  as the filtrations considered are finite, one can apply Bourbaki, \emph{Alg.\ comm.}, chap.~III, \S\,2, n$^\circ$\,8, cor.~3 of th.~1.

  \item[\rm(ii)] The hypothesis that $B$ is preadmissible means that one can suppose all the $\mathfrak{b}_\alpha$ contained in the same $\mathfrak{b}_{\alpha_0}$, whose powers tend to $0$.
  This implies in particular that the kernel of the transition homomorphism $E_{\gamma,n}\to E_{\alpha,n}$, equal to $(\mathfrak{b}_\alpha+\mathfrak{J}^{n+1})/(\mathfrak{b}_\gamma+\mathfrak{J}^{n+1})$, is \emph{nilpotent};
  applying the lemma \sref{0.19.3.10.1}, one sees that one can choose $f_\gamma$ so that the diagram
  
\[
\xymatrix{
      & E_{\alpha,n} \\
      C \ar[ur]^{f_\alpha} \ar[r]_{f_\gamma} & E_{\gamma,n} \ar[u]
    }
\]

  (where the vertical arrow is the transition homomorphism) is commutative.
  The hypothesis that $\mathfrak{J}/\mathfrak{J}^2$ is a formally projective $C$-module furthermore allows one to choose $g_\gamma:\mathfrak{J}/\mathfrak{J}^2\to\mathfrak{J}_\gamma/\mathfrak{J}_\gamma^{n+1}$ so that the diagram
  \[
    \xymatrix@C=5em@R=3em{
      \mathfrak{J}_\alpha/\mathfrak{J}_\alpha^{n+1} \ar[r] &
      \mathfrak{J}_\alpha/\mathfrak{J}_\alpha^2 \\
      \mathfrak{J}_\gamma/\mathfrak{J}_\gamma^{n+1} \ar[u] \ar[r] &
      \mathfrak{J}_\gamma/\mathfrak{J}_\gamma^2 \ar[u] \\
      & \mathfrak{J}/\mathfrak{J}^2
        \ar[ul]_{g_\gamma} \ar[u]_{p_\gamma} \ar[uul]^{g_\alpha}
    }
  \]
  is commutative.
  Indeed, the image of $\mathfrak{J}/\mathfrak{J}^2$ under $(g_\alpha,p_\gamma)$ in
  $(\mathfrak{J}_\alpha/\mathfrak{J}_\alpha^{n+1})
  \times(\mathfrak{J}_\gamma/\mathfrak{J}_\gamma^2)$ is contained, by the definition of $g_\alpha$ and the relation $\beta\leq\gamma$, in the canonical image $Q$ of
  $P=\mathfrak{J}_\gamma/\mathfrak{J}_\gamma^{n+1}$ in this product module.
  Apply the definition \sref{0.19.2.1} to the surjective homomorphism $P\to Q$ and to $(g_\alpha,p_\gamma):\mathfrak{J}/\mathfrak{J}^2\to Q$.
  This choice of $f_\gamma$ and $g_\gamma$ allows $v_{\gamma\delta}$ to be constructed as in~{\rm(i)} while also making the diagram \sref{0.19.5.6.4} commutative.
\end{enumerate}
\end{env}

\oldpage[0\textsubscript{IV-1}]{97}

\begin{lemma}[19.5.6.5]
\label{0.19.5.6.5}
Suppose that the $A$-algebras $B$ and $C$ are formally smooth and that $B$ is preadmissible (so that, by \sref{0.19.5.6.1}, the conditions of \sref{0.19.5.6.3} are satisfied).
For every pair of indices $\alpha\leq\beta$ and every homomorphism
$v_{\alpha\beta}:F_{\beta,n}\to E_{\alpha,n}$ satisfying the conditions of
\sref{0.19.5.6.3},~{\rm(i)}, there exists an index $\lambda\geq\beta$ and an
$A$-algebra homomorphism
\[
  w_{\beta\lambda}:E_{\lambda,n}\to F_{\beta,n}
\]
such that:
$1^\circ$ $\gr^0(w_{\beta\lambda}):C_\lambda\to C_\beta$ and
$\gr^1(w_{\beta\lambda}):\mathfrak{J}_\lambda/\mathfrak{J}_\lambda^2
\to\mathfrak{J}_\beta/\mathfrak{J}_\beta^2$ are the transition homomorphisms;
$2^\circ$ the composite
$E_{\lambda,n}\xrightarrow{w_{\beta\lambda}}F_{\beta,n}
\xrightarrow{v_{\alpha\beta}}E_{\alpha,n}$ is the transition homomorphism.
\end{lemma}

\begin{proof}
Apply \sref{0.19.5.6.3},~{\rm(ii)}, with $\gamma=\beta$.
This gives $\delta\geq\beta$ and
$v_{\beta\delta}:F_{\delta,n}\to E_{\beta,n}$.
Recall that $v_{\beta\delta}$ is surjective.
Its kernel $\mathfrak{R}$ is nilpotent: the augmentation ideal
$F_{\delta,n}^+$ is nilpotent by definition, while the kernel of
$\gr^0(v_{\beta\delta}):C_\delta\to C_\beta$ is nilpotent because $B$ is
preadmissible.
Thus $E_{\beta,n}$ identifies with $F_{\delta,n}/\mathfrak{R}$.
Since $B$ is a formally smooth $A$-algebra, the continuous canonical
$A$-homomorphism $p_{\beta,n}:B\to E_{\beta,n}$ factorizes as
\[
  B\xrightarrow{w}F_{\delta,n}
  \xrightarrow{v_{\beta\delta}}E_{\beta,n},
\]
where $w$ is continuous.
As $F_{\delta,n}$ is discrete, there exists $\lambda\geq\delta$ such that $w$
vanishes on $\mathfrak{b}_\lambda$, and hence $w$ factorizes as
\[
  B\to B/\mathfrak{b}_\lambda
  \xrightarrow{w'_\lambda}F_{\delta,n}.
\]
Consider the composite
\[
  B/\mathfrak{b}_\lambda\xrightarrow{w'_\lambda}F_{\delta,n}
  \xrightarrow{q_{\beta\delta}}F_{\beta,n},
\]
where $q_{\beta\delta}$ is the transition homomorphism.
We have
$\gr^0(q_{\beta\delta})=\gr^0(v_{\beta\delta}):C_\delta\to C_\beta$,
the transition homomorphism.
Since
$v_{\alpha\beta}\circ q_{\beta\delta}\circ w'_\lambda$
is the canonical homomorphism
$B/\mathfrak{b}_\lambda\to
B/(\mathfrak{b}_\alpha+\mathfrak{J}^{n+1})$,
the image of
$\mathfrak{J}_\lambda=(\mathfrak{b}_\lambda+\mathfrak{J})/
\mathfrak{b}_\lambda$
under $q_{\beta\delta}\circ w'_\lambda$ is contained in the augmentation ideal
$F_{\beta,n}^+$.
Consequently the image of
$\mathfrak{J}_\lambda^{n+1}
=(\mathfrak{b}_\lambda+\mathfrak{J}^{n+1})/\mathfrak{b}_\lambda$
is zero.
Thus $q_{\beta\delta}\circ w'_\lambda$ factorizes as
\[
  B/\mathfrak{b}_\lambda\to
  B/(\mathfrak{b}_\lambda+\mathfrak{J}^{n+1})=E_{\lambda,n}
  \xrightarrow{w_{\beta\lambda}}F_{\beta,n},
\]
and
$E_{\lambda,n}\xrightarrow{w_{\beta\lambda}}F_{\beta,n}
\xrightarrow{v_{\alpha\beta}}E_{\alpha,n}$
is the transition homomorphism.
The preceding argument also shows that $\gr^0(w_{\beta\lambda})$ is the
transition homomorphism $C_\lambda\to C_\beta$, because it is the composite
\[
  B/(\mathfrak{b}_\lambda+\mathfrak{J})
  \to\gr^0(F_{\delta,n})
  \xrightarrow{\gr^0(q_{\beta\delta})}\gr^0(F_{\beta,n})
\]
and $v_{\beta\delta}\circ w$ is canonical.
Moreover,
$\gr^1(q_{\beta\delta})=\gr^1(v_{\beta\delta})$ by
\sref{0.19.5.6.3}, and the same argument shows that
$\gr^1(w_{\beta\lambda})$ is the transition homomorphism.
\end{proof}

\begin{env}[19.5.6.6]
\label{0.19.5.6.6}
\emph{Proof of} \sref{0.19.5.3},~(ii). ---
To show that $\varphi_n$ is a formal bimorphism, we apply the criterion of \sref{0.19.1.3}; the conditions of \sref{0.19.5.6.5} being satisfied by hypothesis, choose, for every $\alpha$, homomorphisms $v_{\alpha\beta}$ and $w_{\beta\lambda}$ satisfying the conclusions of that lemma.
The homomorphism
\[
  \gr^n(v_{\alpha\beta}):\gr^n(F_{\beta,n})\to\gr^n(E_{\alpha,n})
\]
is the homomorphism
\[
  \varphi_{\alpha\beta,n}:\mathbf{S}_{C_\beta}^n(\mathfrak{J}_\beta/\mathfrak{J}_\beta^2)\to\mathfrak{J}_\alpha^n/\mathfrak{J}_\alpha^{n+1}
\]
deduced from \sref{0.19.5.2.1} by passage to quotients.
Indeed, \sref{0.19.5.6.3} and the definition of $\varphi$ show inductively that $\gr^j(v_{\alpha\beta})=\varphi_{\alpha\beta,j}$ for all $j\leq n$.
Since $\varphi_n$ is surjective, it is \emph{a fortiori} a formal epimorphism.
Moreover, the composite
\[
  \gr^j(F_{\lambda,n})\xrightarrow{\varphi_{\lambda\lambda,j}}\gr^j(E_{\lambda,n})\xrightarrow{\gr^j(w_{\beta\lambda})}\gr^j(F_{\beta,n})
\]
is the transition homomorphism: this holds for $j=0,1$ by \sref{0.19.5.6.5}, and then for every $j\leq n$ because $\gr^*(E_{\lambda,n})$ is generated by $\gr^1(E_{\lambda,n})$.
After composing with the transition map to $\gr^n(F_{\alpha,n})$, the transition homomorphism $\gr^n(F_{\lambda,n})\to\gr^n(F_{\alpha,n})$ therefore factors as
\[
  \gr^n(F_{\lambda,n})\xrightarrow{\varphi_{\lambda\lambda,n}}\gr^n(E_{\lambda,n})\to\gr^n(F_{\alpha,n}),
\]
which is precisely the criterion of \sref{0.19.1.3}.
\end{env}

\begin{lemma}[19.5.6.7]
\label{0.19.5.6.7}
\oldpage[0\textsubscript{IV-1}]{98}
Suppose that $B$ is preadmissible, that $\mathfrak{J}/\mathfrak{J}^2$ is a formally projective $C$-module, that $C$ is a formally smooth $A$-algebra, and that the $\varphi_n$ are formal bimorphisms.
Then, for every pair $\alpha\leq\beta$ and every homomorphism $v_{\alpha\beta}:F_{\beta,n}\to E_{\alpha,n}$ satisfying \sref{0.19.5.6.3},~{\rm(i)}, there exists $\gamma\geq\beta$ such that, for every $\mu\geq\gamma$, there is a commutative diagram
\[
\xymatrix@C=4em@R=3em{
  F_{\beta,n} \ar[r]^{v_{\alpha\beta}} & E_{\alpha,n} \\
  & E_{\mu,n} \ar[ul]_{u_{\beta\mu}} \ar[u]_{p_{\alpha\mu}}
}
\]
of $A$-homomorphisms, where $p_{\alpha\mu}$ is the transition homomorphism.
\end{lemma}

\begin{proof}
Apply the criterion of \sref{0.19.1.3} to each $\varphi_j$, $0\leq j\leq n$.
There exist $\gamma\geq\beta$ and uniquely determined surjective $C_\gamma$-linear homomorphisms
\[
  w_{\beta\gamma}^j:\gr^j(E_{\gamma,n})\to\gr^j(F_{\beta,n})
\]
such that
\[
  \gr^j(F_{\gamma,n})\xrightarrow{\varphi_{\gamma\gamma,j}}\gr^j(E_{\gamma,n})\xrightarrow{w_{\beta\gamma}^j}\gr^j(F_{\beta,n})
\]
is the transition homomorphism.
The same $\gamma$ may be chosen for all $j$, since there are only finitely many of them.
By uniqueness and because $\varphi$ is a homomorphism of graded algebras, the family
$w_{\beta\gamma}=(w_{\beta\gamma}^j)_{0\leq j\leq n}$ is a homomorphism of graded $C_\gamma$-algebras
\[
  \gr^\bullet(E_{\gamma,n})\to\gr^\bullet(F_{\beta,n})=F_{\beta,n}.
\]
Since $\varphi_{\gamma\gamma,0}$ and $\varphi_{\gamma\gamma,1}$ are the identity homomorphisms,
$w_{\beta\gamma}^0:C_\gamma\to C_\beta$ and
$w_{\beta\gamma}^1:\mathfrak{J}_\gamma/\mathfrak{J}_\gamma^2
\to\mathfrak{J}_\beta/\mathfrak{J}_\beta^2$ are the transition homomorphisms.
The same therefore holds for
$\gr^0(v_{\alpha\beta}\circ w_{\beta\gamma})$ and
$\gr^1(v_{\alpha\beta}\circ w_{\beta\gamma})$; since
$\gr^\bullet(E_{\gamma,n})$ is generated by $\gr^1(E_{\gamma,n})$, it follows that
$\gr^j(v_{\alpha\beta}\circ w_{\beta\gamma})$ is also the transition homomorphism for $0\leq j\leq n$.

Apply now \sref{0.19.5.6.3},~(ii), to $\alpha$, $\beta$, and $\gamma$;
this gives diagram \sref{0.19.5.6.4} with $\delta\geq\gamma$.
Repeat the argument at the beginning after replacing $\alpha$ and $\beta$ by $\gamma$ and $\delta$.
One obtains $\lambda\geq\delta$ and a commutative diagram
\[
\xymatrix@C=4em@R=4em{
  \gr^\bullet(E_{\gamma,n})\ar[r]^{w_{\beta\gamma}} &
  F_{\beta,n}\ar[r]^{v_{\alpha\beta}} & E_{\alpha,n} \\
  \gr^\bullet(E_{\lambda,n})\ar[u]_{\gr^\bullet(p_{\gamma\lambda})}
    \ar[r]_{w_{\delta\lambda}} &
  F_{\delta,n}\ar[u]_{q_{\beta\delta}}\ar[r]_{v_{\gamma\delta}} &
  E_{\gamma,n}\ar[u]_{p_{\alpha\gamma}}
    \ar@{-->}[ul]_{u_{\beta\gamma}}
}
\]
whose vertical arrows are the transition homomorphisms.
It remains to prove the existence of $u_{\beta\gamma}$ making the diagram commute, and for this it plainly suffices to show
\[
  \Ker(v_{\gamma\delta})\subset\Ker(q_{\beta\delta}),
\]
since $v_{\gamma\delta}$ and $q_{\beta\delta}$ are surjective.
As $w_{\delta\lambda}$ is surjective, this relation is equivalent to
\[
\begin{split}
  \Ker(v_{\gamma\delta}\circ w_{\delta\lambda})
  &\subset\Ker(q_{\beta\delta}\circ w_{\delta\lambda})\\
  &=\Ker\bigl(w_{\beta\gamma}\circ\gr(p_{\gamma\lambda})\bigr).
\end{split}
\]
But we saw above that
$\gr(p_{\gamma\lambda})=\gr(v_{\gamma\delta}\circ w_{\delta\lambda})$, and clearly
\[
\begin{split}
  \Ker(v_{\gamma\delta}\circ w_{\delta\lambda})
  &\subset\Ker\bigl(\gr(v_{\gamma\delta}\circ w_{\delta\lambda})\bigr)\\
  &=\Ker(\gr(p_{\gamma\lambda}))
  \subset\Ker\bigl(w_{\beta\gamma}\circ\gr(p_{\gamma\lambda})\bigr).
\end{split}
\]
This proves the lemma: for every $\mu\geq\gamma$, define $u_{\beta\mu}$ as the composite
\[
  E_{\mu,n}\xrightarrow{p_{\gamma\mu}}E_{\gamma,n}
  \xrightarrow{u_{\beta\gamma}}F_{\beta,n}.
\]
\end{proof}

\begin{env}[19.5.6.8]
\label{0.19.5.6.8}
\emph{Proof of} \sref{0.19.5.3},~(iii). ---
Let $G$ be a discrete topological $A$-algebra, $\mathfrak{R}$ an ideal of square zero in $G$, and $f:B\to G/\mathfrak{R}$ a continuous $A$-homomorphism of algebras.
As $G/\mathfrak{R}$ is discrete, there exists $\alpha$ such that $f$ annihilates $\mathfrak{b}_\alpha$; moreover, as $\mathfrak{J}^m\subset\mathfrak{b}_\alpha$ for a sufficiently large $m$, $f$ factors as
\[
  B\xrightarrow{p_\alpha}E_{\alpha,n}\xrightarrow{f_{\alpha,n}}G/\mathfrak{R},
\]
where $n=2m$.
By \sref{0.19.5.6.3}, there exist $\beta\geq\alpha$ and $v_{\alpha\beta}:F_{\beta,n}\to E_{\alpha,n}$; thus
\begin{equation}
\label{0.19.5.6.9}
  F_{\beta,n}\xrightarrow{v_{\alpha\beta}}E_{\alpha,n}\xrightarrow{f_{\alpha,n}}G/\mathfrak{R}
  \tag{19.5.6.9}
\end{equation}
is a continuous $A$-homomorphism.
Since $F_{\beta,n}$ is a $C_\beta$-algebra, hence \emph{a fortiori} a discrete topological $C$-algebra, \sref{0.19.5.6.9} gives by composition a continuous $A$-homomorphism
$r:C\to F_{\beta,n}\to G/\mathfrak{R}$.
Since, on the other hand, $C$ is a formally smooth $A$-algebra,

\oldpage[0\textsubscript{IV-1}]{99}

this homomorphism factorizes as
$r:C\xrightarrow{r'}G\to G/\mathfrak{R}$, where $r'$ is a continuous $A$-homomorphism, thereby giving $G$ the structure of a discrete topological $C$-algebra.
Composing \sref{0.19.5.6.9} with
\[
  \mathfrak{J}/\mathfrak{J}^2\to\mathfrak{J}_\beta/\mathfrak{J}_\beta^2\to F_{\beta,n}
\]
gives a continuous $C$-homomorphism $g:\mathfrak{J}/\mathfrak{J}^2\to G/\mathfrak{R}$.
Since $\mathfrak{J}/\mathfrak{J}^2$ is formally projective, $g$ lifts to a continuous $C$-homomorphism $h:\mathfrak{J}/\mathfrak{J}^2\to G$.
As $G$ is discrete, there exists $\gamma\geq\beta$ such that $h$ annihilates the image of $((\mathfrak{b}_\gamma\cap\mathfrak{J})+\mathfrak{J}^2)/\mathfrak{J}^2$ and $r'$ annihilates the image of $(\mathfrak{b}_\gamma+\mathfrak{J})/\mathfrak{J}$.
Thus $h$ factors through a $C_\gamma$-homomorphism $\mathfrak{J}_\gamma/\mathfrak{J}_\gamma^2\to G$, and hence extends to
\[
  w_\gamma:\mathbf{S}_{C_\gamma}^*(\mathfrak{J}_\gamma/\mathfrak{J}_\gamma^2)\to G.
\]
Every element of degree $m$ has image in $\mathfrak{R}$, so every element of degree $n=2m$ has image zero; consequently $w_\gamma$ factors as $F_{\gamma,n}\xrightarrow{w'_\gamma}G$.

Applying \sref{0.19.5.6.7} to
$v_{\alpha\gamma}=v_{\alpha\beta}\circ q_{\beta\gamma}:F_{\gamma,n}\to E_{\alpha,n}$
gives $\delta\geq\gamma$ and $u_{\gamma\delta}:E_{\delta,n}\to F_{\gamma,n}$ such that
$v_{\alpha\gamma}\circ u_{\gamma\delta}=p_{\alpha\delta}$.
One obtains the commutative diagram
\[
\xymatrix@C=3em@R=2.5em{
  && E_{\alpha,n} \ar[r]^{f_{\alpha,n}} & G/\mathfrak{R} \\
  && F_{\beta,n} \ar[u]_{v_{\alpha\beta}} & \\
  B \ar[r]_{p_\delta} & E_{\delta,n} \ar[uur]^{p_{\alpha\delta}} \ar[r]_{u_{\gamma\delta}} & F_{\gamma,n} \ar[u]_{q_{\beta\gamma}} \ar[r]^{w'_\gamma} & G \ar[uu]
}
\]
of continuous $A$-homomorphisms.
The composite $f':B\to G$ along the bottom row satisfies $f=(G\to G/\mathfrak{R})\circ f'$, proving that $B$ is a formally smooth $A$-algebra.
The proof of \sref{0.19.5.3} is complete.
\end{env}

\begin{corollary}[19.5.7]
\label{0.19.5.7}
Let $A$ be a topological ring, $B$ a topological $A$-algebra, $(\mathfrak{b}_\lambda)$ a fundamental system of open ideals of $B$, $\mathfrak{J}$ an ideal of $B$, and $C=B/\mathfrak{J}$ the topological quotient $A$-algebra.
Set $C_\lambda=B/(\mathfrak{b}_\lambda+\mathfrak{J})$.
Suppose that:
\begin{enumerate}
  \item[$1^\circ$] for all $n$, the topology induced on $\mathfrak{J}^n$ by that of $B$ is also the topology of the $B$-module $\mathfrak{J}^n$ deduced from the topology of $B$ \sref{0.19.0.2}; this condition is satisfied in particular if $B$ is Noetherian and its topology is preadic (\textbf{0\textsubscript{I}}, 7.3.2);
  \item[$2^\circ$] $C$ is a formally smooth $A$-algebra.
\end{enumerate}
Under these conditions:
\begin{enumerate}
  \item[\rm(i)] If $B$ is a formally smooth $A$-algebra, then, for all $\lambda$, $(\mathfrak{J}/\mathfrak{J}^2)\otimes_C C_\lambda$ is a projective $C_\lambda$-module.
  \item[\rm(ii)] If $B$ is a formally smooth $A$-algebra and a preadmissible ring, then, for all $\lambda$, the canonical homomorphism
  \[
    \label{0.19.5.7.1}
    \varphi_\lambda=\varphi\otimes 1_{C_\lambda}:
    \mathbf{S}_{C_\lambda}^*((\mathfrak{J}/\mathfrak{J}^2)\otimes_C C_\lambda)
    \to \gr_\mathfrak{J}(B)\otimes_C C_\lambda
    \tag{19.5.7.1}
  \]
  is bijective.
  \oldpage[0\textsubscript{IV-1}]{100}
  \item[\rm(iii)] Conversely, suppose that $B$ is preadmissible, that the sequence $(\mathfrak{J}^n)$ tends to $0$, and that, for all $\lambda$, $(\mathfrak{J}/\mathfrak{J}^2)\otimes_C C_\lambda$ is a projective $C_\lambda$-module and the homomorphism \sref{0.19.5.7.1} is bijective.
  Then $B$ is a formally smooth $A$-algebra.
\end{enumerate}
\end{corollary}

\begin{proof}
Indeed, the topology of $\mathfrak{J}/\mathfrak{J}^2$ and those of the $\mathbf{S}_C^n(\mathfrak{J}/\mathfrak{J}^2)$ are then also deduced from that of $B$ \sref{0.19.5.1};
the conclusions are immediate consequences of \sref{0.19.5.3} and of the criteria of \sref{0.19.1.4} and \sref{0.19.2.4} specialized to this type of topology.
\end{proof}

\begin{remark}[19.5.8]
\label{0.19.5.8}
Suppose that $\mathfrak{J}/\mathfrak{J}^2$ is a $C$-module of finite type and that, for the quotient topology inherited from $B$, $C$ is a Zariski ring.
Let $\mathfrak{r}$ be an ideal of definition of $C$, so that the $\mathfrak{r}^n$ form a fundamental system of neighbourhoods of $0$ in $C$.
Then the conclusions of {\rm(i)} and {\rm(ii)} in \sref{0.19.5.7} may be replaced by the following:
\begin{enumerate}
  \item[\rm(i')] $\mathfrak{J}/\mathfrak{J}^2$ is a projective $C$-module.
  \item[\rm(ii')] The canonical homomorphism
  \[
    \varphi:\mathbf{S}_C^*(\mathfrak{J}/\mathfrak{J}^2)\to\gr_\mathfrak{J}(B)
  \]
  is bijective.
\end{enumerate}
Indeed, it is clear that {\rm(i')} implies the conclusion of {\rm(i)} in \sref{0.19.5.7}.
Conversely, if $(\mathfrak{J}/\mathfrak{J}^2)\otimes_C C_\lambda$ is a projective $C_\lambda$-module for all $\lambda$, then $(\mathfrak{J}/\mathfrak{J}^2)\otimes_C(C/\mathfrak{r}^n)$ is a projective, hence flat, $(C/\mathfrak{r}^n)$-module for every $n$;
one concludes that $\mathfrak{J}/\mathfrak{J}^2$ is a flat $C$-module (\textbf{0\textsubscript{III}}, 10.2.2), hence projective since it is finitely presented (Bourbaki, \emph{Alg. comm.}, chap.~II, \textsection5, n\textsuperscript{o}~2, cor.~2 of th.~1).
On the other hand, the $C$-modules $\mathbf{S}_C^n(\mathfrak{J}/\mathfrak{J}^2)$ and $\gr_\mathfrak{J}^n(B)$ are of finite type, and when $C$ is a Zariski ring it is equivalent to saying that $\varphi_n$ is bijective or that $\varphi_n\otimes 1_{C_\lambda}$ is bijective for all $\lambda$ (Bourbaki, \emph{Alg. comm.}, chap.~III, \textsection3, n\textsuperscript{o}~5, prop.~9);
therefore {\rm(ii)} is equivalent to {\rm(ii')}.
\end{remark}

\subsection{The case of algebras over a field}
\label{subsection:0.19.6}

\begin{theorem}[19.6.1) (Cohen]
\label{0.19.6.1}
Let $k$ be a field, $K$ an extension of $k$, $k$ and $K$ being equipped with the discrete topologies.
For $K$ to be a formally smooth $k$-algebra, it is necessary and sufficient that $K$ be a separable extension of $k$.
\end{theorem}

\begin{proof}
The necessity of the condition will be established in \sref{0.19.6.5.1} (and of course is not used until then);
we will content ourselves here with proving that the condition is \emph{sufficient}.
We distinguish two cases:

\emph{I. --- $K$ is a separable extension of finite type of $k$.}
We know then (Bourbaki, \emph{Alg.}, chap.~V, \S\,9, n$^\circ$\,3, th.~2) that there exists a pure sub-extension $K'=k(T_1,\ldots,T_n)$ of $K$ such that $K$ is a finite separable algebraic extension of $K'$.
  In view of \sref{0.19.3.5},~(iii), one can therefore reduce to the case where $K=K'$, or to the case where $K$ is algebraically finite over $k$.
In the first case, $A=k[T_1,\ldots,T_n]$ is a formally smooth $k$-algebra \sref{0.19.3.3}, hence $k(T_1,\ldots,T_n)=S^{-1}A$ where $S=A-\{0\}$ is also formally smooth \sref{0.19.3.5},~(iv).
In the second case, \emph{all} Hochschild cohomology groups $\mathrm{H}^n_k(K,L)$ are zero for every $(K,K)$-bimodule $L$:
indeed, if one considers the $k$-algebra tensor product $C=K\otimes_k K$, $L$ is a $C$-module on the left and the cohomology group $\mathrm{H}^n_k(K,L)$ is equal to $\operatorname{Ext}^n_C(K,L)$, where $K$ is also viewed as $(K,K)$-bimodule,
\oldpage[0\textsubscript{IV-1}]{101}
hence as a left $C$-module (M, IX, 4).
Now, as $K$ is a finite separable extension of $k$, we know that $K\otimes_k K$ is a direct product of extensions of $k$, one of which is $K$ itself (Bourbaki, \emph{Alg.}, chap.~VIII, \S\,8, prop.~3);
this implies that $K$ is a \emph{projective} $C$-module, whence our assertion.
Every $k$-extension of $K$ of zero kernel $L$ is therefore $k$-trivial, and \emph{a fortiori} all commutative $k$-extensions are trivial, whence the theorem in this case \sref{0.19.4.4}.

\emph{II. --- General case.} ---
With the notation of \sref{0.18.4.5}, it suffices to prove that $\Hom_K(\mathrm{H}_2(\mathrm{P}'_\bullet),L)=0$ for \emph{every} $K$-vector space $L$, which means that $\mathrm{H}_2(\mathrm{P}'_\bullet)=0$.
If $K$ is a union of a filtering family of sub-extensions $K^{(\alpha)}$ of $k$, $\mathrm{P}'_\bullet$ is the \emph{inductive limit} of the corresponding complexes $\mathrm{P}'^{(\alpha)}_\bullet$, since the $\varinjlim$ functor commutes with the tensor product in the category of $k$-modules;
by the exactness of the $\varinjlim$ functor in this category, one has $\mathrm{H}_2(\mathrm{P}'_\bullet)=\varinjlim\mathrm{H}_2(\mathrm{P}'^{(\alpha)}_\bullet)$.
As $K$ is separable, it follows that the same holds for every sub-extension of $K$ and that $K$ is a union of sub-extensions of finite type, so the first part of the proof implies that $K$ is a formally smooth $k$-algebra, Q.E.D.
\end{proof}

\begin{corollary}[19.6.2]
\label{0.19.6.2}
Let $A$ be a separated complete local ring, $K$ its residue field, $k$ a subfield of $K$ such that $K$ is a separable extension of $k$.
Then there exists a subfield $K'$ of $A$ containing $k$, such that the restriction to $K'$ of the canonical homomorphism $A\to K$ is an isomorphism of $K'$ onto $K$.
\end{corollary}

\begin{proof}
Let $\mathfrak{m}$ be the maximal ideal of $A$.
By the hypothesis and \sref{0.19.6.1}, $K$ is a formally smooth $k$-algebra;
applying \sref{0.19.3.10} (replacing $C$ by $A$ and $\mathfrak{J}$ by $\mathfrak{m}$), one sees that the identity automorphism of $K=A/\mathfrak{m}$ factorizes as $K\xrightarrow{u}A\to A/\mathfrak{m}$, where $u$ is a $k$-homomorphism, hence necessarily a $k$-isomorphism of $K$ onto a subfield $K'$ of $A$ containing $k$.
\end{proof}

\begin{corollary}[19.6.3]
\label{0.19.6.3}
Let $A$ be a Noetherian complete local ring, $\mathfrak{m}$ its maximal ideal, $k$ its residue field.
The following conditions are equivalent:
\begin{enumerate}
  \item[a)] $A$ contains a subfield.
  \item[b)] If $p$ is the characteristic of $k$, one has $pA=0$.
  \item[c)] There exists a field $K$ such that $A$ is isomorphic to a quotient ring of a formal power series ring $B=K[[T_1,\ldots,T_n]]$.
\end{enumerate}
When this is the case, one can suppose that $A$ is isomorphic to $B/\mathfrak{b}$, where $\mathfrak{b}$ is contained in the square of the maximal ideal of $B$.
\end{corollary}

\begin{proof}
It is immediate that $c)$ implies $a)$, since $A$ is $\neq 0$ because it is a local ring;
as $K$ is a field and the composite homomorphism $K\xrightarrow{\varphi}B\to A$ (where $\varphi$ is the canonical injection) is not zero, it is injective.
To see that $a)$ implies $b)$, it suffices to remark that if $K$ is a subfield of $A$, $K$ is isomorphic to a subfield of $k$ and has therefore the same characteristic $p$;
as $p.1=0$ in $K$, hence in $A$, one has $pA=0$.
Conversely, $b)$ implies $a)$, for the composite of the canonical homomorphisms $\mathbf{Z}\xrightarrow{j}A\to k$ has kernel $p\mathbf{Z}$, and one has $j(p\mathbf{Z})=0$, hence $\Ker(j)=p\mathbf{Z}$ and $A$ contains a ring $A'$ isomorphic to $\mathbf{Z}/p\mathbf{Z}$ and not meeting $\mathfrak{m}$;
if $p>0$, $A'$ is a field; otherwise,

\oldpage[0\textsubscript{IV-1}]{102}

every element of $A'$ being invertible in $A$, the field of fractions of $A'$ (isomorphic to $\mathbf{Q}$) is also contained in $A$.

Let us prove finally that $a)$ implies $c)$;
the condition $a)$ implies that $A$ contains a prime subfield $k_0$, isomorphic to the prime subfield of $k$, and whose $k$ is therefore a separable extension.
Applying \sref{0.19.6.2}, one sees that there exists an isomorphism $f$ of $k$ onto a subfield $K$ of $A$.
On the other hand, let $(x_i)_{1\leq i\leq n}$ be a system of elements of $\mathfrak{m}$ such that the classes mod.\ $\mathfrak{m}^2$ of the $x_i$ form a basis (over $k$) of $\mathfrak{m}/\mathfrak{m}^2$.
Since $A$ is complete, there exists a continuous ring homomorphism $u:k[[T_1,\ldots,T_n]]\to A$ such that $u$ equals $f$ on $k$ and $u(T_i)=x_i$ for all $i$, and this homomorphism is surjective by virtue of the choice of the $x_i$ (Bourbaki, \emph{Alg.\ comm.}, chap.~III, \S\,2, n$^\circ$\,9, prop.~11).
\end{proof}

\begin{theorem}[19.6.4]
\label{0.19.6.4}
Let $k$ be a field, $A$ a local Noetherian $k$-algebra, $\mathfrak{m}$ its maximal ideal, $K=A/\mathfrak{m}$ its residue field, $A$ being equipped with the $\mathfrak{m}$-preadic topology.
Suppose that $K$ is a separable extension of $k$.
Then the following conditions are equivalent:
\begin{enumerate}
  \item[a)] $A$ is a formally smooth $k$-algebra.
  \item[b)] The completion $\widehat{A}$ of $A$ is $k$-isomorphic to a formal power series ring $K[[T_1,\ldots,T_n]]$ (where the $k$-algebra structure is defined by the composite homomorphism $k\xrightarrow{\varphi}K\xrightarrow{\psi}K[[T_1,\ldots,T_n]]$, where $\psi$ is the canonical injection and $\varphi$ the homomorphism defining the extension structure of $K$).
  \item[b')] $\widehat{A}$ is a formal power series ring isomorphic to $K[[T_1,\ldots,T_n]]$.
  \item[c)] $A$ is a regular local ring.
\end{enumerate}
\end{theorem}

\begin{proof}
The fact that $a)$ implies $b)$ is the particular case of \sref{0.19.5.4} (equivalence of $a)$ and $c)$) obtained by replacing $A$ by $k$, $B$ by $A$, and $\mathfrak{J}$ by $\mathfrak{m}$, taking into account \sref{0.19.6.1}.
It is clear that $b)$ implies $b'$) and $b')$ implies $c)$ by \sref{0.17.1.1}.
Finally, if $c)$ holds, it follows from \sref{0.17.1.1},~$a$), and from \sref{0.19.5.4} (equivalence of $b)$ and $a)$), applied with $A=k$, $B=A$, $\mathfrak{J}=\mathfrak{m}$, and $C=K=A/\mathfrak{m}$, that $A$ is a formally smooth $k$-algebra.
\end{proof}

\begin{corollary}[19.6.5]
\label{0.19.6.5}
Let $k$ be a field, $A$ a local Noetherian $k$-algebra, $\mathfrak{m}$ its maximal ideal, $A$ being equipped with the $\mathfrak{m}$-preadic topology.
Suppose that $A$ is a formally smooth $k$-algebra;
then $A$ is \emph{geometrically regular} over $k$, in other words (cf.~$\mathbf{IV}$, 6.7.6), for every finite extension $k'$ of $k$, the semi-local ring $A'=A\otimes_k k'$ is regular \sref{0.17.3.6}.
Furthermore, if $K$ is the residue field of $A$, $\widehat{A}$ is isomorphic to a formal power series ring $K[[T_1,\ldots,T_n]]$.
\end{corollary}

\begin{proof}
As $\widehat{A}'=\widehat{A}\otimes_k k'$, it follows from \sref{0.19.3.6} and \sref{0.17.1.5} (applied to the maximal ideals of $A'$) that one can reduce to the case where $A$ is complete;
then $A'$ is a formally smooth $k'$-algebra \sref{0.19.3.5},~(iii) and is moreover a direct product of $k'$-local algebras, which are also formally smooth \sref{0.19.3.5},~(v).
One is therefore reduced to proving that $A$ is \emph{regular}.
Let $k_0$ be the prime subfield of $k$;
as $k$ is a separable extension of $k_0$, it is a formally smooth $k_0$-algebra \sref{0.19.6.1}, and by the hypothesis, so is $A$ \sref{0.19.3.5},~(ii).
As the residue field $K$ of $A$ is a separable extension of $k_0$, one can then apply \sref{0.19.6.4} to $A$ and $k_0$, whence the conclusion.
\end{proof}

\oldpage[0\textsubscript{IV-1}]{103}

\begin{corollary}[19.6.5.1]
\label{0.19.6.5.1}
Let $A$ be a Noetherian local ring, $k$ a subfield of $A$ such that $A$ is a formally smooth $k$-algebra (for its preadic topology).
Then every field $K$ such that $k\subset K\subset A$ is a separable extension of $k$.
\end{corollary}

\begin{proof}
Indeed, for every finite extension $k'$ of $k$, the ring $K\otimes_k k'$ is identified with a subring of $A'=A\otimes_k k'$;
as $A'$ is a regular ring, hence reduced, so is $K\otimes_k k'$, which proves that $K$ is a separable extension of $k$ (Bourbaki, \emph{Alg.}, chap.~VIII, \S\,7, n$^\circ$\,3, th.~1).
\end{proof}

We note that this proves the necessity of the condition in the statement of \sref{0.19.6.1}.

\begin{remark}[19.6.5.2]
\label{0.19.6.5.2}
If $K$ is a \emph{non-separable} extension of $k$, the formal power series ring $K[[T_1,\ldots,T_n]]$ is therefore \emph{not} a formally smooth $k$-algebra (for the usual $k$-algebra structure recalled in \sref{0.19.6.4}).
On the other hand, there exist formally smooth $k$-algebras $A$ which are complete Noetherian local rings whose residue field $K$ is a non-separable extension of $k$ (for example the completions of the localizations of $B=k[T_1,\ldots,T_n]$ at maximal ideals $\mathfrak{n}$ such that $B/\mathfrak{n}$ is a finite non-separable algebraic extension of $k$);
such a ring cannot therefore be $k$-isomorphic to $K[[T_1,\ldots,T_n]]$ even though it is \emph{isomorphic} to it by virtue of \sref{0.19.6.5}.

One can in certain cases sharpen the corollary \sref{0.19.6.5}:
\end{remark}

\begin{proposition}[19.6.6]
\label{0.19.6.6}
Let $k$ be a field, $A$ a local Noetherian $k$-algebra, $\mathfrak{m}$ its maximal ideal, $A$ being equipped with the $\mathfrak{m}$-preadic topology.
Suppose that the residue field $K=A/\mathfrak{m}$ of $A$ is such that there exists a finite radicial extension $k_0$ of $k$ for which the residue field of the local ring $K\otimes_k k_0$ is a separable extension of $k_0$ (we will express this property later by saying that $K$ is of \emph{finite radicial multiplicity} over $k$ ($\mathbf{IV}$, 4.7.4) and we will see that this condition is always satisfied when $K$ is an extension of finite type of $k$).
Then the following conditions are equivalent:
\begin{enumerate}
  \item[a)] $A$ is a formally smooth $k$-algebra.
  \item[b)] There exists a finite radicial extension $k'$ of $k$ such that $\widehat{A}\otimes_k k'$ is $k'$-isomorphic to a formal power series algebra $K'[[T_1,\ldots,T_n]]$, where $K'$ is a separable extension of $k'$.
  \item[b')] There exists a finite extension $k'$ of $k$ such that the semi-local ring $\widehat{A}\otimes_k k'$ has at least one local component which is $k'$-isomorphic to a formal power series algebra $K'[[T_1,\ldots,T_n]]$, where $K'$ is a separable extension of $k'$.
  \item[c)] $A$ is geometrically regular over $k$ \sref{0.19.6.5}.
\end{enumerate}
\end{proposition}

\begin{proof}
Note first that if $k'$ is a finite radicial extension of $k$, there is only a single maximal ideal of $A'=A\otimes_k k'$ above $\mathfrak{m}$, formed of elements whose $p^h$-th power ($p$ being the characteristic of $k$) is in $\mathfrak{m}$ for a suitable $h$ (Bourbaki, \emph{Alg.\ comm.}, chap.~V, \S\,2, n$^\circ$\,3, lemma~4);
$A'$ is therefore a local ring, and so is
$K\otimes_k k'=(A\otimes_k k')/(\mathfrak{m}\otimes_k k')$; moreover, these two local rings have the same residue field.
We recall, on the other hand, that if $K$ is a separable extension of $k$, then, for every finite extension $k''$ of $k$, $K\otimes_k k''$ is a direct product of fields (Bourbaki, \emph{Alg.}, chap.~VIII, \S\,7, n$^\circ$\,3, cor.~1 of th.~1);
therefore $\mathfrak{m}\otimes_k k''$ is the radical of $A\otimes_k k''$, and the field components of $K\otimes_k k''$ are the residue fields at the maximal ideals of $A\otimes_k k''$;

\oldpage[0\textsubscript{IV-1}]{104}

furthermore, these fields are separable extensions of $k''$, since for every extension $k_1$ of $k''$,
$K\otimes_k k_1=(K\otimes_k k'')\otimes_{k''}k_1$ has no radical (\emph{loc.\ cit.}, th.~1).
Applying these remarks to the field $K$ of finite radicial multiplicity over $k$, one sees that for \emph{every} finite extension $k''$ of $k_0$, $(K\otimes_k k'')_{\mathrm{red}}$ is separable over $k''$.

Let us pass to the proof of \sref{0.19.6.6}.
It is trivial that $b)$ implies $b')$.
Let us show that $b')$ implies $a)$.
Set $B'=K'[[T_1,\ldots,T_n]]$.
The hypothesis that $K'$ is separable over $k'$ entails, by the preceding remarks, that for every finite extension $k''$ of $k'$, the components of the complete semi-local ring
$B'\otimes_{k'}k''$ (equal to the formal power series ring
$(K'\otimes_{k'}k'')[[T_1,\ldots,T_n]]$) are the rings
$K_i''[[T_1,\ldots,T_n]]$, where the $K_i''$ are the field components of
$K'\otimes_{k'}k''$.
Since the local components of $B'\otimes_{k'}k''$ are also local components of
\[
  \widehat{A}\otimes_k k''
  =(\widehat{A}\otimes_k k')\otimes_{k'}k'',
\]
the hypothesis $b')$ remains true when $k'$ is replaced by any one of its finite extensions.
In particular, one can suppose $k'$ \emph{quasi-Galois}, hence a Galois extension of a radicial extension $k_0'$ of $k$;
$A'=\widehat{A}\otimes_k k'$ can then be written
$A_0'\otimes_{k_0'}k'$, where $A_0'=\widehat{A}\otimes_k k_0'$ is a complete local ring.
We know (Bourbaki, \emph{Alg.}, chap.~VIII, \S\,8, prop.~4) that $A'$ is a direct product of local rings isomorphic to $A_0'$ as $k_0'$-algebras.
Now one of these rings is, by hypothesis, $k'$-isomorphic to a formal power series ring
$K'[[T_1,\ldots,T_n]]$, where $K'$ is a separable extension of $k'$, hence also a separable extension of $k_0'$.
In view of \sref{0.19.6.4}, $A_0'$ is therefore a formally smooth $k_0'$-algebra.
But as $k_0'$ is a $k$-module which is faithfully flat, projective and of finite type, one concludes by \sref{0.19.4.7} that $\widehat{A}$ is a formally smooth $k$-algebra.

We have already seen \sref{0.19.6.5} that $a)$ implies $c)$; it remains to verify that $c)$ implies $b)$.
If $c)$ holds, then, for every finite radicial extension $k'$ of $k$ containing $k_0$, $A\otimes_k k'$ is a regular local ring whose residue field $K'$ is a separable extension of $k'$.
It follows from \sref{0.19.6.4} that its completion
$\widehat{A}\otimes_k k'$ is $k'$-isomorphic to a formal power series ring
$K'[[T_1,\ldots,T_n]]$.
\end{proof}

\begin{remarks}[19.6.7]
\label{0.19.6.7}
\begin{enumerate}
  \item[\rm(i)] We note that the hypothesis that $K$ is of finite radicial multiplicity over $k$ was only used in the proof of the implication $b')\Rightarrow a)$.
  \item[\rm(ii)] We will prove later (22.5.8) that $a)$ and $c)$ are equivalent, \emph{without} any hypothesis on the extension $K$ of $k$.
\end{enumerate}
\end{remarks}

\subsection{The case of local homomorphisms; existence and uniqueness theorems}
\label{subsection:0.19.7}

In this section, whenever a semi-local ring is considered as a topological ring, it is always understood that it is endowed with its $\mathfrak{r}$-preadic topology, where $\mathfrak{r}$ is its radical. Every local homomorphism of local rings is thus automatically continuous.

\begin{theorem}[19.7.1]
\label{0.19.7.1}
Let $A$, $B$ be two local Noetherian rings, $\mathfrak{m}$, $\mathfrak{n}$ their respective maximal ideals, $k=A/\mathfrak{m}$ the residue field of $A$; suppose that $A$ and $B$ are equipped respectively with the $\mathfrak{m}$-preadic and $\mathfrak{n}$-preadic topologies.
Let $\varphi:A\to B$ be a local homomorphism, and set $B_0=B\otimes_A k$.
The following properties are equivalent:
\begin{enumerate}
  \item[\rm a)] $B$ is a formally smooth $A$-algebra.
  \oldpage[0\textsubscript{IV-1}]{105}
  \item[\rm b)] $B$ is a flat $A$-module and $B_0=B/\mathfrak{m}B$ (equipped with the quotient topology) is a formally smooth $k$-algebra.
\end{enumerate}
\end{theorem}

\begin{proof}
The proof is carried out in several steps.

\medskip
\noindent\phantomsection\label{0.19.7.1.1}\textbf{(19.7.1.1)}\enspace
Let us first show that $b)$ implies $a)$; we will apply the criterion \sref{0.19.4.7}, with $\mathfrak{J}=\mathfrak{m}$; by virtue of the second hypothesis in $b)$, it suffices to show that $B$ is a \emph{formally projective} $A$-module.
The hypothesis implies that for all $h>0$, $B/\mathfrak{m}^h B$ is a flat $(A/\mathfrak{m}^h)$-module ($\mathbf{0_{III}}$, 10.2.1); as the $\mathfrak{m}^h$ form a fundamental system of neighbourhoods of $0$ in $A$ and $(B/\mathfrak{m}^h B)\otimes_{A/\mathfrak{m}^h} k = B_0$, we can replace $A$ and $B$ by $A/\mathfrak{m}^h$ and $B/\mathfrak{m}^h B$ respectively, and thus suppose that $A$ is \emph{Artinian} (hence discrete).
As $B_0$ is a $k$-algebra formally smooth, it is a regular ring \sref{0.19.6.5}; let $(x_i^0)_{1\leq i\leq n}$ be a regular system of parameters for $B_0$ \sref{0.17.1.6},
and for each $i$, let $x_i\in B$ be such that $x_i^0$ is its image in $B_0=B/\mathfrak{m}B$; as the $x_i^0$ generate the maximal ideal $\mathfrak{n}_0=\mathfrak{n}/\mathfrak{m}B$ of $B_0$, the ideals $\mathfrak{Q}_m=\sum_{i=1}^n (x_i^0)^m B_0$ (for $m>0$) form a fundamental system of neighbourhoods of $0$ in $B_0$, since $\mathfrak{Q}_m$ is evidently contained in $\mathfrak{n}_0^m$, and on the other hand contains $\mathfrak{n}_0^{mn}$.
Set $\mathfrak{J}_m=\sum_{i=1}^n x_i^m B$ for all $m>0$;
it is clear that $\mathfrak{n}=\mathfrak{J}_1+\mathfrak{m}B$; as there exists $h>0$ such that $\mathfrak{m}^h=0$, one has $\mathfrak{n}^m\subset\mathfrak{J}_1^{m-h}\subset\mathfrak{n}^{m-h}$ for $m>h$, and as we have seen that $\mathfrak{J}_m\supset\mathfrak{J}_1^{mn}$, the $\mathfrak{J}_m$ form a fundamental system of neighbourhoods of $0$ in $B$.
The whole argument then amounts to proving that the $B/\mathfrak{J}_m$ are \emph{free} $A$-modules, and it amounts to the same to verify that these are \emph{flat} $A$-modules ($\mathbf{0_{III}}$, 10.1.3).
Now, the hypothesis that $(x_i^0)$ is a $B_0$-regular sequence of elements of the maximal ideal of $B_0$ implies the same property for the sequence $(x_i^0)^m$ ($1\leq i\leq n$) for all $m>0$ \sref{0.15.1.20}; the conclusion thus follows from \sref{0.15.1.16}, $b)$ and $c)$.
\end{proof}

\begin{lemma}[19.7.1.2]
\label{0.19.7.1.2}
Let $A$ be a topological ring, $B$, $C$ two topological $A$-algebras that are local Noetherian rings.
Suppose further that $C$ is complete and that the residue field $B/\mathfrak{m}$ of $B$ is an $A$-module of finite type.
Let $E$ be the completed tensor product $B\widehat{\otimes}_A C$.
Then:
\begin{enumerate}
  \item[\rm(i)] $E$ is a complete semi-local Noetherian ring.
  \item[\rm(ii)] The ideal $\mathfrak{m}E$ is contained in the radical of $E$, and for all $h>0$, $E/\mathfrak{m}^h E$ is isomorphic to $(B/\mathfrak{m}^h)\otimes_A C$.
  \item[\rm(iii)] If $C$ is a flat $A$-module, $E$ is a flat $B$-module.
\end{enumerate}
\end{lemma}

\begin{proof}
By definition, $E$ is the separated completion of the tensor product $B\otimes_A C$ for the topology defined by the ideals $\operatorname{Im}(\mathfrak{m}^i\otimes_A C)+\operatorname{Im}(B\otimes_A \mathfrak{n}'{}^j)$ ($\mathbf{0_I}$, 7.7.5).
If one sets $\mathfrak{r}=\operatorname{Im}(\mathfrak{m}\otimes_A C)+\operatorname{Im}(B\otimes_A \mathfrak{n}')$, one has $\mathfrak{r}^{2i}\subset\operatorname{Im}(\mathfrak{m}^i\otimes_A C)+\operatorname{Im}(B\otimes_A \mathfrak{n}'{}^i)\subset\mathfrak{r}^i$, so $E$ is also the separated completion of $B\otimes_A C$ for the $\mathfrak{r}$-preadic topology.
By hypothesis, $(B\otimes_A C)/\mathfrak{r}=(B/\mathfrak{m})\otimes_A (C/\mathfrak{n}')$ is a $(C/\mathfrak{n}')$-module of finite type, hence an Artinian ring;
moreover, $\mathfrak{r}/\mathfrak{r}^2$, being a quotient of $((\mathfrak{m}/\mathfrak{m}^2)\otimes_A C)\oplus(B\otimes_A (\mathfrak{n}'/\mathfrak{n}'{}^2))$, is a $(B\otimes_A C)$-module of finite type; applying ($\mathbf{0_I}$, 7.2.11), one sees that $E$ is a \emph{Noetherian} ring; moreover, $E/\mathfrak{r}E$, isomorphic to $(B\otimes_A C)/\mathfrak{r}$, being Artinian, $E$ is \emph{semi-local}.
Note now that $E$, which is isomorphic to $\varprojlim_{i,j}((B/\mathfrak{m}^i)\otimes_A(C/\mathfrak{n}'{}^j))$, is also isomorphic, by the theorem of the double projective limit, to $\varprojlim_i(\varprojlim_j((B/\mathfrak{m}^i)\otimes_A (C/\mathfrak{n}'{}^j)))$; but $\varprojlim_j((B/\mathfrak{m}^i)\otimes_A (C/\mathfrak{n}'{}^j))$
\oldpage[0\textsubscript{IV-1}]{106}
is the separated completion of $(B/\mathfrak{m}^i)\otimes_A C$, and as $C$ is complete, this is none other than $(B/\mathfrak{m}^i)\otimes_A C$ itself, since $B/\mathfrak{m}^i$ is an $A$-module of finite type (as $\mathfrak{m}/\mathfrak{m}^i$ is a $(B/\mathfrak{m})$-module of finite type ($\mathbf{0_I}$, 7.3.6)).
One therefore has $E=\varprojlim_i((B/\mathfrak{m}^i)\otimes_A C)$.
For every integer $h>0$, $\mathfrak{m}^h E$, being an ideal of $E$, is closed in $E$ ($\mathbf{0_I}$, 7.3.5), hence complete, and on the other hand it is evidently dense in $\varprojlim_i(\operatorname{Im}(\mathfrak{m}^h/\mathfrak{m}^{h+i})\otimes_A C)$, hence equal to this projective limit.
Moreover, all the projective systems considered are defined by \emph{surjective} homomorphisms; it follows therefore from ($\mathbf{0_{III}}$, 13.2.2) that $E/\mathfrak{m}^h E$ is isomorphic to $(B/\mathfrak{m}^h)\otimes_A C = (B\otimes_A C)/\operatorname{Im}(\mathfrak{m}^h\otimes_A C)$.
In particular $\operatorname{Im}(\mathfrak{m}\otimes_A C)\subset\mathfrak{r}$, which shows that $\mathfrak{m}E$ is contained in $\mathfrak{r}E$, hence in the radical of $E$.
Finally, the hypothesis that $C$ is a flat $A$-module implies that $(B/\mathfrak{m}^h)\otimes_A C = E/\mathfrak{m}^h E$ is a flat $(B/\mathfrak{m}^h)$-module for all $h>0$: since $B$ and $E$ are Noetherian and $\mathfrak{m}E$ is contained in the radical of $E$, it follows from ($\mathbf{0_{III}}$, 10.2.2) that $E$ is a flat $B$-module.
\end{proof}

\begin{lemma}[19.7.1.3]
\label{0.19.7.1.3}
Let $A$ be a local Noetherian ring, $\mathfrak{m}$ its maximal ideal, $k$ its residue field, $B_0$ a local Noetherian $k$-algebra; suppose that $B_0$ is complete and regular.
Then there exists a topological $A$-algebra $B$ that is a complete local Noetherian ring, a flat $A$-module, and such that $B_0$ is $k$-isomorphic to $B\otimes_A k = B/\mathfrak{m}B$.
\end{lemma}

\begin{proof}
As $\widehat{A}$ is flat over $A$ and has the same residue field, we can restrict to the case where $A$ is \emph{complete}.

Let $K$ be the residue field of $B_0$, and let us distinguish two cases:

\medskip
\noindent\emph{Case~I}: $K$ is a \emph{separable} extension of $k$.
By virtue of \sref{0.19.6.4}, $B_0$ is $k$-isomorphic to a formal power series ring $K[[T_1,\ldots,T_n]]$.
When $B_0=K$, the lemma has already been proved ($\mathbf{0_{III}}$, 10.3.1); let $C$ be a complete local Noetherian ring that is a flat $A$-module and such that $C\otimes_A k$ is isomorphic to $K$.
For $n\geq 1$, it suffices to take (using the preceding notation) $B=C[[T_1,\ldots,T_n]]$; indeed it is known (Bourbaki, \emph{Alg.\ comm.}, chap.~III, \S\,3, n$^\circ$\,4, cor.~3 of th.~1) that $B$ is a flat $C$-module, hence also a flat $A$-module, and on the other hand, it is immediate that $C[[T_1,\ldots,T_n]]\otimes_A k$ is isomorphic to $(C/\mathfrak{m}C)[[T_1,\ldots,T_n]]=B_0$.

\medskip
\noindent\emph{Case~II}: $K$ is of \emph{characteristic} $p>0$, and is therefore of the same characteristic as $k$.
Denote by $P$ the prime field $\mathbf{F}_p$, and by $W(P)$ the complete local ring of $p$-adic numbers $\mathbf{Z}_p$, which is a complete discrete valuation ring (hence \emph{regular}), with $P$ as residue field.
Let us first show that there exists a continuous ring homomorphism $W(P)\to A$, thus making $A$ into a $W(P)$-\emph{topological algebra}.
Indeed, if $j:\mathbf{Z}\to A$ is the canonical homomorphism, one has $j(p\mathbf{Z})\subset\mathfrak{m}$ by hypothesis, whence $j^{-1}(\mathfrak{m})=p\mathbf{Z}$, and $j$ therefore factors as $\mathbf{Z}\to\mathbf{Z}_{p\mathbf{Z}}\xrightarrow{f}A$, where $f$ is a local homomorphism, hence continuous, which (since $A$ is complete) extends by continuity to the desired homomorphism $W(P)\to A$.

As $k$ is a separable extension of $P$, case~I) shows that there is a local homomorphism $W(P)\to W(k)$, where $W(k)$ is a complete local Noetherian ring and a flat $W(P)$-module, such that $W(k)\otimes_{W(P)}P$ is isomorphic to $k$.
Moreover, as the uniformiser $p$ of $W(P)$ is a $W(k)$-regular element by flatness ($\mathbf{0_I}$, 6.3.4), and as
\oldpage[0\textsubscript{IV-1}]{107}
$W(k)/pW(k)=k$, the maximal ideal of $W(k)$ is $pW(k)$, which implies that this ring is a \emph{complete discrete valuation ring} (Bourbaki, \emph{Alg.\ comm.}, chap.~VI, \S\,3, n$^\circ$\,5, prop.~9).
By \sref{0.19.7.1.1} (as $k$ is separable over $P$, hence a formally smooth $P$-algebra \sref{0.19.6.1}), one sees that $W(k)$ is a $W(P)$-algebra \emph{formally smooth}.
The continuous $W(P)$-homomorphism $W(k)\to k$ therefore factors as $W(k)\to A\to k$ \sref{0.19.3.11},
which allows us to consider $A$ as a $W(k)$-topological algebra.
Applying now case~I) to $B_0$ considered as a $P$-algebra and to $W(P)$, one sees that there exists a $W(P)$-algebra $B_P$ which is a complete local Noetherian ring, a flat $W(P)$-module, and such that $B_P\otimes_{W(P)}P$ is $P$-isomorphic to $B_0$.
Using again the fact that $W(k)$ is a $W(P)$-algebra formally smooth, one sees by \sref{0.19.3.11} that the homomorphism $W(k)\to k\to B_0$ factors as $W(k)\to B_P\to B_0$; moreover, as $k=W(k)/pW(k)$, one has $B_P\otimes_{W(k)}k = B_P/pB_P = B_P\otimes_{W(P)}P = B_0$.
Let us show that $B_P$ is a flat $W(k)$-module; as $W(k)$ is a discrete valuation ring with $p$ as uniformiser, it suffices to verify that $B_P$ is a $W(k)$-module \emph{without torsion} ($\mathbf{0_I}$, 6.3.4), or equivalently that $p$ is a $B_P$-regular element, which follows from the fact that $B_P$ is a flat $W(P)$-module ($\mathbf{0_I}$, 6.3.4).

Set now
\[
  B = B_P\widehat{\otimes}_{W(k)} A
\]
and note that the residue field of $A$ being equal to that of $W(k)$, is \emph{a fortiori} a $W(k)$-module of finite type.
It follows first from \sref{0.19.7.1.2} that $B$ is a complete semi-local Noetherian ring, $\mathfrak{m}B$ is contained in the radical of $B$; moreover $B/\mathfrak{m}B$ is $k$-isomorphic to $B_P\otimes_{W(k)}k = B_0$, hence $B$ is in fact a \emph{local} ring.
As $B_P$ is a flat $W(k)$-module, \sref{0.19.7.1.2} shows finally that $B$ is a flat $A$-module. \textsc{q.e.d.}
\end{proof}

\begin{lemma}[19.7.1.4]
\label{0.19.7.1.4}
Let $A$ be a ring, $\mathfrak{J}$ an ideal of $A$, $M$, $N$ two $A$-modules separated for the $\mathfrak{J}$-preadic topology.
Suppose further that $N$ is complete for the $\mathfrak{J}$-preadic topology and that $M$ is a flat $A$-module.
Let $u:N\to M$ be an $A$-homomorphism; if $u\otimes 1:N\otimes_A(A/\mathfrak{J})\to M\otimes_A(A/\mathfrak{J})$ is bijective, then $u$ is bijective.
\end{lemma}

\begin{proof}
The graded modules being taken relative to the $\mathfrak{J}$-preadic filtrations, it follows from the hypotheses on $M$ and $N$ relative to the $\mathfrak{J}$-preadic topologies that it suffices to prove that $\gr(u):\gr_\bullet(N)\to\gr_\bullet(M)$ is \emph{bijective} (Bourbaki, \emph{Alg.\ comm.}, chap.~III, \S\,2, n$^\circ$\,8, cor.~3 of th.~1).
Now, there is a commutative diagram
\[
\xymatrix{
  \gr_0(N)\otimes_{A_0}\gr_\bullet(A) \ar[r]^-{\gr_0(u)\otimes 1} \ar[d]_{\varphi_N} & \gr_0(M)\otimes_{A_0}\gr_\bullet(A) \ar[d]^{\varphi_M} \\
  \gr_\bullet(N) \ar[r]_-{\gr(u)} & \gr_\bullet(M)
}
\]
\oldpage[0\textsubscript{IV-1}]{108}
where $A_0=A/\mathfrak{J}=\gr_0(A)$, and $\varphi_M$ and $\varphi_N$ are the canonical maps ($\mathbf{0_{III}}$, 10.1.1.2).
By hypothesis, $\gr_0(u)$ is bijective as well as $\varphi_M$ ($\mathbf{0_{III}}$, 10.2.1), and $\varphi_N$ is surjective; one deduces first that $\gr_0(u)\otimes 1$ is bijective, then that $\varphi_N$ is injective, hence bijective, and finally that $\gr(u)$ is bijective.
\end{proof}

\begin{lemma}[19.7.1.5]
\label{0.19.7.1.5}
Let $A$ be a Noetherian ring, $\mathfrak{J}$ an ideal of $A$, $B$, $B'$ two $A$-algebras that are local Noetherian rings, the homomorphisms $A\to B$, $A\to B'$ being continuous for the $\mathfrak{J}$-preadic topology on $A$.
Suppose that:
\begin{enumerate}
  \item[$1^\circ$] $B$ and $B'$ are complete for the $\mathfrak{J}$-preadic topologies;
  \item[$2^\circ$] $B$ is a formally smooth $A$-algebra;
  \item[$3^\circ$] $B'$ is a flat $A$-module.
\end{enumerate}
Set $A_0=A/\mathfrak{J}$, and let $u_0:B\otimes_A A_0\to B'\otimes_A A_0$ be an $A_0$-isomorphism; then there exists an $A$-isomorphism $u:B\to B'$ such that $u_0=u\otimes 1$ (which implies that $B'$ is also a formally smooth $A$-algebra and $B$ is a flat $A$-module).
\end{lemma}

\begin{proof}
Set $B_0=B\otimes_A A_0$, $B_0'=B'\otimes_A A_0$.
Note that if $\mathfrak{m}$ and $\mathfrak{m}'$ are the maximal ideals of $B$ and $B'$, the $\mathfrak{J}$-preadic topologies on $B$ and $B'$ are separated since $\mathfrak{J}B\subset\mathfrak{m}$ and $\mathfrak{J}B'\subset\mathfrak{m}'$; moreover, as $\mathfrak{J}B'$ is closed in $B'$ for the $\mathfrak{J}$-preadic topology, the composite homomorphism $B\to B_0\xrightarrow{u_0}B_0'$, which is continuous for the $\mathfrak{J}$-preadic topologies, factors as $B\xrightarrow{u}B'\to B'_0$, where $u$ is a continuous $A$-homomorphism \sref{0.19.3.10}.
One has $u_0=u\otimes 1$, and the hypothesis that $u_0$ is bijective implies that $u$ is bijective, by virtue of \sref{0.19.7.1.4}.
\end{proof}

\medskip
\noindent\textbf{(19.7.1.6)}\enspace
\emph{End of the proof of} \sref{0.19.7.1}. ---
To finish proving \sref{0.19.7.1}, it remains to show that $a)$ implies $b)$; it is already known that $a)$ implies that $B_0$ is a $k$-algebra formally smooth \sref{0.19.3.5}, (iii), hence it suffices to prove that $B$ is a flat $A$-module.
It amounts to the same to establish that $\widehat{B}$ is a flat $\widehat{A}$-module (Bourbaki, \emph{Alg.\ comm.}, chap.~III, \S\,5, n$^\circ$\,4, prop.~4), and one knows that $\widehat{B}$ is an $\widehat{A}$-algebra formally smooth \sref{0.19.3.6}; one can therefore restrict to the case where $A$ and $B$ are \emph{complete}.
As $B_0$ is a $k$-algebra formally smooth, it is a regular ring \sref{0.19.6.5} and complete ($\mathbf{0_I}$, 6.3.5); applying \sref{0.19.7.1.3}, one sees that there exists an $A$-algebra $B'$ which is a complete local Noetherian ring and a flat $A$-module, a local homomorphism $A\to B'$, and a $k$-isomorphism $B_0\simeq B_0'=B'\otimes_A k$.
It then suffices to apply \sref{0.19.7.1.5} (taking $\mathfrak{J}$ to be the maximal ideal of $A$) to obtain that $B$ is $A$-isomorphic to $B'$, hence is a flat $A$-module. \textsc{q.e.d.}

\begin{theorem}[19.7.2]
\label{0.19.7.2}
Let $A$ be a local Noetherian ring, $\mathfrak{J}$ an ideal contained in the maximal ideal of $A$, $A_0=A/\mathfrak{J}$, $B_0$ a complete local Noetherian ring, $A_0\to B_0$ a local homomorphism making $B_0$ into a formally smooth $A_0$-algebra.
Then there exists a complete local Noetherian ring $B$, a local homomorphism $A\to B$ making $B$ into a flat $A$-module, and an $A_0$-isomorphism $u:B\otimes_A A_0\xrightarrow{\sim} B_0$.
If $(B',u')$ is a pair satisfying the same conditions as $(B,u)$, there exists an $A$-isomorphism $v:B\xrightarrow{\sim} B'$ making the diagram
\[
\xymatrix{
  B\otimes_A A_0 \ar[rr]^{v\otimes 1}_{\sim} \ar[dr]^{\sim}_u && B'\otimes_A A_0 \ar[dl]^{u'}_{\sim} \\
  & B_0
}
\]
commutative.
\end{theorem}

\oldpage[0\textsubscript{IV-1}]{109}
\begin{proof}
Let $\mathfrak{m}$ be the maximal ideal of $A$, so that $\mathfrak{m}_0=\mathfrak{m}/\mathfrak{J}$ is the maximal ideal of $A_0$, $A$ and $A_0$ having the same residue field $k$.
Set $B_{00}=B_0\otimes_{A_0} k$; as $B_{00}$ is a formally smooth $k$-algebra \sref{0.19.3.5}, (iii), it is a regular local ring \sref{0.19.6.5}; applying \sref{0.19.7.1.3}, one sees that there exists a topological $A$-algebra $B$ which is a complete local Noetherian ring, a flat $A$-module, and for which one has a $k$-isomorphism $u_{00}:B\otimes_A k\xrightarrow{\sim} B_{00}$.
Note that by virtue of \sref{0.19.7.1}, $B$ is a formally smooth $A$-algebra, hence $B\otimes_A A_0=B/\mathfrak{J}B$ is a formally smooth $A_0$-algebra \sref{0.19.3.5}, (iii) and a complete local Noetherian ring; moreover, $B_0$ is a flat $A_0$-module by virtue of the hypothesis and of \sref{0.19.7.1}; as one has a $k$-isomorphism $u_{00}:B\otimes_A k=(B\otimes_A A_0)\otimes_{A_0} k\xrightarrow{\sim} B_0\otimes_{A_0} k = B_{00}$, one deduces from \sref{0.19.7.1.5}, applied by replacing $A$ by $A_0$ and $\mathfrak{J}$ by $\mathfrak{m}_0$, that there exists an $A_0$-isomorphism $u:B\otimes_A A_0\xrightarrow{\sim} B_0$ such that $u_{00}=u\otimes 1$.
As for the uniqueness assertion, note that the ideals $\mathfrak{J}^h B$ (resp.\ $\mathfrak{J}^h B'$) are closed in $B$ (resp.\ $B'$) ($\mathbf{0_I}$, 7.3.5), hence $B$ and $B'$ are separated and complete for the $\mathfrak{J}$-preadic topologies (Bourbaki, \emph{Top.\ g\'{e}n.}, chap.~III, 3$^\text{e}$ \'{e}d., \S\,3, n$^\circ$\,5, cor.~2 of prop.~9); one has by hypothesis an $A_0$-isomorphism $v_0:B\otimes_A A_0\xrightarrow{\sim} B'\otimes_A A_0$ such that $u'v_0=u$; as $B$ is a formally smooth $A$-algebra and $B'$ a flat $A$-module, one can apply \sref{0.19.7.1.5}, whence the existence of the $A$-isomorphism $v$ answering the question.
\end{proof}

\begin{remarks}[19.7.3]
\label{0.19.7.3}
\begin{enumerate}
  \item[\rm(i)] We note that the uniqueness assertion in \sref{0.19.7.2} is still valid if one only supposes that $B$ and $B'$ are complete \emph{for the $\mathfrak{J}$-preadic topologies}.
We do not know whether one can similarly improve the existence assertion, that is, whether one can dispense with assuming the local ring $B_0$ to be complete (for its $\mathfrak{n}_0$-preadic topology, where $\mathfrak{n}_0$ denotes its maximal ideal) by only requiring that $B$ be complete for the $\mathfrak{J}$-preadic topology.
When $A_0$ is complete for the $\mathfrak{m}$-preadic topology, one can see that this problem reduces to the following: if $B_0$ is a regular local Noetherian ring (not necessarily complete) containing the prime field $\mathbf{F}_p=\mathbf{Z}/p\mathbf{Z}$, does there exist for every $n>1$ a flat $(\mathbf{Z}/p^n\mathbf{Z})$-algebra $B$ such that $B/pB$ is isomorphic to $B_0$?
  \item[\rm(ii)] We note that in general, the isomorphism $v$ whose existence is asserted in \sref{0.19.7.2} is \emph{not unique} (cf.\ \sref{0.19.8.7}).
\end{enumerate}
\end{remarks}

\subsection{Cohen algebras and Cohen rings; application to the structure of complete local rings}
\label{subsection:0.19.8}
\label{0.19.8}

The results of this section are immediate applications of the theorems of \sref{subsection:0.19.7}, but deserve to be stated explicitly because of their practical importance.

\begin{definition}[19.8.1]
\label{0.19.8.1}
Let $A$, $B$ be two local Noetherian rings, $\mathfrak{m}$ the maximal ideal of $A$, $k=A/\mathfrak{m}$ its residue field, $\varphi:A\to B$ a local homomorphism, making $B$ into an $A$-algebra.
We say that $B$ is a \emph{Cohen $A$-algebra} if it satisfies the following conditions:
\begin{enumerate}
  \item[\rm(i)] $B$ is a complete ring.
  \item[\rm(ii)] $B$ is a flat $A$-module.
  \item[\rm(iii)] $B\otimes_A k$ is a field (in other words, $\mathfrak{m}B$ is the maximal ideal of $B$) which is a separable extension of $k$.
\end{enumerate}
\end{definition}

\oldpage[0\textsubscript{IV-1}]{110}
\begin{theorem}[19.8.2]
\label{0.19.8.2}
Let $A$ be a local Noetherian ring, $k$ its residue field.
\begin{enumerate}
  \item[\rm(i)] If $B$ is a Cohen $A$-algebra, $B$ is a formally smooth $A$-algebra.
For every complete local Noetherian ring $C$, every local homomorphism $A\to C$ and every ideal $\mathfrak{J}\neq C$ in $C$, every $A$-homomorphism $B\to C/\mathfrak{J}$ therefore factors as $B\xrightarrow{v}C\to C/\mathfrak{J}$, where $v$ is an $A$-homomorphism (necessarily local).
  \item[\rm(ii)] For every field $K$, separable extension of $k$, there exists a Cohen $A$-algebra $B$ such that $B\otimes_A k$ is $k$-isomorphic to $K$, and such an $A$-algebra is unique up to isomorphism.
\end{enumerate}
\end{theorem}

\begin{proof}
As $B\otimes_A k$ is a separable field extension of $k$, it is a formally smooth $k$-algebra \sref{0.19.6.1}, so assertion (i) follows from \sref{0.19.7.1}.
To prove (ii), one can restrict to the case where $A$ is \emph{complete}, for it amounts to the same to say that $B$ is a flat $A$-module or a flat $\widehat{A}$-module ($\mathbf{0_{III}}$, 10.2.3), one has $\mathfrak{m}B=\mathfrak{m}\widehat{A}B$ and $k$ is the residue field of $\widehat{A}$.
It then suffices to apply \sref{0.19.7.2} taking $\mathfrak{J}=\mathfrak{m}$ and $B_0=K$ (and using \sref{0.19.6.1}).
\end{proof}

\begin{definition}[19.8.3]
\label{0.19.8.3}
A \emph{prime local ring} is a local ring of the form $\mathbf{Z}_{p\mathbf{Z}}$, where $p\mathbf{Z}$ is a prime ideal of $\mathbf{Z}$.
A \emph{complete prime local ring} is the completion of a prime local ring.
\end{definition}

The prime local rings are therefore of two kinds:
\begin{enumerate}
  \item[$1^\circ$] Those corresponding to maximal ideals $p\mathbf{Z}$ where $p\neq 0$ is a prime number; $\mathbf{Z}_{p\mathbf{Z}}$ is a discrete valuation ring, whose completion is the ring of $p$-adic integers usually denoted $\mathbf{Z}_p$.\footnote{This notation, now universally used, is in conflict with the notation $A_f$ adopted in ($\mathbf{0_I}$, 1.2.3): with $A=\mathbf{Z}$ and $f=p$, $A_f$ denotes the ring of rational numbers of the form $k/p^n$ ($k\in\mathbf{Z}$, $n$ an integer $\geq 0$); we will always avoid using the notation $\mathbf{Z}_p$ to denote this latter ring.}
  \item[$2^\circ$] For the prime ideal $p\mathbf{Z}=(0)$, $\mathbf{Z}_{p\mathbf{Z}}$ is the field of rational numbers $\mathbf{Q}$, which is identical to its completion (the topology being naturally the topology of a local Noetherian ring, hence discrete here).
\end{enumerate}

The terminology of \sref{0.19.8.3}, analogous to that of ``prime fields'', is justified in the same way: for every local ring $A$, consider the canonical homomorphism $\varphi:\mathbf{Z}\to A$, and let $p\mathbf{Z}$ be the inverse image under this homomorphism of the maximal ideal $\mathfrak{m}$ of $A$; $p\mathbf{Z}$ is a prime ideal of $\mathbf{Z}$ and the preceding homomorphism therefore factors as $\mathbf{Z}\to\mathbf{Z}_{p\mathbf{Z}}\xrightarrow{\psi}A$; moreover, as $\varphi$ is the \emph{unique} homomorphism of $\mathbf{Z}$ into $A$, $p$ and $\psi$ are determined in a unique way.
In other words, for every local ring $A$, there is a unique homomorphism $\psi:P\to A$, where $P$ is a prime local ring; if moreover $A$ is separated and complete, one can extend this homomorphism by completion, and there is thus a unique homomorphism $\hat{\psi}:\widehat{P}\to A$, where $\widehat{P}$ is a complete prime local ring.
Moreover, $\psi$ gives by passage to quotients a homomorphism of the residue field $\mathbf{F}_p$ if $p>0$ (resp.\ $\mathbf{Q}$ if $p=0$) into the residue field $k$ of $A$, and $p$ is therefore the \emph{characteristic} of $k$.

If one takes in particular $A$ to be a prime local ring (resp.\ complete prime local ring), one sees that there exists in such a ring only \emph{a single endomorphism}, namely the identity.

\begin{definition}[19.8.4]
\label{0.19.8.4}
\oldpage[0\textsubscript{IV-1}]{111}
Let $A$ be a local ring, $P\to A$ the unique homomorphism from a prime local ring $P$ into $A$, $p$ the characteristic of the residue fields of $P$ and $A$.
We say that $A$ is a \emph{Cohen ring} if it is a Cohen $P$-algebra, that is \sref{0.19.8.1}, if:
\begin{enumerate}
  \item[$1^\circ$] $A$ is Noetherian and complete.
  \item[$2^\circ$] $A$ is a flat $P$-module (which also amounts to saying that $A$ is a flat $\widehat{P}$-module (Bourbaki, \emph{Alg.\ comm.}, chap.~III, \S\,5, n$^\circ$\,4, prop.~4)).
  \item[$3^\circ$] $A/pA$ is a field (necessarily separable over the residue field of $P$, this field being prime).
\end{enumerate}
\end{definition}

If $p=0$, these conditions are equivalent to saying that $A$ is a field of characteristic $0$.
If $p>0$, one necessarily has $pA\neq 0$; condition $3^\circ$ means that $pA$ is the maximal ideal $\mathfrak{m}$ of $A$; condition $2^\circ$ means that $p$ is $A$-regular, since $P$ is a discrete valuation ring ($\mathbf{0_I}$, 6.3.4).
Hence $A$ is a regular ring \sref{0.17.1.1}, $d)$) of dimension $1$, and therefore a discrete valuation ring, complete by virtue of $1^\circ$; in summary:

\begin{proposition}[19.8.5]
\label{0.19.8.5}
The Cohen rings are the fields of characteristic $0$ and the complete discrete valuation rings whose residue field has characteristic $p>0$, and whose maximal ideal is generated by $p\cdot 1$ ($1$ being the unity of the ring).
\end{proposition}

We note that in the second case, $p\cdot 1\neq 0$ since $p$ is $A$-regular, hence one can identify $p\cdot 1$ with the integer $p$; the canonical homomorphism $\mathbf{Z}_p\to A$ is injective, and one identifies $p\cdot 1$ with the element $p$ of $\mathbf{Z}_p$; one says in this case that $A$ is a \emph{$p$-ring} (or \emph{Cohen $p$-ring}).

\begin{theorem}[19.8.6]
\label{0.19.8.6}
\textup{(Cohen).} ---
\begin{enumerate}
  \item[\rm(i)] Let $W$ be a Cohen ring, $C$ a complete local Noetherian ring, $\mathfrak{J}$ an ideal of $C$ distinct from $C$.
Then every local homomorphism $u:W\to C/\mathfrak{J}$ factors as $W\xrightarrow{v}C\to C/\mathfrak{J}$ where $v$ is a local homomorphism.
  \item[\rm(ii)] Let $K$ be a field.
There exists a Cohen ring $W$ whose residue field is isomorphic to $K$.
If $W'$ is a second Cohen ring, $K'$ its residue field, every isomorphism $u:K\xrightarrow{\sim}K'$ arises by passage to quotients from an isomorphism $v:W\xrightarrow{\sim}W'$.
\end{enumerate}
\end{theorem}

\begin{proof}
This is nothing other than \sref{0.19.8.2} applied to the case where $A$ is a prime local ring.
\end{proof}

\begin{remarks}[19.8.7]
\label{0.19.8.7}
\begin{enumerate}
  \item[\rm(i)] When $K$ is of characteristic $0$, part (ii) of \sref{0.19.8.6} becomes trivial.
  \item[\rm(ii)] The homomorphism $v$ of \sref{0.19.8.6}, (i) is \emph{not} necessarily determined uniquely by $u$, as is shown already by the case where $W$ is a field of characteristic $0$, $\mathfrak{J}^2=0$ and $u$ is an isomorphism (cf.\ (21.5.5)).
Similarly, in \sref{0.19.8.6}, (ii), the isomorphism $v$ is not necessarily determined uniquely by $u$ (cf.\ (21.5.5)).
However, when $K$ is \emph{perfect} and of characteristic $p>0$, one will see (21.5.5) that in \sref{0.19.8.6}, (ii) the isomorphism $v$ is unique.
One will also see later that in this case $W$ identifies with the ring $W_\infty(K)$ of \emph{Witt vectors of infinite length} over $K$.
  \item[\rm(iii)] In \sref{0.19.8.6}, (i), one can weaken the hypotheses on $C$ by using \sref{0.19.3.10} and \sref{0.19.3.12}.
\end{enumerate}
\end{remarks}

\begin{theorem}[19.8.8]
\label{0.19.8.8}
\oldpage[0\textsubscript{IV-1}]{112}
\textup{(Cohen).} ---
Let $A$ be a complete local Noetherian ring, $k$ its residue field.
\begin{enumerate}
  \item[\rm(i)] There exists a Cohen ring $W$ such that $A$ is isomorphic to a quotient ring of a formal power series ring $W[[T_1,\ldots,T_n]]$ (and in particular $A$ is isomorphic to a quotient of a complete regular local ring \sref{0.17.3.8}).
If $A$ contains a field, it is isomorphic to a quotient ring of $k[[T_1,\ldots,T_n]]$.
  \item[\rm(ii)] Suppose moreover that $A$ is an integral domain.
Then there exists a subring $B$ of $A$ such that: $1^\circ$ $B$ is isomorphic to a formal power series ring over a ring $C$ which is a field or a Cohen ring (which implies that $B$ is a complete regular local ring \sref{0.17.3.8}); $2^\circ$ $B$ has the same residue field as $A$ and the injection $B\to A$ is a local homomorphism; $3^\circ$ $A$ is a finite $B$-algebra.
\end{enumerate}
\end{theorem}

\begin{proof}
Let $\mathfrak{m}$ be the maximal ideal of $A$.
There exists a Cohen ring $W$ whose residue field is isomorphic to $k$ \sref{0.19.8.6}, (ii); one therefore has a homomorphism $W\to A/\mathfrak{m}$, which then factors as $W\xrightarrow{u}A\to A/\mathfrak{m}$, where $u$ is a local homomorphism \sref{0.19.8.6}, (i).
For every finite family $(x_i)_{1\leq i\leq n}$ of elements of $\mathfrak{m}$, there then exists a local homomorphism $v:W[[T_1,\ldots,T_n]]\to A$ extending $u$ and such that $v(T_i)=x_i$ for all $i$ (Bourbaki, \emph{Alg.\ comm.}, chap.~III, \S\,4, n$^\circ$\,5, prop.~6).
When $A$ contains a field, it contains a prime field $P$, of which $k$ is an extension (necessarily separable), and therefore $A$ contains a field isomorphic to $k$ \sref{0.19.6.2}; one can then replace $W$ by $k$ in the preceding definition of $v$.

\medskip
\noindent(i)\enspace
Let us first take the $x_i$ to be a system of generators of $\mathfrak{m}$.
As $W$ has the same residue field as $A$, and the classes of the $x_i$ generate the graded ring $\gr_\bullet(A)$ as a $k$-algebra, $\gr(v):\gr_\bullet(W[[T_1,\ldots,T_n]])\to\gr_\bullet(A)$ is surjective; one deduces that $v$ is itself surjective (Bourbaki, \emph{Alg.\ comm.}, chap.~III, \S\,2, n$^\circ$\,8, cor.~2 of th.~1).
We recall that the case where $A$ contains a field has already been seen and is included here only for convenience \sref{0.19.6.3}.

\medskip
\noindent(ii)\enspace
If $A$ contains a field, it contains a field $k'$ isomorphic to $k$ as we have seen; one then considers a system of parameters $(y_j)_{1\leq j\leq m}$ of $A$ \sref{0.16.3.6}, one takes $B=k[[T_1,\ldots,T_m]]$ and one considers the local homomorphism $w:B\to A$ which coincides in $k$ with an isomorphism $k\to k'$ and is such that $w(T_j)=y_j$ for $1\leq j\leq m$.
If $A$ does not contain a field, the unique homomorphism $\mathbf{Z}_{p\mathbf{Z}}\to A$ \sref{0.19.8.3} is necessarily injective (for otherwise, as $A$ is an integral domain, its kernel would be the maximal ideal of $\mathbf{Z}_{p\mathbf{Z}}$ and its image would be isomorphic to a field; moreover one has $p>0$ by hypothesis, $\mathbf{Z}_{p\mathbf{Z}}$ being a field if $p=0$); the element $p\cdot 1$ of $A$ (identified with $p$) is a non-zero-divisor in $A$, and is contained in $\mathfrak{m}$, hence \sref{0.16.3.4} and \sref{0.16.3.7} there exist $m-1$ other elements $x_i$ ($2\leq i\leq m$) of $\mathfrak{m}$ forming, with $p$, a system of parameters of $A$.
The Cohen ring $W$ considered at the beginning of the proof is then a discrete valuation ring whose residue field is isomorphic to $k$, in which $p$ generates the maximal ideal \sref{0.19.8.5}, and the homomorphism $u:W\to A$ defined at the beginning maps $p$ to itself.
One then takes $B=W[[T_1,\ldots,T_{m-1}]]$ and considers the local homomorphism $w:B\to A$ which coincides with $u$ in $W$ and is such that $w(T_i)=x_i$ for $1\leq i\leq m-1$.
\oldpage[0\textsubscript{IV-1}]{113}
In both cases, if $\mathfrak{n}$ is the maximal ideal of $B$, it is clear that $\mathfrak{n}A$ is an ideal of definition of $A$; as moreover $B/\mathfrak{n}$ and $A/\mathfrak{m}$ are isomorphic, $A$ is a quasi-finite $B$-module ($\mathbf{0_I}$, 7.4.4), hence a $B$-module of finite type since $B$ is complete and $A$ is separated for the $\mathfrak{n}$-preadic topologies ($\mathbf{0_I}$, 7.4.1).
On the other hand, in both cases one has $\dim(B)=\dim(A)=m$; in the first case, this follows from \sref{0.17.1.4}, (iii); in the second, one sees directly that $p$ and the $T_i$ ($1\leq i\leq m-1$) form a $B$-regular sequence generating $\mathfrak{n}$, or one can also use the fact that these elements generate $\mathfrak{n}$ and that $\dim(B)\geq\dim(A)$ by \sref{0.16.3.10}.
As $A$ and $B$ are integral domains, one finally deduces from \sref{0.16.3.10} that $w$ is injective, which completes the proof.
\end{proof}

\begin{corollary}[19.8.9]
\label{0.19.8.9}
Let $A$ be a complete local Noetherian integral domain containing a field $k_0$; let $k$ be the residue field of $A$, and suppose that $k$ is finite over $k_0$.
Then, in the conclusion of \sref{0.19.8.8}, (ii), one can replace $1^\circ$ and $2^\circ$ by the condition that $B$ is of the form $k_0[[T_1,\ldots,T_n]]$, the canonical injection $B\to A$ being a local $k_0$-homomorphism (for the usual $k_0$-algebra structure of $B$).
\end{corollary}

\begin{proof}
Indeed, taking up the proof of \sref{0.19.8.8}, (ii), one defines this time $w:k_0[[T_1,\ldots,T_n]]\to A$ as coinciding in $k_0$ with the identity and mapping $T_j$ to $y_j$ for $1\leq j\leq n$.
The hypothesis that $k$ is of finite degree over $k_0$ again implies that $A$ is a quasi-finite $B$-module, hence of finite type by ($\mathbf{0_I}$, 7.4.1), and one concludes as in \sref{0.19.8.8}.
\end{proof}

\begin{corollary}[19.8.10]
\label{0.19.8.10}
Let $A$ be a local Artinian ring whose maximal ideal $\mathfrak{m}$ is of square zero; then there exists a regular local Noetherian ring $B$, with maximal ideal $\mathfrak{n}$, such that $A$ is isomorphic to $B/\mathfrak{n}^2$.
\end{corollary}

\begin{proof}
Let $K$ be the residue field $A/\mathfrak{m}$ of $A$, and $n$ the rank of $\mathfrak{m}/\mathfrak{m}^2=\mathfrak{m}$ over $K$.
If $A$ contains a field, it follows from \sref{0.19.6.3} that $A$ is isomorphic to $B/\mathfrak{b}$, where $B=K[[T_1,\ldots,T_n]]$ and $\mathfrak{b}$ is contained in the square of the maximal ideal $\mathfrak{n}$ of $B$; but as $\operatorname{long}(B/\mathfrak{n}^2)=n+1=\operatorname{long}(A)$, one necessarily has $\mathfrak{b}=\mathfrak{n}^2$.

Suppose next that $A$ does not contain a field; this implies that $K$ is of characteristic $p>0$ and that $p\cdot 1\neq 0$ in $A$ \sref{0.19.6.3}; hence $p\cdot 1$ is an element of $\mathfrak{m}$, and there are therefore $n-1$ other elements $x_i$ ($2\leq i\leq n$) of $\mathfrak{m}$ forming with $p\cdot 1$ a basis of $\mathfrak{m}$ over $K$.
Let $W$ be a Cohen ring whose residue field is isomorphic to $K$; $W$ is a discrete valuation ring in which $p$ generates the maximal ideal; one has seen in the proof of \sref{0.19.8.8} that there is a homomorphism $u:W\to A$ mapping $p$ to itself and which by passage to quotients gives the identity on $K$.
One takes $B=W[[T_2,\ldots,T_n]]$ and considers the local homomorphism $w:B\to A$ which coincides with $u$ in $W$ and is such that $w(T_i)=x_i$ for $2\leq i\leq n$.
It is clear that $w$ is surjective and that its kernel $\mathfrak{b}$ is contained in the square of the maximal ideal $\mathfrak{n}=pB+BT_2+\ldots+BT_n$ of $B$; as $\operatorname{long}(B/\mathfrak{n}^2)=n+1=\operatorname{long}(A)$, one again has $\mathfrak{b}=\mathfrak{n}^2$.
\end{proof}

\begin{proposition}[19.8.11]
\label{0.19.8.11}
Let $A$ be a local Artinian ring, $\mathfrak{m}$ its maximal ideal, $k$ its residue field.
For $A$ to be isomorphic to a quotient ring of a Cohen ring, it is necessary and sufficient that $\mathfrak{m}$ be generated by $p\cdot 1$, where $p$ is the characteristic of $k$.
\end{proposition}

\begin{proof}
\oldpage[0\textsubscript{IV-1}]{114}
The condition is evidently necessary \sref{0.19.8.5}.
To see that it is sufficient, one observes, as at the beginning of the proof of \sref{0.19.8.8}, that there exists a Cohen ring $W$ whose residue field is isomorphic to $k$ and a local homomorphism $u:W\to A$.
Moreover, if one considers the composite homomorphism $\mathbf{Z}_{p\mathbf{Z}}\to W\xrightarrow{u}A$ (which is necessarily the unique homomorphism of $\mathbf{Z}_{p\mathbf{Z}}$ into $A$), one sees that the image under $u$ of the element $p\cdot 1$ of $W$ is the element $p\cdot 1$ of $A$; as the element $p\cdot 1$ of $W$ generates the maximal ideal of this ring, one deduces immediately from the hypothesis that $\gr(u):\gr_\bullet(W)\to\gr_\bullet(A)$ is surjective, and therefore so is $u$ (Bourbaki, \emph{Alg.\ comm.}, chap.~III, \S\,2, n$^\circ$\,8, cor.~2 of th.~1).
\end{proof}

\subsection{Relatively formally smooth algebras}
\label{subsection:0.19.9}

\begin{definition}[19.9.1]
\label{0.19.9.1}
Let $\Lambda$ be a topological ring, $A$ a topological $\Lambda$-algebra, $B$ a topological $A$-algebra.
We say that $B$ is a \emph{formally smooth $A$-algebra relative to $\Lambda$} if, for every discrete topological $A$-algebra $C$, and every nilpotent ideal $\mathfrak{J}$ of $C$, every continuous $A$-homomorphism $u_0:B\to C/\mathfrak{J}$ which factors as $B\xrightarrow{u}C\to C/\mathfrak{J}$, where $u$ is a continuous $\Lambda$-homomorphism, also factors as $B\xrightarrow{v}C\to C/\mathfrak{J}$, where $v$ is a continuous $A$-homomorphism.
\end{definition}

It follows from this definition that if $B$ is a \emph{formally smooth} $A$-algebra, then $B$ is also formally smooth relative to $\Lambda$, for every structure of topological $\Lambda$-algebra defined on $A$ (in other words, every continuous ring homomorphism $\Lambda\to A$).

\begin{proposition}[19.9.2]
\label{0.19.9.2}
Let $\Lambda$ be a topological ring, $A$ a topological $\Lambda$-algebra.
\begin{enumerate}
  \item[\rm(i)] $A$ is a $\Lambda$-algebra formally smooth relative to $\Lambda$.
  \item[\rm(ii)] If $B$ is an $A$-algebra formally smooth relative to $\Lambda$ and $C$ is a $B$-algebra formally smooth relative to $\Lambda$, then $C$ is an $A$-algebra formally smooth relative to $\Lambda$.
  \item[\rm(iii)] Let $B$ be an $A$-algebra formally smooth relative to $\Lambda$, $A'$ a topological $A$-algebra, then the topological $A'$-algebra $B\otimes_A A'$ is formally smooth relative to $\Lambda$.
  \item[\rm(iv)] Let $B$ be a topological $A$-algebra, $S$ (resp.\ $T$) a multiplicative subset of $A$ (resp.\ $B$) such that the canonical image of $S$ in $B$ is contained in $T$.
If $B$ is an $A$-algebra formally smooth relative to $\Lambda$, then $T^{-1}B$ is an $S^{-1}A$-algebra formally smooth relative to $\Lambda$.
  \item[\rm(v)] Let $B_i$ ($1\leq i\leq n$) be topological $A$-algebras.
For $\prod_{i=1}^n B_i$ to be an $A$-algebra formally smooth relative to $\Lambda$, it is necessary and sufficient that each of the $B_i$ be so.
\end{enumerate}
\end{proposition}

\begin{proof}
Assertion (i) is trivial, and the proof of the others closely follows the proofs of \sref{0.19.3.5}; it is therefore left to the reader.
\end{proof}

\begin{corollary}[19.9.3]
\label{0.19.9.3}
Let $A$ be a topological ring, $A'$ and $B$ two topological $A$-algebras.
Then the topological $A$-algebra $B'=B\otimes_A A'$ is a $B$-algebra formally smooth relative to $A$.
\end{corollary}

\begin{proof}
This follows from \sref{0.19.9.2}, (i) and (iii).
\end{proof}

\begin{proposition}[19.9.4]
\label{0.19.9.4}
Let $\Lambda$ be a topological ring, $A$ a topological $\Lambda$-algebra, $B$ a topological $A$-algebra.
The following conditions are equivalent:
\begin{enumerate}
  \item[\rm a)] $B$ is an $A$-algebra formally smooth relative to $\Lambda$.
  \item[\rm b)] $\widehat{B}$ is an $A$-algebra formally smooth relative to $\Lambda$.
  \item[\rm c)] $\widehat{B}$ is an $\widehat{A}$-algebra formally smooth relative to $\Lambda$.
  \item[\rm d)] $\widehat{B}$ is an $\widehat{A}$-algebra formally smooth relative to $\widehat{\Lambda}$.
\end{enumerate}
\end{proposition}

\begin{proof}
The proof is again left to the reader, modelled on that of \sref{0.19.3.6}.
\end{proof}

\oldpage[0\textsubscript{IV-1}]{115}
\begin{env}[19.9.5]
\label{0.19.9.5}
Similarly, the statement \sref{0.19.3.8} is still valid (with the same proof) when one replaces ``formally smooth'' by ``formally smooth relative to $\Lambda$''.
If in the statement of \sref{0.19.3.10} one replaces ``formally smooth'' by ``formally smooth relative to $\Lambda$'', the conclusion is replaced by the following (the proof remaining essentially unchanged): \emph{every continuous $A$-homomorphism $u_0:B\to C/\mathfrak{J}$ which factors as $B\xrightarrow{u}C\to C/\mathfrak{J}$, where $u$ is a continuous $\Lambda$-homomorphism, also factors as $B\xrightarrow{v}C\to C/\mathfrak{J}$, where $v$ is a continuous $A$-homomorphism}.
\end{env}

\begin{env}[19.9.6]
\label{0.19.9.6}
The criteria for formal smoothness \sref{0.19.4.1} and \sref{0.19.4.2} are also valid when one replaces ``formally smooth'' by ``formally smooth relative to $\Lambda$'', the proofs remaining practically unchanged.
\end{env}

\begin{proposition}[19.9.7]
\label{0.19.9.7}
Let $\Lambda$ be a topological ring, $A$ a topological $\Lambda$-algebra, $B$ a topological $A$-algebra.
Suppose that for every discrete $A$-algebra $C$ and every ideal $\mathfrak{J}$ of $C$ such that $\mathfrak{J}^2=0$, every continuous $A$-homomorphism $u_0:B\to C/\mathfrak{J}$ which factors as $B\xrightarrow{u}C\to C/\mathfrak{J}$, where $u$ is a continuous $\Lambda$-homomorphism, also factors as $B\xrightarrow{v}C\to C/\mathfrak{J}$, where $v$ is a continuous $A$-homomorphism.
Then $B$ is an $A$-algebra formally smooth relative to $\Lambda$.
\end{proposition}

\begin{proof}
The proof of \sref{0.19.4.3} transcribes immediately.
\end{proof}

\begin{proposition}[19.9.8]
\label{0.19.9.8}
Let $\Lambda$ be a topological ring, $A$ a topological $\Lambda$-algebra, $B$ a topological $A$-algebra.
For $B$ to be an $A$-algebra formally smooth relative to $\Lambda$, it is necessary and sufficient that for every discrete topological $B$-module $L$, annihilated by an open ideal of $B$, one has (cf.\ 18.4.2)
\[
  \operatorname{Exalcotop}_{A/\Lambda}(B,L) = 0.
\]
\end{proposition}

\begin{proof}
Using the notation of the proof of \sref{0.19.4.4}, it suffices to note here that one can suppose that the extension $E_\Lambda$ of $B_\Lambda$ is $\Lambda$-trivial; the rest of the proof is then unchanged.
\end{proof}

When $\Lambda$, $A$ and $B$ are discrete rings, the criterion \sref{0.19.9.8} reduces to
\begin{equation}
\tag{19.9.8.1}
  \label{0.19.9.8.1}
  \operatorname{Exalcom}_{A/\Lambda}(B,L) = 0 \quad\text{for every $B$-module $L$};
\end{equation}
in other words, \emph{every commutative $A$-extension of $B$ by a $B$-module, which is also $\Lambda$-trivial, is also $A$-trivial}.

\subsection{Formally unramified algebras and formally \'{e}tale algebras}
\label{subsection:0.19.10}
\label{0.19.10}

\begin{env}[19.10.1]
\label{0.19.10.1}
We will also need two notions closely related to that of a formally smooth algebra.
\end{env}

\begin{definition}[19.10.2]
\label{0.19.10.2}
Let $A$ be a topological ring, $B$ a topological $A$-algebra.
We say that $B$ is a \emph{formally unramified} $A$-algebra if, for every $A$-homomorphism $p:E\to C$ of discrete $A$-algebras, whose kernel is nilpotent, and every continuous $A$-homomorphism $u:B\to C$, there exists \emph{at most one} continuous $A$-homomorphism $v:B\to E$ such that $u$ factors as $B\xrightarrow{v}E\xrightarrow{p}C$.
We say that $B$ is a \emph{formally \'{e}tale} $A$-algebra if it is formally smooth and formally unramified, in other words, if, for every surjective $A$-homomorphism $p:E\to C$ of discrete $A$-algebras, whose kernel is nilpotent, and for every continuous $A$-homomorphism $u:B\to C$, there exists \emph{one and only one} continuous $A$-homomorphism $v:B\to E$ such that $u$ factors as $B\xrightarrow{v}E\xrightarrow{p}C$.
\end{definition}

\oldpage[0\textsubscript{IV-1}]{116}
\begin{examples}[19.10.3]
\label{0.19.10.3}
We restrict here to \emph{discrete} rings.
\begin{enumerate}
  \item[\rm(i)] If the structural homomorphism $\rho:A\to B$ is \emph{surjective}, $B$ is a \emph{formally unramified} $A$-algebra, for with the notation of \sref{0.19.10.2}, the composite homomorphism $A\xrightarrow{\rho}B\xrightarrow{v}E$ must be the structural homomorphism, hence $v$ is unique (if it exists).
On the other hand, $B$ is formally smooth if and only if $\ker(\rho)=\ker(\rho)^2$.
  \item[\rm(ii)] Let $A$ be a ring, $S$ a multiplicative subset of $A$; then $B=S^{-1}A$ is a \emph{formally \'{e}tale} $A$-algebra.
Indeed, let $j:A\to B$ be the canonical homomorphism, $\rho:A\to E$ a ring homomorphism, $\mathfrak{J}$ a nilpotent ideal of $E$, $p:E\to C=E/\mathfrak{J}$ the canonical homomorphism.
If $u:B\to C$ is an $A$-homomorphism, $u(s/1)$ is invertible in $C$ for $s\in S$ since $s/1$ is invertible in $B$; but $u(s/1)=p(\rho(s))$, and as $\mathfrak{J}$ is nilpotent, the fact that the image of $\rho(s)$ in $E/\mathfrak{J}$ is invertible implies that $\rho(s)$ is itself invertible in $E$ ($\mathbf{0_I}$, 7.1.12); one concludes that $\rho$ factors in a unique way as $A\xrightarrow{j}B\xrightarrow{v}E$, and one evidently has $u\circ j=(p\circ v)\circ j$, whence $u=p\circ v$, which proves our assertion.
  \item[\rm(iii)] A finite separable extension $k'$ of a field $k$ is a formally smooth $k$-algebra \sref{0.19.6.1}; we will see later that for a \emph{finite type} extension $K$ of a field $k$ to be a formally \'{e}tale $k$-algebra, it is necessary and sufficient that it be formally unramified, and this is also equivalent to $K$ being a finite and separable extension of $k$ (21.7.4).
  \item[\rm(iv)] A polynomial algebra $A[X_\alpha]_{\alpha\in I}$ over a ring $A$ is \emph{formally smooth} \sref{0.19.3.3} but \emph{not formally \'{e}tale} if $I$ is non-empty, as follows immediately from the definitions.
\end{enumerate}
\end{examples}

\begin{remark}[19.10.4]
\label{0.19.10.4}
The same reasoning as in \sref{0.19.4.3} shows that in order to verify that a topological $A$-algebra $B$ is formally unramified, one can, in the definition \sref{0.19.10.2}, restrict to the case where the kernel of $p:E\to C$ is of \emph{square zero}; one can moreover restrict to the case where $u(B)\subset p(E)$ (for otherwise there is no $v$ factoring $u$) and, replacing $C$ by $u(B)$ and $E$ by $p^{-1}(u(B))$, suppose $u$ and $p$ \emph{surjective}.
\end{remark}
